% =====================================================================================================
% Thesis section: construction of the STARFORGE + POLARIS dust-emission and polarization dataset.
% Include with % =====================================================================================================
% Thesis section: construction of the STARFORGE + POLARIS dust-emission and polarization dataset.
% Include with % =====================================================================================================
% Thesis section: construction of the STARFORGE + POLARIS dust-emission and polarization dataset.
% Include with % =====================================================================================================
% Thesis section: construction of the STARFORGE + POLARIS dust-emission and polarization dataset.
% Include with \input{dataset_construction}. Required packages in the preamble:
%   \usepackage{amsmath}  \usepackage{booktabs}  \usepackage{graphicx}  \usepackage{longtable}
% References are in dataset_construction.bib (\bibliography{dataset_construction}).
% Every number below is taken from the pipeline code, its logs, manifests and the audit report
% (paths in Sect. "Data products and code"); dates are UTC.
% =====================================================================================================

\section{A synthetic dust-emission and polarization dataset of star-forming cores}
\label{sec:dataset}

This section describes how the training data for this thesis were constructed. Starting from a
magnetohydrodynamic simulation of a star-forming molecular cloud, dense cores were identified in every
simulation snapshot, a $2.2\,\mathrm{pc}$ cube around each selected core was converted into a radiative-transfer
model, and the dust temperature, grain alignment and polarized dust emission of every cube were computed with
the Monte Carlo code POLARIS. The resulting images were combined with the three-dimensional state of the gas
into a table with one row per line of sight. Table~\ref{tab:pipeline} summarises the five steps.

\begin{table}[ht]
  \centering
  \caption{The five steps of the pipeline. Each step is a separate script; versions refer to the
  version strings recorded in each step's output.}
  \label{tab:pipeline}
  \small
  \begin{tabular}{@{}llll@{}}
    \toprule
    Step & Task & Product & Code (version) \\
    \midrule
    1 & Dense cores in every snapshot   & 11\,599 cores in 51 snapshots & \texttt{find\_starforge\_cores\_dendrogram\_v5.py} (5.2.0) \\
    2 & Non-overlapping boxes, families & 1908 boxes, 317 families     & \texttt{select\_boxes.py} \\
    3 & POLARIS grids and commands      & 1908 grids of $64^3$ cells    & \texttt{prepare\_grids.py} (3.5.1) \\
    4 & Radiative transfer              & 1908 runs, 217\,512 images    & \texttt{run\_polaris.py}, POLARIS \\
    5 & Sightline table                 & 23\,445\,504 rows, 1443 columns & \texttt{build\_dataset.py} (1.1.1 / 1.2.0) \\
    \bottomrule
  \end{tabular}
\end{table}

% -----------------------------------------------------------------------------------------------------
\subsection{The STARFORGE simulation}
\label{sec:dataset:sim}

The gas data come from the STARFORGE \citep{Grudic2021} simulation \texttt{M2e5\_R30\_Res585}, run with the
meshless magnetohydrodynamics code GIZMO \citep{Hopkins2015} (header \texttt{GIZMO\_version} 2022). The initial
condition is a uniform-density spherical cloud of $2\times10^{5}\,M_\odot$ and radius $30\,\mathrm{pc}$ (measured:
$1.993\times10^{5}\,M_\odot$ inside $R=30\,\mathrm{pc}$, median density $1.19\times10^{-22}\,\mathrm{g\,cm^{-3}}$)
embedded in a periodic box of side $300\,\mathrm{code\ units}=300.42\,\mathrm{pc}$. The initial three-dimensional
velocity dispersion of the cloud is $5.86\,\mathrm{km\,s^{-1}}$, which corresponds to a virial parameter
$\alpha=2$; the initial gas temperature is $\approx12\,\mathrm{K}$; gas cells have a median mass of
$1.001\times10^{-3}\,M_\odot$ (range $1.001$--$1.301\times10^{-3}\,M_\odot$). The interstellar radiation field
strength parameter in the snapshot header is $1$; solar abundances are $Z_\odot=0.0142$, $Y_\odot=0.2703$.

\paragraph{Units.} The code units are $1.0014\,\mathrm{pc}$ (length), $1.0008\,M_\odot$ (mass) and
$1\,\mathrm{m\,s^{-1}}$ (velocity), so one code time unit is $979.16\,\mathrm{Myr}$. The units were verified
against the initial conditions (\texttt{check\_snapshot\_units.py}): with these units the gravitational constant
derived from the header units agrees with the header value to $3.4\times10^{-4}$, the enclosed mass and density
match the nominal cloud, and the velocity dispersion matches $\alpha=2$. The saved \texttt{MagneticField} is in
tesla, \emph{not} in the unit given by the header field \texttt{Internal\_UnitB\_In\_Gauss}
($2.91\times10^{-9}\,\mathrm{G}$): read in tesla, the initial cloud field is $7.19\,\mu\mathrm{G}$, consistent
(within a factor $1.26$, the uncertainty of the critical mass-to-flux normalisation) with the nominal
normalised mass-to-flux ratio $4.2$, whereas the gauss or header-unit readings are off by four or more
orders of magnitude. POLARIS expects the field in tesla, so the field is passed on unconverted.

\paragraph{Snapshots.} Snapshots \texttt{snapshot\_000} to \texttt{snapshot\_500} in steps of 10 were analysed
(51 snapshots, each $\approx 31\,\mathrm{GB}$ with $\approx 2.5\times10^{8}$ gas cells). Dense cores exist from
snapshot 130 onwards. Snapshots 130--420 are $0.0406\,\mathrm{Myr}$ apart (0.528--1.707\,Myr); snapshots
420--500 are $0.406\,\mathrm{Myr}$ apart (1.707--4.957\,Myr). Sink particles (stars) exist only in the late
snapshots.

% -----------------------------------------------------------------------------------------------------
\subsection{Step 1: dense cores}
\label{sec:dataset:cores}

\paragraph{Molecular density.} For every gas cell the molecular hydrogen number density is
\begin{equation}
  n_{\mathrm{H_2}} = \frac{\rho\, X_\mathrm{H}\, f_\mathrm{neutral}\, f_\mathrm{mol}}{2\,m_\mathrm{H}},\qquad
  X_\mathrm{H} = 1 - Z - Y,
  \label{eq:nh2}
\end{equation}
where $\rho$ is the cell density, $Z$ and $Y$ are the first two columns of \texttt{Metallicity},
$f_\mathrm{neutral}$ is \texttt{NeutralHydrogenAbundance}, $f_\mathrm{mol}$ is \texttt{MolecularMassFraction},
and $m_\mathrm{H}=1.6735575\times10^{-24}\,\mathrm{g}$. Invalid or out-of-range inputs abort the run instead of
being clipped.

\paragraph{Dendrogram.} Cells with $n_{\mathrm{H_2}} > 10^{4}\,\mathrm{cm^{-3}}$ were selected. Two selected cells
are neighbours if either lies within the kernel support radius $h$ (\texttt{SmoothingLength}) of the other,
using periodic distances; the resulting graph was symmetrised. A dendrogram \citep{Rosolowsky2008} was computed
with \texttt{astrodendro} 0.3.1 on this particle graph with
$\texttt{min\_value}=10^{4}\,\mathrm{cm^{-3}}$, $\texttt{min\_delta}=10^{4}\,\mathrm{cm^{-3}}$ and
$\texttt{min\_npix}=100$ cells. Every leaf of the dendrogram, i.e.\ a local maximum of $n_{\mathrm{H_2}}$
with at least 100 cells and a contrast of at least $10^{4}\,\mathrm{cm^{-3}}$ above the level at which it
merges with a neighbour, is a \emph{core}. A leaf is a density peak; it is not required to be gravitationally
bound or associated with a star.

\paragraph{Core properties.} For each core the catalogue records, among others, the peak cell (position,
$n_{\mathrm{H_2}}$, particle ID), mass, volume and radii, shape, internal velocity dispersion, mean gas and dust
temperature, mass-weighted field, the gas-only virial ratio $2(K+U)/|W|$ (direct pair summation with
$10^{-3}\,\mathrm{pc}$ Plummer softening, up to $2\times10^{4}$ cells), and the sink particles near the core,
including the number of sinks within a $2.2\,\mathrm{pc}$ cube centred on the peak.
Step 1 found 11\,599 cores in total (none before snapshot 130; 1038 in snapshot 500) with a total mass of
$4535\,M_\odot$.

% -----------------------------------------------------------------------------------------------------
\subsection{Step 2: boxes and families}
\label{sec:dataset:boxes}

\paragraph{Box selection.} A \emph{box} is an axis-aligned cube of side $S=2.2\,\mathrm{pc}$ centred on a core's
density peak. Within each snapshot, boxes were selected by greedy non-maximum suppression: cores were sorted by
mass (ties, which are common because core masses are multiples of the cell mass, broken by higher peak density,
then lower \texttt{core\_id}), and a core was kept unless its cube overlapped an already kept cube. Two cubes
overlap if the periodic Chebyshev distance between their centres is $<S$. This gives 1908 boxes in 38 snapshots
(130--500; 1 to 104 boxes per snapshot, median 55). The central cores have masses of $0.10$--$9.09\,M_\odot$
(median $0.41\,M_\odot$) and peak densities $n_{\mathrm{H_2}}=2.0\times10^{4}$--$7.0\times10^{9}\,\mathrm{cm^{-3}}$
(median $3.9\times10^{4}\,\mathrm{cm^{-3}}$). Of all cores, 3604 ($1249\,M_\odot$) lie in no box: they were
suppressed by a more massive neighbour without lying inside its cube.

\paragraph{Families.} Boxes of consecutive snapshots largely image the same gas. Two boxes in consecutive
snapshots (consecutive among snapshots that contain boxes) were therefore linked if their cubes overlap
(periodic Chebyshev distance $<S$), and the connected components of these links, found with a union--find
structure, were labelled \texttt{family\_id}. The 1908 boxes form 317 families; the median family has 2 boxes,
the 90th percentile 9.4, the 99th percentile 107, and the largest family has 260 boxes. Six families have more
than 50 boxes. Families are the unit of the train/validation/test split (Sect.~\ref{sec:dataset:splits}).

% -----------------------------------------------------------------------------------------------------
\subsection{Step 3: radiative-transfer models}
\label{sec:dataset:grids}

\paragraph{Grid geometry.} Each box was converted into a POLARIS octree grid that is a uniform cube of
$64^{3}$ cells (octree level 6) of size $\Delta=2.2/64=0.034375\,\mathrm{pc}$. The grid is shifted by half a
cell so that the core's peak cell lies at the centre of cell $(32,32,32)$; centring the cube exactly on the peak
would place the peak on a corner shared by eight cells. Each snapshot was read once and the gas of all its boxes
was cut out in that single pass; boxes were then built in parallel.

\paragraph{Deposition.} Every gas cell whose kernel overlaps the cube (centre within $S/2+h$ of the box centre
along each axis) was deposited with GIZMO's compact cubic-spline kernel,
\begin{equation}
  W(q) \propto \begin{cases} 1-6q^{2}+6q^{3}, & 0\le q<\tfrac12,\\ 2(1-q)^{3}, & \tfrac12\le q<1,\\ 0, & q\ge 1,\end{cases}
  \qquad q = \frac{|\mathbf{x}_c-\mathbf{x}_i|}{h_i},
\end{equation}
sampled at the centres $\mathbf{x}_c$ of all grid cells within the kernel, \emph{including cells outside the cube}.
The weights $w_{ic}=W(q_{ic})/\sum_{c'}W(q_{ic'})$ are normalised over the full kernel and only the in-cube part
is kept, so each grid holds exactly the mass inside the cube and the faces are not under-dense. If a kernel
contains no cell centre ($h_i$ well below $\Delta$), the cell's mass goes to the grid cell containing its centre
(this applied to 1.07\% of the $2.73\times10^{8}$ deposited gas cells). Cell quantities are
\begin{equation}
  \rho_c = \frac{\sum_i m_i w_{ic}}{\Delta^{3}},\qquad
  \langle X\rangle_c = \frac{\sum_i m_i w_{ic} X_i}{\sum_i m_i w_{ic}},
\end{equation}
with $X$ the dust temperature, gas temperature and the three field components (mass-weighted means). A cell
receiving no mass would get a density floor of $10^{-6}$ times the smallest occupied density and the fields of
the nearest occupied cell; this happened for a single cell in a single box ($3.6\times10^{-14}\,M_\odot$).
The grid mass agrees with the mass of the gas cells centred inside the cube to between $-0.12\%$ and $+0.08\%$.
Boxes contain $3.5\times10^{4}$--$7.9\times10^{5}$ contributing gas cells (median $1.3\times10^{5}$) and
$27.6$--$748.6\,M_\odot$ of gas (median $113\,M_\odot$).

\paragraph{Grid fields.} The grid stores (POLARIS field IDs in brackets) the gas number density
$n=\rho/(\mu m_u)$ in $\mathrm{m^{-3}}$ [0], the STARFORGE dust temperature [2], the gas temperature [3] and
$B_x, B_y, B_z$ in tesla [4--6]. The POLARIS branch used reads the gas density only as a number density and
converts it back with $\rho=n\,\mu\,m_u$, where $m_u=1.66053904\times10^{-27}\,\mathrm{kg}$ (POLARIS's
\texttt{m\_H}); the same $n$ and $\mu$ enter the radiative-torque calculation (gas damping, thermal speed).
$\mu$ is therefore the mean mass per gas particle. We use $\mu=2.33$: the mass-weighted median of the particle
mean molecular weight of box gas, computed from the snapshot chemistry, is 2.32--2.35 for
$n_\mathrm{H}=10^{3}$--$10^{6}\,\mathrm{cm^{-3}}$ (diagnostic \texttt{check\_mu.py}, snapshot 500), and
$\rho$, hence the dust mass, is reproduced exactly for any $\mu$.

\paragraph{Dust.} One dust component was used: astrodust \citep{Hensley2023} without iron inclusions
($f_\mathrm{Fe}=0$), with the catalogue's magnetic susceptibility $\chi=4.47\times10^{-4}$ and material density
$2898\,\mathrm{kg\,m^{-3}}$, a power-law size distribution $\mathrm{d}n/\mathrm{d}a\propto a^{-3.5}$
\citep{Mathis1977} between $a_\mathrm{min}=5\,\mathrm{nm}$ and $a_\mathrm{max}=250\,\mathrm{nm}$, and a
dust-to-gas mass ratio of $0.01$ (POLARIS \texttt{<mass\_fraction>}). The ratio is the one STARFORGE itself
uses: its field \texttt{DustToGasRatio\_Local} stores $Z$ times the local dust-to-metals ratio relative to solar,
with a median of $0.01416=Z_\odot$ over all boxes, which corresponds to an absolute dust-to-gas ratio
$0.01\,Z/Z_\odot = 0.01$. The catalogue covers $1\,\mathrm{nm}$--$1\,\mu\mathrm{m}$ in size and
$0.05\,\mu\mathrm{m}$--$5\,\mathrm{mm}$ in wavelength, i.e.\ POLARIS's internal wavelength grid
(100 points, $0.0912$--$2000\,\mu\mathrm{m}$) and all detector wavelengths. With this susceptibility the
Larmor-precession limit of POLARIS's alignment model is never reached: the smallest Larmor-limited grain size
over all $5.0\times10^{8}$ grid cells is $408\,\mathrm{nm}>a_\mathrm{max}$ (\texttt{check\_larmor\_limit.py}).

\paragraph{Radiation sources.} Every box is illuminated by the interstellar radiation field (ISRF) of
\citet{Mathis1983} with strength $G_0=1$, launched from a sphere of 10 times the circumscribed radius of the grid
($10\times\tfrac{\sqrt3}{2}\times2.2\,\mathrm{pc}=19.1\,\mathrm{pc}$). In addition, every sink particle inside the
cube with positive radius and luminosity was added as a blackbody point source at its position, with the
STARFORGE protostellar radius $R_\ast$ and an effective temperature from
$L_\ast = 4\pi R_\ast^{2}\sigma_\mathrm{SB}T_\mathrm{eff}^{4}$. Fifty-two boxes contain stars, 275 sources in total
(up to 39 in one box): $L_\ast=6.8\times10^{-3}$--$1.0\times10^{4}\,L_\odot$ (median $0.16\,L_\odot$),
$R_\ast=0.31$--$20.8\,R_\odot$, $T_\mathrm{eff}=2990$--$27\,900\,\mathrm{K}$, sink masses
$0.02$--$11.7\,M_\odot$. (Step 1 counted 273 sinks in these boxes; in one box, 1851, the half-cell shift of the
grid moves two further sinks inside the cube.) Boxes without stars are heated by the ISRF only.

\paragraph{Photon numbers and Monte Carlo noise.} With radiative-torque alignment, POLARIS launches the
requested number of photons for each of its 100 wavelengths. Production used $10^{7}$ ISRF photons and $10^{6}$
photons per star (per wavelength). These values were chosen from two convergence studies, each running every
configuration twice with independent random numbers and taking the per-cell noise as
$\mathrm{median}\,|x_1-x_2|/\sqrt2$:
(i) ISRF photons $10^{6}$--$10^{7}$ and launch radii 3--30 on box \texttt{snapshot\_440/core\_000395}; at $10^{7}$
photons and radius 10 the dust-temperature noise is $0.0078\,\mathrm{K}$ in the interior, $0.0087\,\mathrm{K}$ on
faces, $0.0103\,\mathrm{K}$ on edges and $0.0137\,\mathrm{K}$ in corners ($0.0081\,\mathrm{K}$ overall), the
alignment radius varies by $0.63\%$, and at $850\,\mu\mathrm{m}$ the noise is $0.026\%$ in $I$, $0.20\%$ in $p$
and $0.035^\circ$ in the polarization angle. Smaller radii concentrate the noise in the corners
($0.073\,\mathrm{K}$ at radius 3), larger radii raise it in the interior ($0.0105\,\mathrm{K}$ at radius 30).
(ii) Star photons $10^{5}$ vs.\ $10^{6}$ on box \texttt{snapshot\_500/core\_000873}: $10^{6}$ halves the
temperature noise in gas below $25\,\mathrm{K}$ (from 0.077--0.084 to 0.039--0.045\,K) and changes the
$53\,\mu\mathrm{m}$ intensity by a median of $2.0\%$ ($0.22\%$ at $850\,\mu\mathrm{m}$).

\paragraph{POLARIS tasks.} Each box has a command file with two tasks, both with radiative-torque alignment
(\texttt{ALIG\_RAT}) and 8 threads:
\begin{enumerate}
  \item \texttt{CMD\_TEMP\_RAT}: Monte Carlo radiative transfer of the ISRF and stellar photons; computes the
  equilibrium dust temperature of every cell and the radiative-torque alignment radius $a_\mathrm{alig}$ (grains
  larger than $a_\mathrm{alig}$ are aligned with the magnetic field). The STARFORGE dust temperature in the input
  grid is replaced by this temperature. Writes \texttt{temp/grid\_temp\_rat.dat}.
  \item \texttt{CMD\_DUST\_EMISSION} on that grid: ray-traced thermal emission of the aligned grains, giving the
  Stokes $I,Q,U,V$, the optical depth and the column density for 38 wavelengths and 3 viewing directions
  (114 detectors) with $64\times64$ pixels each (\texttt{max\_subpixel\_lvl} 0). The detector covers the grid, so
  one pixel corresponds to one grid column ($0.034375\,\mathrm{pc}$); the source distance is
  $140\,\mathrm{pc}$, so a pixel subtends $50.6''$ and
  $\Omega_\mathrm{pix}=(\Delta/d)^{2}=6.03\times10^{-8}\,\mathrm{sr}$.
\end{enumerate}
The 38 wavelengths (Table~\ref{tab:bands}) are those of common far-infrared to millimetre instruments and are
read from one table (\texttt{wavelength\_telescope\_mapping.csv}) by both the grid preparation and the dataset
builder. The three viewing directions are the POLARIS detector rotations
$(\vartheta_1,\vartheta_2)=(0^\circ,0^\circ)$, $(-90^\circ,0^\circ)$ and $(0^\circ,90^\circ)$, which look along the
grid $z$, $y$ and $x$ axes (Sect.~\ref{sec:dataset:geometry}). POLARIS was the \texttt{experimental} branch,
commit \texttt{6c5e1a89c6e14ce806c47bbb87947b98e3f761b0} \citep{Reissl2016}.

% -----------------------------------------------------------------------------------------------------
\subsection{Step 4: radiative-transfer runs}
\label{sec:dataset:runs}

All 1908 boxes were run between 2026-10-31 22:47 and 2026-11-09 12:41 on one machine (AMD EPYC 9554P, 64 cores
/ 128 hardware threads, $377\,\mathrm{GB}$ memory), each box in its own directory, at the lowest scheduling
priority (\texttt{nice 19}), with up to 16 boxes at a time (8 threads each). The runner measured the CPU use of
other users every $15\,\mathrm{s}$ and kept 8 hardware threads free above their peak use over the previous
10 minutes; when the machine was needed, the most recently started boxes were suspended (\texttt{SIGSTOP}) and
resumed 10 minutes after the demand had gone, so no run was ever killed or restarted. Boxes containing stars
were run first. A box took $0.92$--$5.98\,\mathrm{h}$ of wall-clock time (median $1.46\,\mathrm{h}$ at up to 14
concurrent boxes), $2810$ box-hours in total, of which boxes spent $5.3\,\mathrm{h}$ suspended.

Every run was validated before being accepted: the temperature grid exists and all dust temperatures are finite
and within $2$--$3000\,\mathrm{K}$; all 114 detector files exist; $I$, $Q$ and $U$ are finite everywhere; and at
wavelengths $\geq50\,\mu\mathrm{m}$, $I>0$ in every pixel and $Q,U$ are not identically zero. All 1908 runs
finished with exit code 0 and passed; none failed. POLARIS's Monte Carlo is seeded by thread, and the assignment
of photons to threads varies between runs, so runs are statistically but not bitwise reproducible.

% -----------------------------------------------------------------------------------------------------
\subsection{Step 5: the sightline dataset}
\label{sec:dataset:table}

\subsubsection{Geometry and orientation}
\label{sec:dataset:geometry}

Grid arrays are indexed $[i_x,i_y,i_z]$ and images $[\mathrm{row},\mathrm{col}]$, both from 0 to 63. The mapping
between them was determined empirically (\texttt{check\_view\_geometry.py}): for each view, the POLARIS
column-density image was compared with $\sum n\,\Delta$ along each grid axis in all 48 combinations of
line-of-sight axis, transposition and flips; the best orientation matches to a relative difference of
$\approx10^{-13}$. Which end of the line of sight faces the observer cannot be seen in optically thin emission; it
was determined where the dust is optically thick ($21\,\mu\mathrm{m}$, pixels with $\tau>1$), where
$I=\sum_k B_\nu(T_k)\,\Delta\tau_k\,e^{-\tau_k}$ depends on the side from which the absorption is counted
(\texttt{check\_observer\_side.py}). The model built from POLARIS's own temperatures reproduces the image with a
near-constant ratio only if the high-index end of the line-of-sight axis faces the observer, which holds in every
view. Table~\ref{tab:views} gives the result.

\begin{table}[ht]
  \centering
  \caption{Mapping between a dataset row (view $v$, pixel row $r$, pixel column $c$, profile index $k$, with $k=0$
  the cell nearest the observer) and the grid, and the field components in the image frame.}
  \label{tab:views}
  \small
  \begin{tabular}{@{}cccccc@{}}
    \toprule
    View & POLARIS rotation & Line of sight & Image $+$col, $+$row & Grid cell $(i_x,i_y,i_z)$ & $(B_\mathrm{col},B_\mathrm{row},B_\mathrm{los})$ \\
    \midrule
    0 & $(0^\circ,0^\circ)$   & $z$ & $+x$, $+y$ & $(c,\ r,\ 63-k)$    & $(B_x,\ B_y,\ -B_z)$ \\
    1 & $(-90^\circ,0^\circ)$ & $y$ & $+x$, $-z$ & $(c,\ 63-k,\ 63-r)$ & $(B_x,\ -B_z,\ -B_y)$ \\
    2 & $(0^\circ,90^\circ)$  & $x$ & $-z$, $+y$ & $(63-k,\ r,\ 63-c)$ & $(-B_z,\ B_y,\ -B_x)$ \\
    \bottomrule
  \end{tabular}
\end{table}

Angles in the image plane are measured from the image $+$col axis towards $+$row. With Stokes $U$ as written by
POLARIS, the plane-of-sky field direction is the polarization angle plus $90^\circ$: comparing
$\tfrac12\arctan(U,Q)+90^\circ$ with a density-weighted Stokes model of the grid field in this frame (in the more
polarized half of the pixels) gives median differences of $0.2^\circ$--$0.8^\circ$, against $22^\circ$--$54^\circ$ if the
sign of $U$ is reversed. $B_\mathrm{los}>0$ points away from the
observer. Every box and view was checked during the build: $\sum_k n_k\Delta$ along the rotated profiles
reproduces POLARIS's column-density image to a maximum relative difference of $5.4\times10^{-13}$
(the build aborts above $10^{-5}$).

\subsubsection{Contents}

One row is one line of sight: one pixel of one view of one box. The dataset has
$1908\times3\times64\times64=23\,445\,504$ rows and 1443 columns (Table~\ref{tab:columns}).

\begin{table}[ht]
  \centering
  \caption{Columns of the dataset. $k=00\ldots63$ runs along the line of sight from the observer; columns
  marked (+log) also exist as \texttt{<name>\_log} $=\log_{10}$ (NaN where the value is $\le0$). Values are
  \texttt{float32} except $I,Q,U$ (\texttt{float64}) and the integer metadata.}
  \label{tab:columns}
  \small
  \begin{tabular}{@{}p{0.17\textwidth}p{0.47\textwidth}p{0.12\textwidth}r@{}}
    \toprule
    Group & Columns & Unit & Number \\
    \midrule
    metadata & \texttt{snapshot, time\_Myr, core\_id, box\_id, family\_id, n\_stars, split, view, pix\_row, pix\_col} & & 10 \\
    sightline & \texttt{N\_col, Sigma, Sigma\_dust} (+log) & $\mathrm{cm^{-2}}$, $\mathrm{g\,cm^{-2}}$ & 6 \\
              & \texttt{temp\_polaris\_mean, temp\_starforge\_mean} (+log) & K & 4 \\
              & \texttt{B\_tot\_mean} (+log) & $\mu$G & 2 \\
              & \texttt{cos2psi\_B\_mean, sin2psi\_B\_mean, B\_coherence} & -- & 3 \\
    bands ($\times38$) & \texttt{<band>\_I, \_Q, \_U} & MJy\,sr$^{-1}$ & 114 \\
              & \texttt{<band>\_p, \_polangle, \_I\_log, \_p\_log} & --, deg & 152 \\
    profiles ($\times64$) & \texttt{dens\_k, dens\_dust\_k} (+log) & $\mathrm{cm^{-3}}$, $\mathrm{g\,cm^{-3}}$ & 256 \\
              & \texttt{temp\_polaris\_k, temp\_starforge\_k} (+log) & K & 256 \\
              & \texttt{B\_tot\_k, B\_pos\_k} (+log) & $\mu$G & 256 \\
              & \texttt{cos2psi\_B\_k, sin2psi\_B\_k, sin2gamma\_B\_k} & -- & 192 \\
              & \texttt{B\_col\_k, B\_row\_k, B\_los\_k} & $\mu$G & 192 \\
    \midrule
    total & & & 1443 \\
    \bottomrule
  \end{tabular}
\end{table}

\paragraph{Profiles.} For each cell $k$ along the line of sight:
$\texttt{dens}=n$ in $\mathrm{cm^{-3}}$ (gas particles, $\mu=2.33$);
$\texttt{dens\_dust}=0.01\,n\,\mu\,m_u$ in $\mathrm{g\,cm^{-3}}$;
$\texttt{temp\_polaris}$, the dust temperature computed by POLARIS; $\texttt{temp\_starforge}$, the deposited
STARFORGE dust temperature; the field in $\mu$G ($1\,\mathrm{T}=10^{10}\,\mu\mathrm{G}$) in the image frame
(Table~\ref{tab:views}) and the derived quantities
\begin{gather}
  B_\mathrm{pos}^{2}=B_\mathrm{col}^{2}+B_\mathrm{row}^{2},\qquad B_\mathrm{tot}^{2}=B_\mathrm{pos}^{2}+B_\mathrm{los}^{2},\\
  \cos2\psi_B=\frac{B_\mathrm{col}^{2}-B_\mathrm{row}^{2}}{B_\mathrm{pos}^{2}},\qquad
  \sin2\psi_B=\frac{2B_\mathrm{col}B_\mathrm{row}}{B_\mathrm{pos}^{2}},\qquad
  \sin^{2}\gamma_B=\frac{B_\mathrm{pos}^{2}}{B_\mathrm{tot}^{2}},
\end{gather}
where $\psi_B$ is the plane-of-sky field angle and $\gamma_B$ the angle between the field and the line of sight.
The double-angle form removes the $180^\circ$ ambiguity of the field direction, which dust polarization cannot
measure.

\paragraph{Sightline quantities.} $N_\mathrm{col}$ is POLARIS's column-density image converted from
$\mathrm{m^{-2}}$ to $\mathrm{cm^{-2}}$; $\Sigma=N_\mathrm{col}\,\mu m_u$ and $\Sigma_\mathrm{dust}=0.01\,\Sigma$.
Means are density-weighted, $\langle X\rangle=\sum_k n_kX_k/\sum_k n_k$, for both dust temperatures and
$B_\mathrm{tot}$. The field orientation that a polarization map measures is described by
\begin{equation}
  z = \sum_k n_k\,\sin^{2}\gamma_{B,k}\;e^{2i\psi_{B,k}},\qquad
  (\texttt{cos2psi\_B\_mean},\texttt{sin2psi\_B\_mean})=\frac{(\operatorname{Re}z,\operatorname{Im}z)}{|z|},\qquad
  \texttt{B\_coherence}=\frac{|z|}{\sum_k n_k}\in[0,1],
  \label{eq:coherence}
\end{equation}
so \texttt{B\_coherence} combines field tangling along the line of sight and inclination; it correlates with the
$850\,\mu\mathrm{m}$ polarization fraction with $r=0.87$--$0.94$ (Pearson, per file, in the eight files checked
by \texttt{validate\_dataset.py}).

\paragraph{Bands.} POLARIS returns $I,Q,U$ in Jy per pixel; they are converted to $\mathrm{MJy\,sr^{-1}}$ by
dividing by $10^{6}\,\Omega_\mathrm{pix}$. The polarization fraction is $p=\sqrt{Q^{2}+U^{2}}/I$ (NaN where
$I\le0$) and \texttt{polangle} $=\tfrac12\arctan(U,Q)$ in degrees, mapped to $(-90^\circ,90^\circ]$; it is the
E-vector angle, and the inferred field angle is $\texttt{polangle}+90^\circ$. $I,Q,U$ are stored in double
precision because at $3.4$--$5.8\,\mu\mathrm{m}$ the intensities reach $10^{-108}\,\mathrm{MJy\,sr^{-1}}$, below
the single-precision range. POLARIS was run without stochastic heating of small grains, so the 14 bands at
$\le25\,\mu\mathrm{m}$ are flagged \texttt{reliable: false} in the manifest; the 24 bands at
$53\,\mu\mathrm{m}$--$1.38\,\mathrm{mm}$ are reliable. The images contain no observational noise and no beam.

\begin{table}[ht]
  \centering
  \caption{The 38 detector wavelengths and band names (column prefixes). Bands at $\le25\,\mu\mathrm{m}$ are
  flagged unreliable.}
  \label{tab:bands}
  \footnotesize
  \begin{tabular}{@{}rl rl@{}}
    \toprule
    $\lambda$ [$\mu$m] & Band & $\lambda$ [$\mu$m] & Band \\
    \midrule
    3.4   & \texttt{wise\_1}            & 89    & \texttt{sofia\_hawc\_plus\_c} \\
    3.6   & \texttt{spitzer\_irac\_1}   & 90    & \texttt{akari\_fis\_90} \\
    4.5   & \texttt{spitzer\_irac\_2}   & 100   & \texttt{herschel\_pacs\_2\_or\_iras\_100} \\
    4.6   & \texttt{wise\_2}            & 140   & \texttt{akari\_fis\_140} \\
    5.8   & \texttt{spitzer\_irac\_3}   & 154   & \texttt{sofia\_hawc\_plus\_d} \\
    8.0   & \texttt{spitzer\_irac\_4}   & 160   & \texttt{herschel\_pacs\_3} \\
    8.28  & \texttt{msx\_1}             & 214   & \texttt{sofia\_hawc\_plus\_e} \\
    12.0  & \texttt{wise\_3\_or\_iras\_12} & 250 & \texttt{herschel\_spire\_1} \\
    12.13 & \texttt{msx\_2}             & 350   & \texttt{herschel\_spire\_2\_or\_alma\_band\_10\_c43\_8\_or\_planck\_hfi\_857ghz} \\
    14.65 & \texttt{msx\_3}             & 460   & \texttt{alma\_band\_9\_c43\_8} \\
    21.34 & \texttt{msx\_4}             & 500   & \texttt{herschel\_spire\_3} \\
    22.0  & \texttt{wise\_4}            & 550   & \texttt{planck\_hfi\_545ghz} \\
    24.0  & \texttt{spitzer\_mips\_24}  & 690   & \texttt{alma\_band\_8\_c43\_8} \\
    25.0  & \texttt{iras\_25}           & 849   & \texttt{planck\_hfi\_353ghz} \\
    53    & \texttt{sofia\_hawc\_plus\_a} & 870 & \texttt{apex\_laboca} \\
    60    & \texttt{iras\_60}           & 945   & \texttt{alma\_band\_7\_c43\_8} \\
    63    & \texttt{sofia\_hawc\_plus\_b} & 1100 & \texttt{cso\_bolocam} \\
    65    & \texttt{akari\_fis\_65}     & 1300  & \texttt{alma\_band\_6\_c43\_8} \\
    70    & \texttt{herschel\_pacs\_1}  & 1380  & \texttt{planck\_hfi\_217ghz} \\
    \bottomrule
  \end{tabular}
\end{table}

\subsubsection{Train, validation and test splits}
\label{sec:dataset:splits}

Boxes in the same family image largely the same gas a short time apart, so splitting by box or by pixel would put
near-identical data into training and test sets. Whole families were therefore assigned to one split, and all
three views of a box share its split. Families were shuffled with a fixed seed (NumPy \texttt{default\_rng},
seed 20261111) and assigned in turn to the split furthest below its target share (80/10/10\% of the boxes) among
the splits the family fits into; families with more than 50 boxes always go to the training set, so that no single
region dominates a validation or test set. Two versions of the dataset were built from identical data
(Table~\ref{tab:splits}):

\begin{itemize}
  \item \texttt{starforge\_polaris\_dataset\_v1} (builder 1.1.1): the split described above. All 52 boxes with
  stars belong to families assigned to the training set.
  \item \texttt{starforge\_polaris\_dataset\_v1\_stars\_eval} (builder 1.2.0, option \texttt{--assign}): as v1, but
  the four small star-forming families 41 and 226 are moved to validation and 133 and 139 to test; every other box
  keeps its split. Star-forming families 140 and 280 stay in training because they lie only $3.3$ and
  $2.8\,\mathrm{pc}$ from the large training family 2. The non-star validation and test families are identical in
  both versions, so models trained on either can be compared on that common subset.
\end{itemize}

\begin{table}[ht]
  \centering
  \caption{Splits of the two dataset versions (boxes; each box gives $3\times4096$ rows).}
  \label{tab:splits}
  \small
  \begin{tabular}{@{}lrrrrrr@{}}
    \toprule
     & \multicolumn{3}{c}{\texttt{v1}} & \multicolumn{3}{c}{\texttt{v1\_stars\_eval}} \\
    \cmidrule(lr){2-4}\cmidrule(l){5-7}
    Split & boxes & families & boxes with stars & boxes & families & boxes with stars \\
    \midrule
    train      & 1538 (80.6\%) & 207 & 52 & 1458 (76.4\%) & 203 & 36 \\
    validation &  187 (9.8\%)  &  58 &  0 &  216 (11.3\%) &  60 &  8 \\
    test       &  183 (9.6\%)  &  52 &  0 &  234 (12.3\%) &  54 &  8 \\
    \bottomrule
  \end{tabular}
\end{table}

In v1 the largest family makes up 17\% of the training set, 16\% of the validation set and 21\% of the test
set; all three splits span 0.5--5.0\,Myr. All star-forming families lie within the same central region of the
cloud and appear at late times (snapshots 430--500), with their nearest boxes 2.6--4.8\,pc apart; scores on the
star-forming test families therefore measure performance on new cores of that region rather than of a different
cloud.

% -----------------------------------------------------------------------------------------------------
\subsection{The dataset}
\label{sec:dataset:contents}

\subsubsection{What one row contains}

Each row describes one line of sight through one box, seen in one of the three viewing directions. It holds
\begin{enumerate}
  \item \emph{identifiers} that locate the sightline in the simulation (snapshot, time, core, box, family), in the
  split, and in its image (view, pixel row and column);
  \item the \emph{observables} a telescope would measure along this sightline: Stokes $I$, $Q$, $U$, the
  polarization fraction and the polarization angle in 38 bands;
  \item \emph{sightline targets}: quantities integrated or averaged along the line of sight (column and surface
  density, mean dust temperatures, mean field strength, field orientation and coherence);
  \item \emph{profile targets}: the three-dimensional state of the 64 grid cells along the line of sight, ordered
  from the observer: gas and dust density, the two dust temperatures and the magnetic field (strength, orientation
  and components).
\end{enumerate}
Observables and targets of the same row refer to exactly the same column of grid cells, and the profiles are
already in the observer's frame (Table~\ref{tab:views}), so no geometric transformation is needed when using the
data. Because each row group holds one box and view, the rows of a row group reassemble into a $64\times64$ image
with \texttt{pix\_row} and \texttt{pix\_col}, and the profile columns of that image into a $64\times64\times64$
cube; the dataset therefore serves sightline-based (per-pixel) as well as image-based models.

\subsubsection{Files}

The dataset is stored in Apache Parquet format (written by \texttt{pyarrow}, zstd compression); Table~\ref{tab:files}
lists the files. Each split directory holds one file per snapshot that has boxes in that split, and each file holds
one row group of 4096 rows per box and view, so a single box and view can be read without reading the rest of the
file. The data occupy $178\,\mathrm{GB}$ ($166\,\mathrm{GiB}$) per dataset version. Building them from the
POLARIS output takes 27 minutes with 16 processes.

\begin{table}[ht]
  \centering
  \caption{Files of a dataset version (sizes and counts for \texttt{v1}).}
  \label{tab:files}
  \small
  \begin{tabular}{@{}p{0.27\textwidth}p{0.52\textwidth}r@{}}
    \toprule
    Path & Content & Size \\
    \midrule
    \texttt{data/train/<snapshot>.parquet} & 37 files; 1538 boxes, 4614 row groups, 18\,898\,944 rows & 143.5\,GB \\
    \texttt{data/val/<snapshot>.parquet}   & 38 files; 187 boxes, 561 row groups, 2\,297\,856 rows & 17.4\,GB \\
    \texttt{data/test/<snapshot>.parquet}  & 33 files; 183 boxes, 549 row groups, 2\,248\,704 rows & 17.1\,GB \\
    \texttt{splits.csv} & one line per box: \texttt{box} (\texttt{snapshot\_NNN/core\_NNNNNN}), \texttt{family\_id}, \texttt{split} & \\
    \texttt{manifest.json} & builder version and path, creation time, input paths, numbers of boxes, rows and columns,
      geometry conventions, units, the band table with wavelength and \texttt{reliable} flag, split seed, target fractions,
      sizes (and, in \texttt{v1\_stars\_eval}, the family overrides), the largest column-density mismatch, and a QA record per box
      (column-density check per view, non-positive intensities per band) & 612\,kB \\
    \bottomrule
  \end{tabular}
\end{table}

\subsubsection{Columns}

Table~\ref{tab:dictionary} defines every column. In the column names, \texttt{<b>} stands for one of the 38 band
names of Table~\ref{tab:bands} and \texttt{<k>} for the two-digit cell index \texttt{00}\ldots\texttt{63} along
the line of sight (\texttt{00} nearest the observer). A column with suffix \texttt{\_log} contains $\log_{10}$ of the
column without it. Of the 1443 columns, 1320 are single precision (\texttt{float32}), 114 double precision
(\texttt{float64}: all $I$, $Q$, $U$), 8 integers and one (\texttt{split}) a dictionary-encoded string. The
dataset contains no missing or non-finite values: \texttt{validate\_dataset.py} scanned every value of every
column and found none, so the NaN conventions of the builder (logarithm of a non-positive value, $p$ where
$I\le0$) never apply.

\begingroup
\small
\setlength{\LTcapwidth}{\textwidth}
\begin{longtable}{@{}p{0.27\textwidth}p{0.08\textwidth}p{0.11\textwidth}p{0.46\textwidth}@{}}
  \caption{Data dictionary. Type: f32/f64 = single/double-precision float, i8/i16/i32 = integers. Counts in
  brackets give the number of columns of a pattern.}
  \label{tab:dictionary}\\
  \toprule
  Column & Type & Unit & Definition \\
  \midrule
  \endfirsthead
  \multicolumn{4}{@{}l}{\small\emph{Table~\ref{tab:dictionary} continued}}\\
  \toprule
  Column & Type & Unit & Definition \\
  \midrule
  \endhead
  \midrule
  \multicolumn{4}{r@{}}{\small\emph{continued on next page}}\\
  \endfoot
  \bottomrule
  \endlastfoot
  \multicolumn{4}{@{}l}{\emph{Identifiers} (10 columns)} \\
  \texttt{snapshot} & i16 & -- & STARFORGE snapshot number, 130--500 \\
  \texttt{time\_Myr} & f32 & Myr & simulation time of the snapshot, 0.528--4.957 \\
  \texttt{core\_id} & i32 & -- & ID of the box's central core in the step-1 catalogue of that snapshot \\
  \texttt{box\_id} & i32 & -- & unique box ID, 0--1907 (step 2) \\
  \texttt{family\_id} & i32 & -- & cross-snapshot family, 0--316 (Sect.~\ref{sec:dataset:boxes}) \\
  \texttt{n\_stars} & i16 & -- & number of stellar sources inside the box, 0--39 \\
  \texttt{split} & string & -- & \texttt{train}, \texttt{val} or \texttt{test} \\
  \texttt{view} & i8 & -- & viewing direction 0, 1, 2 (Table~\ref{tab:views}) \\
  \texttt{pix\_row}, \texttt{pix\_col} & i8 & -- & image pixel, 0--63 \\
  \midrule
  \multicolumn{4}{@{}l}{\emph{Sightline quantities} (15 columns)} \\
  \texttt{N\_col} (+\texttt{\_log}) & f32 & cm$^{-2}$ & column density of gas particles from POLARIS's column-density image \\
  \texttt{Sigma} (+\texttt{\_log}) & f32 & g\,cm$^{-2}$ & gas surface density, $N_\mathrm{col}\,\mu m_u$ \\
  \texttt{Sigma\_dust} (+\texttt{\_log}) & f32 & g\,cm$^{-2}$ & dust surface density, $0.01\,\Sigma$ \\
  \texttt{temp\_polaris\_mean} (+\texttt{\_log}) & f32 & K & density-weighted mean POLARIS dust temperature \\
  \texttt{temp\_starforge\_mean} (+\texttt{\_log}) & f32 & K & density-weighted mean STARFORGE dust temperature \\
  \texttt{B\_tot\_mean} (+\texttt{\_log}) & f32 & $\mu$G & density-weighted mean field strength $|\mathbf{B}|$ \\
  \texttt{cos2psi\_B\_mean}, \texttt{sin2psi\_B\_mean} & f32 & -- & direction of $z$ (Eq.~\ref{eq:coherence}): plane-of-sky field orientation seen in polarization \\
  \texttt{B\_coherence} & f32 & -- & $|z|/\sum_k n_k$, 0--1 (Eq.~\ref{eq:coherence}) \\
  \midrule
  \multicolumn{4}{@{}l}{\emph{Bands} (7 columns $\times$ 38 bands = 266)} \\
  \texttt{<b>\_I} & f64 & MJy\,sr$^{-1}$ & Stokes $I$ \\
  \texttt{<b>\_Q}, \texttt{<b>\_U} & f64 & MJy\,sr$^{-1}$ & Stokes $Q$, $U$ in the image frame (angles from $+$col towards $+$row) \\
  \texttt{<b>\_p} & f32 & -- & polarization fraction $\sqrt{Q^2+U^2}/I$ \\
  \texttt{<b>\_polangle} & f32 & deg & polarization (E-vector) angle $\tfrac12\arctan(U,Q)$ in $(-90,90]$; field angle $=$ \texttt{polangle} $+90^\circ$ \\
  \texttt{<b>\_I\_log}, \texttt{<b>\_p\_log} & f32 & -- & $\log_{10}$ of $I$ [MJy\,sr$^{-1}$] and of $p$ \\
  \midrule
  \multicolumn{4}{@{}l}{\emph{Profiles} (64 cells each, $k=00$ nearest the observer; 1152 columns)} \\
  \texttt{dens\_<k>} (+\texttt{\_log}) & f32 & cm$^{-3}$ & gas particle density $n=\rho/(\mu m_u)$, $\mu=2.33$ \\
  \texttt{dens\_dust\_<k>} (+\texttt{\_log}) & f32 & g\,cm$^{-3}$ & dust mass density, $0.01\,\rho$ \\
  \texttt{temp\_polaris\_<k>} (+\texttt{\_log}) & f32 & K & dust temperature computed by POLARIS (used for the images) \\
  \texttt{temp\_starforge\_<k>} (+\texttt{\_log}) & f32 & K & STARFORGE dust temperature (mass-weighted deposit) \\
  \texttt{B\_tot\_<k>} (+\texttt{\_log}) & f32 & $\mu$G & field strength $|\mathbf{B}|$ \\
  \texttt{B\_pos\_<k>} (+\texttt{\_log}) & f32 & $\mu$G & plane-of-sky field strength $\sqrt{B_\mathrm{col}^2+B_\mathrm{row}^2}$ \\
  \texttt{cos2psi\_B\_<k>}, \texttt{sin2psi\_B\_<k>} & f32 & -- & plane-of-sky field orientation $\cos2\psi_B$, $\sin2\psi_B$ \\
  \texttt{sin2gamma\_B\_<k>} & f32 & -- & $\sin^2\gamma_B=B_\mathrm{pos}^2/B_\mathrm{tot}^2$ (1: field in the sky plane; 0: along the line of sight) \\
  \texttt{B\_col\_<k>}, \texttt{B\_row\_<k>} & f32 & $\mu$G & field components along image $+$col and $+$row \\
  \texttt{B\_los\_<k>} & f32 & $\mu$G & line-of-sight component, positive away from the observer \\
\end{longtable}
\endgroup

\subsubsection{Value ranges}

Table~\ref{tab:ranges} gives the distribution of the main quantities. The sightline and band columns were
evaluated over all 23\,445\,504 rows of \texttt{v1}; the profile columns over all 64 cells of 300 randomly chosen
box views ($7.9\times10^{7}$ cells). All column densities in the dataset correspond to optically thin dust at
long wavelengths ($\tau_{850}\le0.01$, Sect.~\ref{sec:dataset:verification}). Intensities at
$\le5.8\,\mu\mathrm{m}$ are of order $10^{-108}\,\mathrm{MJy\,sr^{-1}}$ and carry no usable signal.

\begin{table}[ht]
  \centering
  \caption{Distribution of selected quantities: minimum, 1st percentile, median, 99th percentile and maximum.}
  \label{tab:ranges}
  \small
  \begin{tabular}{@{}lllllll@{}}
    \toprule
    Column & Unit & min & 1\% & median & 99\% & max \\
    \midrule
    \multicolumn{7}{@{}l}{\emph{Sightline (all rows)}} \\
    \texttt{N\_col}               & cm$^{-2}$     & $2.4\times10^{18}$ & $9.9\times10^{19}$ & $9.2\times10^{20}$ & $7.1\times10^{21}$ & $2.5\times10^{23}$ \\
    \texttt{Sigma}                & g\,cm$^{-2}$  & $9.3\times10^{-6}$ & $3.8\times10^{-4}$ & $3.6\times10^{-3}$ & $2.7\times10^{-2}$ & $0.97$ \\
    \texttt{temp\_polaris\_mean}  & K             & 8.27 & 12.5 & 14.5 & 16.0 & 54.4 \\
    \texttt{temp\_starforge\_mean}& K             & 6.16 & 11.3 & 13.3 & 15.5 & 53.8 \\
    \texttt{B\_tot\_mean}         & $\mu$G        & 0.96 & 9.6 & 25.3 & 58.3 & 1760 \\
    \texttt{B\_coherence}         & --            & $3.6\times10^{-5}$ & 0.029 & 0.47 & 0.96 & 0.999 \\
    \midrule
    \multicolumn{7}{@{}l}{\emph{Bands (all rows)}} \\
    $I$ (53\,$\mu$m)   & MJy\,sr$^{-1}$ & $3.2\times10^{-4}$ & 0.0074 & 0.018 & 0.072 & $2.1\times10^{6}$ \\
    $I$ (214\,$\mu$m)  & MJy\,sr$^{-1}$ & 0.083 & 3.3 & 20.5 & 87.4 & $1.2\times10^{5}$ \\
    $I$ (849\,$\mu$m)  & MJy\,sr$^{-1}$ & 0.0039 & 0.17 & 1.33 & 8.1 & 1380 \\
    $I$ (1380\,$\mu$m) & MJy\,sr$^{-1}$ & $8.0\times10^{-4}$ & 0.034 & 0.28 & 1.8 & 246 \\
    $p$ (53\,$\mu$m)   & --  & $1.4\times10^{-5}$ & 0.0059 & 0.098 & 0.259 & 0.355 \\
    $p$ (214\,$\mu$m)  & --  & $1.0\times10^{-5}$ & 0.0052 & 0.089 & 0.260 & 0.323 \\
    $p$ (849\,$\mu$m)  & --  & $2.2\times10^{-6}$ & 0.0051 & 0.089 & 0.261 & 0.324 \\
    $I$ (3.4\,$\mu$m)  & MJy\,sr$^{-1}$ & $7\times10^{-136}$ & $3\times10^{-112}$ & $8\times10^{-108}$ & $1\times10^{-104}$ & $3\times10^{-12}$ \\
    \midrule
    \multicolumn{7}{@{}l}{\emph{Profiles (300 box views, all cells)}} \\
    \texttt{dens}           & cm$^{-3}$    & 0.010 & 5.8 & 73.8 & 2310 & $9.9\times10^{5}$ \\
    \texttt{dens\_dust}     & g\,cm$^{-3}$ & $3.9\times10^{-28}$ & $2.3\times10^{-25}$ & $2.9\times10^{-24}$ & $8.9\times10^{-23}$ & $3.8\times10^{-20}$ \\
    \texttt{temp\_polaris}  & K            & 8.37 & 12.7 & 15.0 & 16.1 & 80.0 \\
    \texttt{temp\_starforge}& K            & 5.91 & 11.3 & 13.4 & 15.6 & 66.8 \\
    \texttt{B\_tot}         & $\mu$G       & 0.037 & 7.3 & 26.0 & 69.1 & 2380 \\
    \texttt{B\_pos}         & $\mu$G       & 0.0017 & 2.7 & 20.2 & 60.2 & 2370 \\
    \texttt{sin2gamma\_B}   & --           & $2.6\times10^{-8}$ & 0.019 & 0.77 & 1.00 & 1.00 \\
    \texttt{B\_los}         & $\mu$G       & $-1600$ & $-45.8$ & $-6.6$ & 31.8 & 1320 \\
    \bottomrule
  \end{tabular}
\end{table}

The intensity spectra of all sightlines peak between 160 and $250\,\mu\mathrm{m}$ (at $214\,\mu\mathrm{m}$ for
96.8\%), as expected for dust at 12--16\,K. The polarization fraction changes little with wavelength between
$53\,\mu\mathrm{m}$ and $1.38\,\mathrm{mm}$ (median 0.098 at $53\,\mu\mathrm{m}$, 0.089 at 214 and
$849\,\mu\mathrm{m}$; maximum 0.355).

\subsubsection{Reading the data}

Columns can be read selectively, which keeps memory small; for example, with \texttt{pandas} and \texttt{pyarrow}:
\begin{verbatim}
import pandas as pd, pyarrow.dataset as ds
root = "/scratch/sjafaritabar/starforge_polaris_dataset_v1/data"
cols = ["planck_hfi_353ghz_I_log", "planck_hfi_353ghz_p"] + \
       [f"dens_{k:02d}_log" for k in range(64)]
train = pd.read_parquet(f"{root}/train", columns=cols)
full = ds.dataset(root, format="parquet", partitioning=None)   # lazy, all splits
\end{verbatim}
To obtain images, a file is read one row group at a time (\texttt{pyarrow.parquet.ParquetFile.read\_row\_group})
and the 4096 rows are placed at (\texttt{pix\_row}, \texttt{pix\_col}).

% -----------------------------------------------------------------------------------------------------
\subsection{Verification}
\label{sec:dataset:verification}

Each step has unit tests (step 2: 7, step 3: 13, step 5: 4; all pass), and an independent audit
(\texttt{audit\_pipeline.py}) re-derives every stage from the raw files without using the pipeline's code. It
passed all 15 checks for both dataset versions:
\begin{enumerate}
  \item Every boxed core is in its step-1 catalogue with the same peak, peak particle and mass; box times equal
  the snapshot header times; no two boxes of a snapshot overlap; families re-derived from scratch reproduce all
  317 \texttt{family\_id}s.
  \item Exactly one grid per box; for every box the summaries, grids, star lists, command files (photon numbers,
  radius, dust, wavelengths, views) and POLARIS outputs are mutually consistent, the grid density is unchanged by
  POLARIS, and dust temperatures and alignment radii are within range.
  \item The $850\,\mu\mathrm{m}$ intensity in $\mathrm{MJy\,sr^{-1}}$ matches the optically thin model
  $\sum_k B_\nu(T_k)\,\Delta\tau_k$ computed from the grid, which tests the units, pixel solid angle and distance:
  POLARIS/model has a per-box median of $1.002$ (range $1.001$--$1.005$); $\tau_{850}\le0.01$.
  \item No box has an outlying POLARIS-vs-STARFORGE dust-temperature offset ($>5\sigma$).
  \item Every one of the 1443 columns was recomputed from the raw grids and FITS files with explicit index formulas
  (Table~\ref{tab:views}) for 1080 random sightlines in 36 boxes ($1.56\times10^{6}$ values); the worst
  relative (or, for logarithms and angles, absolute) error is $3.8\times10^{-6}$, i.e.\ single-precision rounding.
  \item \texttt{splits.csv} covers every box once with the step-2 family IDs; no family is in more than one split;
  all families with more than 50 boxes are in training.
  \item Polarization fractions at $\ge50\,\mu\mathrm{m}$ are physical (maximum 0.355), and the spectra of all
  sightlines peak at 60--$500\,\mu\mathrm{m}$ (214\,$\mu$m for 96.8\%).
\end{enumerate}
\texttt{validate\_dataset.py} additionally confirms the row counts, the absence of duplicate sightlines and of
boxes or families in more than one split, and scans every column of every row for non-finite values.

Across the whole dataset, column densities range over $2.4\times10^{18}$--$2.5\times10^{23}\,\mathrm{cm^{-2}}$
(median $9.1\times10^{20}$), density-weighted POLARIS dust temperatures over 8.3--54.4\,K (median 14.5\,K),
density-weighted field strengths over $0.96$--$1760\,\mu\mathrm{G}$ (median $25.7\,\mu\mathrm{G}$), and
\texttt{B\_coherence} over $4\times10^{-5}$--$0.999$ (median 0.48).

% -----------------------------------------------------------------------------------------------------
\subsection{Limitations}
\label{sec:dataset:limitations}

\begin{itemize}
  \item \textbf{Radiation field.} Each box is illuminated by an unattenuated Mathis field with $G_0=1$. In
  STARFORGE the field reaching the box surfaces is attenuated by the surrounding cloud (snapshot 440: equivalent
  $G_0$ median 0.20, 5--95\% range 0.11--0.46). POLARIS dust temperatures are therefore higher than the STARFORGE
  ones, by a per-box median of $+1.49\,\mathrm{K}$ (5--95\%: $+0.29$ to $+2.76\,\mathrm{K}$). Both temperatures
  are in the dataset.
  \item \textbf{Isolated boxes.} Only gas inside the $2.2\,\mathrm{pc}$ cube emits and absorbs; foreground and
  background cloud emission and extinction are absent.
  \item \textbf{No stochastic heating}: bands at $\le25\,\mu\mathrm{m}$ are not reliable.
  \item \textbf{Ideal observations.} No noise, no beam, a single distance (140\,pc) and a pixel equal to a grid
  cell. Noise and beam have to be added when training.
  \item \textbf{No gas kinematics.} Velocities are not included, so velocity-based field-strength estimates
  (e.g.\ Davis--Chandrasekhar--Fermi) cannot be reproduced from the dataset alone.
  \item \textbf{One simulation.} All boxes come from one cloud; families make the evaluation sets independent in
  gas, not in environment. Late snapshots, which contain all stars, are $0.41\,\mathrm{Myr}$ apart.
  \item \textbf{Monte Carlo.} Dust temperatures carry $\approx0.008\,\mathrm{K}$ of Monte Carlo noise, and runs
  are not bitwise reproducible.
\end{itemize}

% -----------------------------------------------------------------------------------------------------
\subsection{Data products and code}
\label{sec:dataset:files}

Code: \texttt{STARFORGE/creating\_the\_dataset/} (\texttt{1.finding\_cores}, \texttt{2.selecting\_boxes},
\texttt{3.preparing\_polaris\_grids} with its \texttt{diagnostics}, \texttt{4.running\_polaris},
\texttt{5.building\_dataset}, \texttt{audit\_pipeline.py}); photon studies in
\texttt{STARFORGE/M2e5\_R30\_Res585/POLARIS/isrf\_photon\_noise\_test} and \texttt{STARFORGE/M2e5\_R30\_Res585/POLARIS/star\_photon\_test}.
Data: snapshots in \texttt{/scratch/STARFORGE/M2e5\_R30\_Res585}; core catalogues in
\texttt{1.finding\_cores/dendrogram\_cores\_v5.2}; boxes in \texttt{2.selecting\_boxes/boxes\_2.2pc};
grids in \texttt{3.preparing\_polaris\_grids/grids\_64}; POLARIS output in
\texttt{/scratch/sjafaritabar/polaris\_runs} ($\approx70\,\mathrm{GB}$); datasets in
\texttt{/scratch/sjafaritabar/starforge\_polaris\_dataset\_v1} and
\texttt{/scratch/sjafaritabar/starforge\_polaris\_dataset\_v1\_stars\_eval}; audit reports in
\texttt{/scratch/sjafaritabar/audit}.
. Required packages in the preamble:
%   \usepackage{amsmath}  \usepackage{booktabs}  \usepackage{graphicx}  \usepackage{longtable}
% References are in dataset_construction.bib (\bibliography{dataset_construction}).
% Every number below is taken from the pipeline code, its logs, manifests and the audit report
% (paths in Sect. "Data products and code"); dates are UTC.
% =====================================================================================================

\section{A synthetic dust-emission and polarization dataset of star-forming cores}
\label{sec:dataset}

This section describes how the training data for this thesis were constructed. Starting from a
magnetohydrodynamic simulation of a star-forming molecular cloud, dense cores were identified in every
simulation snapshot, a $2.2\,\mathrm{pc}$ cube around each selected core was converted into a radiative-transfer
model, and the dust temperature, grain alignment and polarized dust emission of every cube were computed with
the Monte Carlo code POLARIS. The resulting images were combined with the three-dimensional state of the gas
into a table with one row per line of sight. Table~\ref{tab:pipeline} summarises the five steps.

\begin{table}[ht]
  \centering
  \caption{The five steps of the pipeline. Each step is a separate script; versions refer to the
  version strings recorded in each step's output.}
  \label{tab:pipeline}
  \small
  \begin{tabular}{@{}llll@{}}
    \toprule
    Step & Task & Product & Code (version) \\
    \midrule
    1 & Dense cores in every snapshot   & 11\,599 cores in 51 snapshots & \texttt{find\_starforge\_cores\_dendrogram\_v5.py} (5.2.0) \\
    2 & Non-overlapping boxes, families & 1908 boxes, 317 families     & \texttt{select\_boxes.py} \\
    3 & POLARIS grids and commands      & 1908 grids of $64^3$ cells    & \texttt{prepare\_grids.py} (3.5.1) \\
    4 & Radiative transfer              & 1908 runs, 217\,512 images    & \texttt{run\_polaris.py}, POLARIS \\
    5 & Sightline table                 & 23\,445\,504 rows, 1443 columns & \texttt{build\_dataset.py} (1.1.1 / 1.2.0) \\
    \bottomrule
  \end{tabular}
\end{table}

% -----------------------------------------------------------------------------------------------------
\subsection{The STARFORGE simulation}
\label{sec:dataset:sim}

The gas data come from the STARFORGE \citep{Grudic2021} simulation \texttt{M2e5\_R30\_Res585}, run with the
meshless magnetohydrodynamics code GIZMO \citep{Hopkins2015} (header \texttt{GIZMO\_version} 2022). The initial
condition is a uniform-density spherical cloud of $2\times10^{5}\,M_\odot$ and radius $30\,\mathrm{pc}$ (measured:
$1.993\times10^{5}\,M_\odot$ inside $R=30\,\mathrm{pc}$, median density $1.19\times10^{-22}\,\mathrm{g\,cm^{-3}}$)
embedded in a periodic box of side $300\,\mathrm{code\ units}=300.42\,\mathrm{pc}$. The initial three-dimensional
velocity dispersion of the cloud is $5.86\,\mathrm{km\,s^{-1}}$, which corresponds to a virial parameter
$\alpha=2$; the initial gas temperature is $\approx12\,\mathrm{K}$; gas cells have a median mass of
$1.001\times10^{-3}\,M_\odot$ (range $1.001$--$1.301\times10^{-3}\,M_\odot$). The interstellar radiation field
strength parameter in the snapshot header is $1$; solar abundances are $Z_\odot=0.0142$, $Y_\odot=0.2703$.

\paragraph{Units.} The code units are $1.0014\,\mathrm{pc}$ (length), $1.0008\,M_\odot$ (mass) and
$1\,\mathrm{m\,s^{-1}}$ (velocity), so one code time unit is $979.16\,\mathrm{Myr}$. The units were verified
against the initial conditions (\texttt{check\_snapshot\_units.py}): with these units the gravitational constant
derived from the header units agrees with the header value to $3.4\times10^{-4}$, the enclosed mass and density
match the nominal cloud, and the velocity dispersion matches $\alpha=2$. The saved \texttt{MagneticField} is in
tesla, \emph{not} in the unit given by the header field \texttt{Internal\_UnitB\_In\_Gauss}
($2.91\times10^{-9}\,\mathrm{G}$): read in tesla, the initial cloud field is $7.19\,\mu\mathrm{G}$, consistent
(within a factor $1.26$, the uncertainty of the critical mass-to-flux normalisation) with the nominal
normalised mass-to-flux ratio $4.2$, whereas the gauss or header-unit readings are off by four or more
orders of magnitude. POLARIS expects the field in tesla, so the field is passed on unconverted.

\paragraph{Snapshots.} Snapshots \texttt{snapshot\_000} to \texttt{snapshot\_500} in steps of 10 were analysed
(51 snapshots, each $\approx 31\,\mathrm{GB}$ with $\approx 2.5\times10^{8}$ gas cells). Dense cores exist from
snapshot 130 onwards. Snapshots 130--420 are $0.0406\,\mathrm{Myr}$ apart (0.528--1.707\,Myr); snapshots
420--500 are $0.406\,\mathrm{Myr}$ apart (1.707--4.957\,Myr). Sink particles (stars) exist only in the late
snapshots.

% -----------------------------------------------------------------------------------------------------
\subsection{Step 1: dense cores}
\label{sec:dataset:cores}

\paragraph{Molecular density.} For every gas cell the molecular hydrogen number density is
\begin{equation}
  n_{\mathrm{H_2}} = \frac{\rho\, X_\mathrm{H}\, f_\mathrm{neutral}\, f_\mathrm{mol}}{2\,m_\mathrm{H}},\qquad
  X_\mathrm{H} = 1 - Z - Y,
  \label{eq:nh2}
\end{equation}
where $\rho$ is the cell density, $Z$ and $Y$ are the first two columns of \texttt{Metallicity},
$f_\mathrm{neutral}$ is \texttt{NeutralHydrogenAbundance}, $f_\mathrm{mol}$ is \texttt{MolecularMassFraction},
and $m_\mathrm{H}=1.6735575\times10^{-24}\,\mathrm{g}$. Invalid or out-of-range inputs abort the run instead of
being clipped.

\paragraph{Dendrogram.} Cells with $n_{\mathrm{H_2}} > 10^{4}\,\mathrm{cm^{-3}}$ were selected. Two selected cells
are neighbours if either lies within the kernel support radius $h$ (\texttt{SmoothingLength}) of the other,
using periodic distances; the resulting graph was symmetrised. A dendrogram \citep{Rosolowsky2008} was computed
with \texttt{astrodendro} 0.3.1 on this particle graph with
$\texttt{min\_value}=10^{4}\,\mathrm{cm^{-3}}$, $\texttt{min\_delta}=10^{4}\,\mathrm{cm^{-3}}$ and
$\texttt{min\_npix}=100$ cells. Every leaf of the dendrogram, i.e.\ a local maximum of $n_{\mathrm{H_2}}$
with at least 100 cells and a contrast of at least $10^{4}\,\mathrm{cm^{-3}}$ above the level at which it
merges with a neighbour, is a \emph{core}. A leaf is a density peak; it is not required to be gravitationally
bound or associated with a star.

\paragraph{Core properties.} For each core the catalogue records, among others, the peak cell (position,
$n_{\mathrm{H_2}}$, particle ID), mass, volume and radii, shape, internal velocity dispersion, mean gas and dust
temperature, mass-weighted field, the gas-only virial ratio $2(K+U)/|W|$ (direct pair summation with
$10^{-3}\,\mathrm{pc}$ Plummer softening, up to $2\times10^{4}$ cells), and the sink particles near the core,
including the number of sinks within a $2.2\,\mathrm{pc}$ cube centred on the peak.
Step 1 found 11\,599 cores in total (none before snapshot 130; 1038 in snapshot 500) with a total mass of
$4535\,M_\odot$.

% -----------------------------------------------------------------------------------------------------
\subsection{Step 2: boxes and families}
\label{sec:dataset:boxes}

\paragraph{Box selection.} A \emph{box} is an axis-aligned cube of side $S=2.2\,\mathrm{pc}$ centred on a core's
density peak. Within each snapshot, boxes were selected by greedy non-maximum suppression: cores were sorted by
mass (ties, which are common because core masses are multiples of the cell mass, broken by higher peak density,
then lower \texttt{core\_id}), and a core was kept unless its cube overlapped an already kept cube. Two cubes
overlap if the periodic Chebyshev distance between their centres is $<S$. This gives 1908 boxes in 38 snapshots
(130--500; 1 to 104 boxes per snapshot, median 55). The central cores have masses of $0.10$--$9.09\,M_\odot$
(median $0.41\,M_\odot$) and peak densities $n_{\mathrm{H_2}}=2.0\times10^{4}$--$7.0\times10^{9}\,\mathrm{cm^{-3}}$
(median $3.9\times10^{4}\,\mathrm{cm^{-3}}$). Of all cores, 3604 ($1249\,M_\odot$) lie in no box: they were
suppressed by a more massive neighbour without lying inside its cube.

\paragraph{Families.} Boxes of consecutive snapshots largely image the same gas. Two boxes in consecutive
snapshots (consecutive among snapshots that contain boxes) were therefore linked if their cubes overlap
(periodic Chebyshev distance $<S$), and the connected components of these links, found with a union--find
structure, were labelled \texttt{family\_id}. The 1908 boxes form 317 families; the median family has 2 boxes,
the 90th percentile 9.4, the 99th percentile 107, and the largest family has 260 boxes. Six families have more
than 50 boxes. Families are the unit of the train/validation/test split (Sect.~\ref{sec:dataset:splits}).

% -----------------------------------------------------------------------------------------------------
\subsection{Step 3: radiative-transfer models}
\label{sec:dataset:grids}

\paragraph{Grid geometry.} Each box was converted into a POLARIS octree grid that is a uniform cube of
$64^{3}$ cells (octree level 6) of size $\Delta=2.2/64=0.034375\,\mathrm{pc}$. The grid is shifted by half a
cell so that the core's peak cell lies at the centre of cell $(32,32,32)$; centring the cube exactly on the peak
would place the peak on a corner shared by eight cells. Each snapshot was read once and the gas of all its boxes
was cut out in that single pass; boxes were then built in parallel.

\paragraph{Deposition.} Every gas cell whose kernel overlaps the cube (centre within $S/2+h$ of the box centre
along each axis) was deposited with GIZMO's compact cubic-spline kernel,
\begin{equation}
  W(q) \propto \begin{cases} 1-6q^{2}+6q^{3}, & 0\le q<\tfrac12,\\ 2(1-q)^{3}, & \tfrac12\le q<1,\\ 0, & q\ge 1,\end{cases}
  \qquad q = \frac{|\mathbf{x}_c-\mathbf{x}_i|}{h_i},
\end{equation}
sampled at the centres $\mathbf{x}_c$ of all grid cells within the kernel, \emph{including cells outside the cube}.
The weights $w_{ic}=W(q_{ic})/\sum_{c'}W(q_{ic'})$ are normalised over the full kernel and only the in-cube part
is kept, so each grid holds exactly the mass inside the cube and the faces are not under-dense. If a kernel
contains no cell centre ($h_i$ well below $\Delta$), the cell's mass goes to the grid cell containing its centre
(this applied to 1.07\% of the $2.73\times10^{8}$ deposited gas cells). Cell quantities are
\begin{equation}
  \rho_c = \frac{\sum_i m_i w_{ic}}{\Delta^{3}},\qquad
  \langle X\rangle_c = \frac{\sum_i m_i w_{ic} X_i}{\sum_i m_i w_{ic}},
\end{equation}
with $X$ the dust temperature, gas temperature and the three field components (mass-weighted means). A cell
receiving no mass would get a density floor of $10^{-6}$ times the smallest occupied density and the fields of
the nearest occupied cell; this happened for a single cell in a single box ($3.6\times10^{-14}\,M_\odot$).
The grid mass agrees with the mass of the gas cells centred inside the cube to between $-0.12\%$ and $+0.08\%$.
Boxes contain $3.5\times10^{4}$--$7.9\times10^{5}$ contributing gas cells (median $1.3\times10^{5}$) and
$27.6$--$748.6\,M_\odot$ of gas (median $113\,M_\odot$).

\paragraph{Grid fields.} The grid stores (POLARIS field IDs in brackets) the gas number density
$n=\rho/(\mu m_u)$ in $\mathrm{m^{-3}}$ [0], the STARFORGE dust temperature [2], the gas temperature [3] and
$B_x, B_y, B_z$ in tesla [4--6]. The POLARIS branch used reads the gas density only as a number density and
converts it back with $\rho=n\,\mu\,m_u$, where $m_u=1.66053904\times10^{-27}\,\mathrm{kg}$ (POLARIS's
\texttt{m\_H}); the same $n$ and $\mu$ enter the radiative-torque calculation (gas damping, thermal speed).
$\mu$ is therefore the mean mass per gas particle. We use $\mu=2.33$: the mass-weighted median of the particle
mean molecular weight of box gas, computed from the snapshot chemistry, is 2.32--2.35 for
$n_\mathrm{H}=10^{3}$--$10^{6}\,\mathrm{cm^{-3}}$ (diagnostic \texttt{check\_mu.py}, snapshot 500), and
$\rho$, hence the dust mass, is reproduced exactly for any $\mu$.

\paragraph{Dust.} One dust component was used: astrodust \citep{Hensley2023} without iron inclusions
($f_\mathrm{Fe}=0$), with the catalogue's magnetic susceptibility $\chi=4.47\times10^{-4}$ and material density
$2898\,\mathrm{kg\,m^{-3}}$, a power-law size distribution $\mathrm{d}n/\mathrm{d}a\propto a^{-3.5}$
\citep{Mathis1977} between $a_\mathrm{min}=5\,\mathrm{nm}$ and $a_\mathrm{max}=250\,\mathrm{nm}$, and a
dust-to-gas mass ratio of $0.01$ (POLARIS \texttt{<mass\_fraction>}). The ratio is the one STARFORGE itself
uses: its field \texttt{DustToGasRatio\_Local} stores $Z$ times the local dust-to-metals ratio relative to solar,
with a median of $0.01416=Z_\odot$ over all boxes, which corresponds to an absolute dust-to-gas ratio
$0.01\,Z/Z_\odot = 0.01$. The catalogue covers $1\,\mathrm{nm}$--$1\,\mu\mathrm{m}$ in size and
$0.05\,\mu\mathrm{m}$--$5\,\mathrm{mm}$ in wavelength, i.e.\ POLARIS's internal wavelength grid
(100 points, $0.0912$--$2000\,\mu\mathrm{m}$) and all detector wavelengths. With this susceptibility the
Larmor-precession limit of POLARIS's alignment model is never reached: the smallest Larmor-limited grain size
over all $5.0\times10^{8}$ grid cells is $408\,\mathrm{nm}>a_\mathrm{max}$ (\texttt{check\_larmor\_limit.py}).

\paragraph{Radiation sources.} Every box is illuminated by the interstellar radiation field (ISRF) of
\citet{Mathis1983} with strength $G_0=1$, launched from a sphere of 10 times the circumscribed radius of the grid
($10\times\tfrac{\sqrt3}{2}\times2.2\,\mathrm{pc}=19.1\,\mathrm{pc}$). In addition, every sink particle inside the
cube with positive radius and luminosity was added as a blackbody point source at its position, with the
STARFORGE protostellar radius $R_\ast$ and an effective temperature from
$L_\ast = 4\pi R_\ast^{2}\sigma_\mathrm{SB}T_\mathrm{eff}^{4}$. Fifty-two boxes contain stars, 275 sources in total
(up to 39 in one box): $L_\ast=6.8\times10^{-3}$--$1.0\times10^{4}\,L_\odot$ (median $0.16\,L_\odot$),
$R_\ast=0.31$--$20.8\,R_\odot$, $T_\mathrm{eff}=2990$--$27\,900\,\mathrm{K}$, sink masses
$0.02$--$11.7\,M_\odot$. (Step 1 counted 273 sinks in these boxes; in one box, 1851, the half-cell shift of the
grid moves two further sinks inside the cube.) Boxes without stars are heated by the ISRF only.

\paragraph{Photon numbers and Monte Carlo noise.} With radiative-torque alignment, POLARIS launches the
requested number of photons for each of its 100 wavelengths. Production used $10^{7}$ ISRF photons and $10^{6}$
photons per star (per wavelength). These values were chosen from two convergence studies, each running every
configuration twice with independent random numbers and taking the per-cell noise as
$\mathrm{median}\,|x_1-x_2|/\sqrt2$:
(i) ISRF photons $10^{6}$--$10^{7}$ and launch radii 3--30 on box \texttt{snapshot\_440/core\_000395}; at $10^{7}$
photons and radius 10 the dust-temperature noise is $0.0078\,\mathrm{K}$ in the interior, $0.0087\,\mathrm{K}$ on
faces, $0.0103\,\mathrm{K}$ on edges and $0.0137\,\mathrm{K}$ in corners ($0.0081\,\mathrm{K}$ overall), the
alignment radius varies by $0.63\%$, and at $850\,\mu\mathrm{m}$ the noise is $0.026\%$ in $I$, $0.20\%$ in $p$
and $0.035^\circ$ in the polarization angle. Smaller radii concentrate the noise in the corners
($0.073\,\mathrm{K}$ at radius 3), larger radii raise it in the interior ($0.0105\,\mathrm{K}$ at radius 30).
(ii) Star photons $10^{5}$ vs.\ $10^{6}$ on box \texttt{snapshot\_500/core\_000873}: $10^{6}$ halves the
temperature noise in gas below $25\,\mathrm{K}$ (from 0.077--0.084 to 0.039--0.045\,K) and changes the
$53\,\mu\mathrm{m}$ intensity by a median of $2.0\%$ ($0.22\%$ at $850\,\mu\mathrm{m}$).

\paragraph{POLARIS tasks.} Each box has a command file with two tasks, both with radiative-torque alignment
(\texttt{ALIG\_RAT}) and 8 threads:
\begin{enumerate}
  \item \texttt{CMD\_TEMP\_RAT}: Monte Carlo radiative transfer of the ISRF and stellar photons; computes the
  equilibrium dust temperature of every cell and the radiative-torque alignment radius $a_\mathrm{alig}$ (grains
  larger than $a_\mathrm{alig}$ are aligned with the magnetic field). The STARFORGE dust temperature in the input
  grid is replaced by this temperature. Writes \texttt{temp/grid\_temp\_rat.dat}.
  \item \texttt{CMD\_DUST\_EMISSION} on that grid: ray-traced thermal emission of the aligned grains, giving the
  Stokes $I,Q,U,V$, the optical depth and the column density for 38 wavelengths and 3 viewing directions
  (114 detectors) with $64\times64$ pixels each (\texttt{max\_subpixel\_lvl} 0). The detector covers the grid, so
  one pixel corresponds to one grid column ($0.034375\,\mathrm{pc}$); the source distance is
  $140\,\mathrm{pc}$, so a pixel subtends $50.6''$ and
  $\Omega_\mathrm{pix}=(\Delta/d)^{2}=6.03\times10^{-8}\,\mathrm{sr}$.
\end{enumerate}
The 38 wavelengths (Table~\ref{tab:bands}) are those of common far-infrared to millimetre instruments and are
read from one table (\texttt{wavelength\_telescope\_mapping.csv}) by both the grid preparation and the dataset
builder. The three viewing directions are the POLARIS detector rotations
$(\vartheta_1,\vartheta_2)=(0^\circ,0^\circ)$, $(-90^\circ,0^\circ)$ and $(0^\circ,90^\circ)$, which look along the
grid $z$, $y$ and $x$ axes (Sect.~\ref{sec:dataset:geometry}). POLARIS was the \texttt{experimental} branch,
commit \texttt{6c5e1a89c6e14ce806c47bbb87947b98e3f761b0} \citep{Reissl2016}.

% -----------------------------------------------------------------------------------------------------
\subsection{Step 4: radiative-transfer runs}
\label{sec:dataset:runs}

All 1908 boxes were run between 2026-10-31 22:47 and 2026-11-09 12:41 on one machine (AMD EPYC 9554P, 64 cores
/ 128 hardware threads, $377\,\mathrm{GB}$ memory), each box in its own directory, at the lowest scheduling
priority (\texttt{nice 19}), with up to 16 boxes at a time (8 threads each). The runner measured the CPU use of
other users every $15\,\mathrm{s}$ and kept 8 hardware threads free above their peak use over the previous
10 minutes; when the machine was needed, the most recently started boxes were suspended (\texttt{SIGSTOP}) and
resumed 10 minutes after the demand had gone, so no run was ever killed or restarted. Boxes containing stars
were run first. A box took $0.92$--$5.98\,\mathrm{h}$ of wall-clock time (median $1.46\,\mathrm{h}$ at up to 14
concurrent boxes), $2810$ box-hours in total, of which boxes spent $5.3\,\mathrm{h}$ suspended.

Every run was validated before being accepted: the temperature grid exists and all dust temperatures are finite
and within $2$--$3000\,\mathrm{K}$; all 114 detector files exist; $I$, $Q$ and $U$ are finite everywhere; and at
wavelengths $\geq50\,\mu\mathrm{m}$, $I>0$ in every pixel and $Q,U$ are not identically zero. All 1908 runs
finished with exit code 0 and passed; none failed. POLARIS's Monte Carlo is seeded by thread, and the assignment
of photons to threads varies between runs, so runs are statistically but not bitwise reproducible.

% -----------------------------------------------------------------------------------------------------
\subsection{Step 5: the sightline dataset}
\label{sec:dataset:table}

\subsubsection{Geometry and orientation}
\label{sec:dataset:geometry}

Grid arrays are indexed $[i_x,i_y,i_z]$ and images $[\mathrm{row},\mathrm{col}]$, both from 0 to 63. The mapping
between them was determined empirically (\texttt{check\_view\_geometry.py}): for each view, the POLARIS
column-density image was compared with $\sum n\,\Delta$ along each grid axis in all 48 combinations of
line-of-sight axis, transposition and flips; the best orientation matches to a relative difference of
$\approx10^{-13}$. Which end of the line of sight faces the observer cannot be seen in optically thin emission; it
was determined where the dust is optically thick ($21\,\mu\mathrm{m}$, pixels with $\tau>1$), where
$I=\sum_k B_\nu(T_k)\,\Delta\tau_k\,e^{-\tau_k}$ depends on the side from which the absorption is counted
(\texttt{check\_observer\_side.py}). The model built from POLARIS's own temperatures reproduces the image with a
near-constant ratio only if the high-index end of the line-of-sight axis faces the observer, which holds in every
view. Table~\ref{tab:views} gives the result.

\begin{table}[ht]
  \centering
  \caption{Mapping between a dataset row (view $v$, pixel row $r$, pixel column $c$, profile index $k$, with $k=0$
  the cell nearest the observer) and the grid, and the field components in the image frame.}
  \label{tab:views}
  \small
  \begin{tabular}{@{}cccccc@{}}
    \toprule
    View & POLARIS rotation & Line of sight & Image $+$col, $+$row & Grid cell $(i_x,i_y,i_z)$ & $(B_\mathrm{col},B_\mathrm{row},B_\mathrm{los})$ \\
    \midrule
    0 & $(0^\circ,0^\circ)$   & $z$ & $+x$, $+y$ & $(c,\ r,\ 63-k)$    & $(B_x,\ B_y,\ -B_z)$ \\
    1 & $(-90^\circ,0^\circ)$ & $y$ & $+x$, $-z$ & $(c,\ 63-k,\ 63-r)$ & $(B_x,\ -B_z,\ -B_y)$ \\
    2 & $(0^\circ,90^\circ)$  & $x$ & $-z$, $+y$ & $(63-k,\ r,\ 63-c)$ & $(-B_z,\ B_y,\ -B_x)$ \\
    \bottomrule
  \end{tabular}
\end{table}

Angles in the image plane are measured from the image $+$col axis towards $+$row. With Stokes $U$ as written by
POLARIS, the plane-of-sky field direction is the polarization angle plus $90^\circ$: comparing
$\tfrac12\arctan(U,Q)+90^\circ$ with a density-weighted Stokes model of the grid field in this frame (in the more
polarized half of the pixels) gives median differences of $0.2^\circ$--$0.8^\circ$, against $22^\circ$--$54^\circ$ if the
sign of $U$ is reversed. $B_\mathrm{los}>0$ points away from the
observer. Every box and view was checked during the build: $\sum_k n_k\Delta$ along the rotated profiles
reproduces POLARIS's column-density image to a maximum relative difference of $5.4\times10^{-13}$
(the build aborts above $10^{-5}$).

\subsubsection{Contents}

One row is one line of sight: one pixel of one view of one box. The dataset has
$1908\times3\times64\times64=23\,445\,504$ rows and 1443 columns (Table~\ref{tab:columns}).

\begin{table}[ht]
  \centering
  \caption{Columns of the dataset. $k=00\ldots63$ runs along the line of sight from the observer; columns
  marked (+log) also exist as \texttt{<name>\_log} $=\log_{10}$ (NaN where the value is $\le0$). Values are
  \texttt{float32} except $I,Q,U$ (\texttt{float64}) and the integer metadata.}
  \label{tab:columns}
  \small
  \begin{tabular}{@{}p{0.17\textwidth}p{0.47\textwidth}p{0.12\textwidth}r@{}}
    \toprule
    Group & Columns & Unit & Number \\
    \midrule
    metadata & \texttt{snapshot, time\_Myr, core\_id, box\_id, family\_id, n\_stars, split, view, pix\_row, pix\_col} & & 10 \\
    sightline & \texttt{N\_col, Sigma, Sigma\_dust} (+log) & $\mathrm{cm^{-2}}$, $\mathrm{g\,cm^{-2}}$ & 6 \\
              & \texttt{temp\_polaris\_mean, temp\_starforge\_mean} (+log) & K & 4 \\
              & \texttt{B\_tot\_mean} (+log) & $\mu$G & 2 \\
              & \texttt{cos2psi\_B\_mean, sin2psi\_B\_mean, B\_coherence} & -- & 3 \\
    bands ($\times38$) & \texttt{<band>\_I, \_Q, \_U} & MJy\,sr$^{-1}$ & 114 \\
              & \texttt{<band>\_p, \_polangle, \_I\_log, \_p\_log} & --, deg & 152 \\
    profiles ($\times64$) & \texttt{dens\_k, dens\_dust\_k} (+log) & $\mathrm{cm^{-3}}$, $\mathrm{g\,cm^{-3}}$ & 256 \\
              & \texttt{temp\_polaris\_k, temp\_starforge\_k} (+log) & K & 256 \\
              & \texttt{B\_tot\_k, B\_pos\_k} (+log) & $\mu$G & 256 \\
              & \texttt{cos2psi\_B\_k, sin2psi\_B\_k, sin2gamma\_B\_k} & -- & 192 \\
              & \texttt{B\_col\_k, B\_row\_k, B\_los\_k} & $\mu$G & 192 \\
    \midrule
    total & & & 1443 \\
    \bottomrule
  \end{tabular}
\end{table}

\paragraph{Profiles.} For each cell $k$ along the line of sight:
$\texttt{dens}=n$ in $\mathrm{cm^{-3}}$ (gas particles, $\mu=2.33$);
$\texttt{dens\_dust}=0.01\,n\,\mu\,m_u$ in $\mathrm{g\,cm^{-3}}$;
$\texttt{temp\_polaris}$, the dust temperature computed by POLARIS; $\texttt{temp\_starforge}$, the deposited
STARFORGE dust temperature; the field in $\mu$G ($1\,\mathrm{T}=10^{10}\,\mu\mathrm{G}$) in the image frame
(Table~\ref{tab:views}) and the derived quantities
\begin{gather}
  B_\mathrm{pos}^{2}=B_\mathrm{col}^{2}+B_\mathrm{row}^{2},\qquad B_\mathrm{tot}^{2}=B_\mathrm{pos}^{2}+B_\mathrm{los}^{2},\\
  \cos2\psi_B=\frac{B_\mathrm{col}^{2}-B_\mathrm{row}^{2}}{B_\mathrm{pos}^{2}},\qquad
  \sin2\psi_B=\frac{2B_\mathrm{col}B_\mathrm{row}}{B_\mathrm{pos}^{2}},\qquad
  \sin^{2}\gamma_B=\frac{B_\mathrm{pos}^{2}}{B_\mathrm{tot}^{2}},
\end{gather}
where $\psi_B$ is the plane-of-sky field angle and $\gamma_B$ the angle between the field and the line of sight.
The double-angle form removes the $180^\circ$ ambiguity of the field direction, which dust polarization cannot
measure.

\paragraph{Sightline quantities.} $N_\mathrm{col}$ is POLARIS's column-density image converted from
$\mathrm{m^{-2}}$ to $\mathrm{cm^{-2}}$; $\Sigma=N_\mathrm{col}\,\mu m_u$ and $\Sigma_\mathrm{dust}=0.01\,\Sigma$.
Means are density-weighted, $\langle X\rangle=\sum_k n_kX_k/\sum_k n_k$, for both dust temperatures and
$B_\mathrm{tot}$. The field orientation that a polarization map measures is described by
\begin{equation}
  z = \sum_k n_k\,\sin^{2}\gamma_{B,k}\;e^{2i\psi_{B,k}},\qquad
  (\texttt{cos2psi\_B\_mean},\texttt{sin2psi\_B\_mean})=\frac{(\operatorname{Re}z,\operatorname{Im}z)}{|z|},\qquad
  \texttt{B\_coherence}=\frac{|z|}{\sum_k n_k}\in[0,1],
  \label{eq:coherence}
\end{equation}
so \texttt{B\_coherence} combines field tangling along the line of sight and inclination; it correlates with the
$850\,\mu\mathrm{m}$ polarization fraction with $r=0.87$--$0.94$ (Pearson, per file, in the eight files checked
by \texttt{validate\_dataset.py}).

\paragraph{Bands.} POLARIS returns $I,Q,U$ in Jy per pixel; they are converted to $\mathrm{MJy\,sr^{-1}}$ by
dividing by $10^{6}\,\Omega_\mathrm{pix}$. The polarization fraction is $p=\sqrt{Q^{2}+U^{2}}/I$ (NaN where
$I\le0$) and \texttt{polangle} $=\tfrac12\arctan(U,Q)$ in degrees, mapped to $(-90^\circ,90^\circ]$; it is the
E-vector angle, and the inferred field angle is $\texttt{polangle}+90^\circ$. $I,Q,U$ are stored in double
precision because at $3.4$--$5.8\,\mu\mathrm{m}$ the intensities reach $10^{-108}\,\mathrm{MJy\,sr^{-1}}$, below
the single-precision range. POLARIS was run without stochastic heating of small grains, so the 14 bands at
$\le25\,\mu\mathrm{m}$ are flagged \texttt{reliable: false} in the manifest; the 24 bands at
$53\,\mu\mathrm{m}$--$1.38\,\mathrm{mm}$ are reliable. The images contain no observational noise and no beam.

\begin{table}[ht]
  \centering
  \caption{The 38 detector wavelengths and band names (column prefixes). Bands at $\le25\,\mu\mathrm{m}$ are
  flagged unreliable.}
  \label{tab:bands}
  \footnotesize
  \begin{tabular}{@{}rl rl@{}}
    \toprule
    $\lambda$ [$\mu$m] & Band & $\lambda$ [$\mu$m] & Band \\
    \midrule
    3.4   & \texttt{wise\_1}            & 89    & \texttt{sofia\_hawc\_plus\_c} \\
    3.6   & \texttt{spitzer\_irac\_1}   & 90    & \texttt{akari\_fis\_90} \\
    4.5   & \texttt{spitzer\_irac\_2}   & 100   & \texttt{herschel\_pacs\_2\_or\_iras\_100} \\
    4.6   & \texttt{wise\_2}            & 140   & \texttt{akari\_fis\_140} \\
    5.8   & \texttt{spitzer\_irac\_3}   & 154   & \texttt{sofia\_hawc\_plus\_d} \\
    8.0   & \texttt{spitzer\_irac\_4}   & 160   & \texttt{herschel\_pacs\_3} \\
    8.28  & \texttt{msx\_1}             & 214   & \texttt{sofia\_hawc\_plus\_e} \\
    12.0  & \texttt{wise\_3\_or\_iras\_12} & 250 & \texttt{herschel\_spire\_1} \\
    12.13 & \texttt{msx\_2}             & 350   & \texttt{herschel\_spire\_2\_or\_alma\_band\_10\_c43\_8\_or\_planck\_hfi\_857ghz} \\
    14.65 & \texttt{msx\_3}             & 460   & \texttt{alma\_band\_9\_c43\_8} \\
    21.34 & \texttt{msx\_4}             & 500   & \texttt{herschel\_spire\_3} \\
    22.0  & \texttt{wise\_4}            & 550   & \texttt{planck\_hfi\_545ghz} \\
    24.0  & \texttt{spitzer\_mips\_24}  & 690   & \texttt{alma\_band\_8\_c43\_8} \\
    25.0  & \texttt{iras\_25}           & 849   & \texttt{planck\_hfi\_353ghz} \\
    53    & \texttt{sofia\_hawc\_plus\_a} & 870 & \texttt{apex\_laboca} \\
    60    & \texttt{iras\_60}           & 945   & \texttt{alma\_band\_7\_c43\_8} \\
    63    & \texttt{sofia\_hawc\_plus\_b} & 1100 & \texttt{cso\_bolocam} \\
    65    & \texttt{akari\_fis\_65}     & 1300  & \texttt{alma\_band\_6\_c43\_8} \\
    70    & \texttt{herschel\_pacs\_1}  & 1380  & \texttt{planck\_hfi\_217ghz} \\
    \bottomrule
  \end{tabular}
\end{table}

\subsubsection{Train, validation and test splits}
\label{sec:dataset:splits}

Boxes in the same family image largely the same gas a short time apart, so splitting by box or by pixel would put
near-identical data into training and test sets. Whole families were therefore assigned to one split, and all
three views of a box share its split. Families were shuffled with a fixed seed (NumPy \texttt{default\_rng},
seed 20261111) and assigned in turn to the split furthest below its target share (80/10/10\% of the boxes) among
the splits the family fits into; families with more than 50 boxes always go to the training set, so that no single
region dominates a validation or test set. Two versions of the dataset were built from identical data
(Table~\ref{tab:splits}):

\begin{itemize}
  \item \texttt{starforge\_polaris\_dataset\_v1} (builder 1.1.1): the split described above. All 52 boxes with
  stars belong to families assigned to the training set.
  \item \texttt{starforge\_polaris\_dataset\_v1\_stars\_eval} (builder 1.2.0, option \texttt{--assign}): as v1, but
  the four small star-forming families 41 and 226 are moved to validation and 133 and 139 to test; every other box
  keeps its split. Star-forming families 140 and 280 stay in training because they lie only $3.3$ and
  $2.8\,\mathrm{pc}$ from the large training family 2. The non-star validation and test families are identical in
  both versions, so models trained on either can be compared on that common subset.
\end{itemize}

\begin{table}[ht]
  \centering
  \caption{Splits of the two dataset versions (boxes; each box gives $3\times4096$ rows).}
  \label{tab:splits}
  \small
  \begin{tabular}{@{}lrrrrrr@{}}
    \toprule
     & \multicolumn{3}{c}{\texttt{v1}} & \multicolumn{3}{c}{\texttt{v1\_stars\_eval}} \\
    \cmidrule(lr){2-4}\cmidrule(l){5-7}
    Split & boxes & families & boxes with stars & boxes & families & boxes with stars \\
    \midrule
    train      & 1538 (80.6\%) & 207 & 52 & 1458 (76.4\%) & 203 & 36 \\
    validation &  187 (9.8\%)  &  58 &  0 &  216 (11.3\%) &  60 &  8 \\
    test       &  183 (9.6\%)  &  52 &  0 &  234 (12.3\%) &  54 &  8 \\
    \bottomrule
  \end{tabular}
\end{table}

In v1 the largest family makes up 17\% of the training set, 16\% of the validation set and 21\% of the test
set; all three splits span 0.5--5.0\,Myr. All star-forming families lie within the same central region of the
cloud and appear at late times (snapshots 430--500), with their nearest boxes 2.6--4.8\,pc apart; scores on the
star-forming test families therefore measure performance on new cores of that region rather than of a different
cloud.

% -----------------------------------------------------------------------------------------------------
\subsection{The dataset}
\label{sec:dataset:contents}

\subsubsection{What one row contains}

Each row describes one line of sight through one box, seen in one of the three viewing directions. It holds
\begin{enumerate}
  \item \emph{identifiers} that locate the sightline in the simulation (snapshot, time, core, box, family), in the
  split, and in its image (view, pixel row and column);
  \item the \emph{observables} a telescope would measure along this sightline: Stokes $I$, $Q$, $U$, the
  polarization fraction and the polarization angle in 38 bands;
  \item \emph{sightline targets}: quantities integrated or averaged along the line of sight (column and surface
  density, mean dust temperatures, mean field strength, field orientation and coherence);
  \item \emph{profile targets}: the three-dimensional state of the 64 grid cells along the line of sight, ordered
  from the observer: gas and dust density, the two dust temperatures and the magnetic field (strength, orientation
  and components).
\end{enumerate}
Observables and targets of the same row refer to exactly the same column of grid cells, and the profiles are
already in the observer's frame (Table~\ref{tab:views}), so no geometric transformation is needed when using the
data. Because each row group holds one box and view, the rows of a row group reassemble into a $64\times64$ image
with \texttt{pix\_row} and \texttt{pix\_col}, and the profile columns of that image into a $64\times64\times64$
cube; the dataset therefore serves sightline-based (per-pixel) as well as image-based models.

\subsubsection{Files}

The dataset is stored in Apache Parquet format (written by \texttt{pyarrow}, zstd compression); Table~\ref{tab:files}
lists the files. Each split directory holds one file per snapshot that has boxes in that split, and each file holds
one row group of 4096 rows per box and view, so a single box and view can be read without reading the rest of the
file. The data occupy $178\,\mathrm{GB}$ ($166\,\mathrm{GiB}$) per dataset version. Building them from the
POLARIS output takes 27 minutes with 16 processes.

\begin{table}[ht]
  \centering
  \caption{Files of a dataset version (sizes and counts for \texttt{v1}).}
  \label{tab:files}
  \small
  \begin{tabular}{@{}p{0.27\textwidth}p{0.52\textwidth}r@{}}
    \toprule
    Path & Content & Size \\
    \midrule
    \texttt{data/train/<snapshot>.parquet} & 37 files; 1538 boxes, 4614 row groups, 18\,898\,944 rows & 143.5\,GB \\
    \texttt{data/val/<snapshot>.parquet}   & 38 files; 187 boxes, 561 row groups, 2\,297\,856 rows & 17.4\,GB \\
    \texttt{data/test/<snapshot>.parquet}  & 33 files; 183 boxes, 549 row groups, 2\,248\,704 rows & 17.1\,GB \\
    \texttt{splits.csv} & one line per box: \texttt{box} (\texttt{snapshot\_NNN/core\_NNNNNN}), \texttt{family\_id}, \texttt{split} & \\
    \texttt{manifest.json} & builder version and path, creation time, input paths, numbers of boxes, rows and columns,
      geometry conventions, units, the band table with wavelength and \texttt{reliable} flag, split seed, target fractions,
      sizes (and, in \texttt{v1\_stars\_eval}, the family overrides), the largest column-density mismatch, and a QA record per box
      (column-density check per view, non-positive intensities per band) & 612\,kB \\
    \bottomrule
  \end{tabular}
\end{table}

\subsubsection{Columns}

Table~\ref{tab:dictionary} defines every column. In the column names, \texttt{<b>} stands for one of the 38 band
names of Table~\ref{tab:bands} and \texttt{<k>} for the two-digit cell index \texttt{00}\ldots\texttt{63} along
the line of sight (\texttt{00} nearest the observer). A column with suffix \texttt{\_log} contains $\log_{10}$ of the
column without it. Of the 1443 columns, 1320 are single precision (\texttt{float32}), 114 double precision
(\texttt{float64}: all $I$, $Q$, $U$), 8 integers and one (\texttt{split}) a dictionary-encoded string. The
dataset contains no missing or non-finite values: \texttt{validate\_dataset.py} scanned every value of every
column and found none, so the NaN conventions of the builder (logarithm of a non-positive value, $p$ where
$I\le0$) never apply.

\begingroup
\small
\setlength{\LTcapwidth}{\textwidth}
\begin{longtable}{@{}p{0.27\textwidth}p{0.08\textwidth}p{0.11\textwidth}p{0.46\textwidth}@{}}
  \caption{Data dictionary. Type: f32/f64 = single/double-precision float, i8/i16/i32 = integers. Counts in
  brackets give the number of columns of a pattern.}
  \label{tab:dictionary}\\
  \toprule
  Column & Type & Unit & Definition \\
  \midrule
  \endfirsthead
  \multicolumn{4}{@{}l}{\small\emph{Table~\ref{tab:dictionary} continued}}\\
  \toprule
  Column & Type & Unit & Definition \\
  \midrule
  \endhead
  \midrule
  \multicolumn{4}{r@{}}{\small\emph{continued on next page}}\\
  \endfoot
  \bottomrule
  \endlastfoot
  \multicolumn{4}{@{}l}{\emph{Identifiers} (10 columns)} \\
  \texttt{snapshot} & i16 & -- & STARFORGE snapshot number, 130--500 \\
  \texttt{time\_Myr} & f32 & Myr & simulation time of the snapshot, 0.528--4.957 \\
  \texttt{core\_id} & i32 & -- & ID of the box's central core in the step-1 catalogue of that snapshot \\
  \texttt{box\_id} & i32 & -- & unique box ID, 0--1907 (step 2) \\
  \texttt{family\_id} & i32 & -- & cross-snapshot family, 0--316 (Sect.~\ref{sec:dataset:boxes}) \\
  \texttt{n\_stars} & i16 & -- & number of stellar sources inside the box, 0--39 \\
  \texttt{split} & string & -- & \texttt{train}, \texttt{val} or \texttt{test} \\
  \texttt{view} & i8 & -- & viewing direction 0, 1, 2 (Table~\ref{tab:views}) \\
  \texttt{pix\_row}, \texttt{pix\_col} & i8 & -- & image pixel, 0--63 \\
  \midrule
  \multicolumn{4}{@{}l}{\emph{Sightline quantities} (15 columns)} \\
  \texttt{N\_col} (+\texttt{\_log}) & f32 & cm$^{-2}$ & column density of gas particles from POLARIS's column-density image \\
  \texttt{Sigma} (+\texttt{\_log}) & f32 & g\,cm$^{-2}$ & gas surface density, $N_\mathrm{col}\,\mu m_u$ \\
  \texttt{Sigma\_dust} (+\texttt{\_log}) & f32 & g\,cm$^{-2}$ & dust surface density, $0.01\,\Sigma$ \\
  \texttt{temp\_polaris\_mean} (+\texttt{\_log}) & f32 & K & density-weighted mean POLARIS dust temperature \\
  \texttt{temp\_starforge\_mean} (+\texttt{\_log}) & f32 & K & density-weighted mean STARFORGE dust temperature \\
  \texttt{B\_tot\_mean} (+\texttt{\_log}) & f32 & $\mu$G & density-weighted mean field strength $|\mathbf{B}|$ \\
  \texttt{cos2psi\_B\_mean}, \texttt{sin2psi\_B\_mean} & f32 & -- & direction of $z$ (Eq.~\ref{eq:coherence}): plane-of-sky field orientation seen in polarization \\
  \texttt{B\_coherence} & f32 & -- & $|z|/\sum_k n_k$, 0--1 (Eq.~\ref{eq:coherence}) \\
  \midrule
  \multicolumn{4}{@{}l}{\emph{Bands} (7 columns $\times$ 38 bands = 266)} \\
  \texttt{<b>\_I} & f64 & MJy\,sr$^{-1}$ & Stokes $I$ \\
  \texttt{<b>\_Q}, \texttt{<b>\_U} & f64 & MJy\,sr$^{-1}$ & Stokes $Q$, $U$ in the image frame (angles from $+$col towards $+$row) \\
  \texttt{<b>\_p} & f32 & -- & polarization fraction $\sqrt{Q^2+U^2}/I$ \\
  \texttt{<b>\_polangle} & f32 & deg & polarization (E-vector) angle $\tfrac12\arctan(U,Q)$ in $(-90,90]$; field angle $=$ \texttt{polangle} $+90^\circ$ \\
  \texttt{<b>\_I\_log}, \texttt{<b>\_p\_log} & f32 & -- & $\log_{10}$ of $I$ [MJy\,sr$^{-1}$] and of $p$ \\
  \midrule
  \multicolumn{4}{@{}l}{\emph{Profiles} (64 cells each, $k=00$ nearest the observer; 1152 columns)} \\
  \texttt{dens\_<k>} (+\texttt{\_log}) & f32 & cm$^{-3}$ & gas particle density $n=\rho/(\mu m_u)$, $\mu=2.33$ \\
  \texttt{dens\_dust\_<k>} (+\texttt{\_log}) & f32 & g\,cm$^{-3}$ & dust mass density, $0.01\,\rho$ \\
  \texttt{temp\_polaris\_<k>} (+\texttt{\_log}) & f32 & K & dust temperature computed by POLARIS (used for the images) \\
  \texttt{temp\_starforge\_<k>} (+\texttt{\_log}) & f32 & K & STARFORGE dust temperature (mass-weighted deposit) \\
  \texttt{B\_tot\_<k>} (+\texttt{\_log}) & f32 & $\mu$G & field strength $|\mathbf{B}|$ \\
  \texttt{B\_pos\_<k>} (+\texttt{\_log}) & f32 & $\mu$G & plane-of-sky field strength $\sqrt{B_\mathrm{col}^2+B_\mathrm{row}^2}$ \\
  \texttt{cos2psi\_B\_<k>}, \texttt{sin2psi\_B\_<k>} & f32 & -- & plane-of-sky field orientation $\cos2\psi_B$, $\sin2\psi_B$ \\
  \texttt{sin2gamma\_B\_<k>} & f32 & -- & $\sin^2\gamma_B=B_\mathrm{pos}^2/B_\mathrm{tot}^2$ (1: field in the sky plane; 0: along the line of sight) \\
  \texttt{B\_col\_<k>}, \texttt{B\_row\_<k>} & f32 & $\mu$G & field components along image $+$col and $+$row \\
  \texttt{B\_los\_<k>} & f32 & $\mu$G & line-of-sight component, positive away from the observer \\
\end{longtable}
\endgroup

\subsubsection{Value ranges}

Table~\ref{tab:ranges} gives the distribution of the main quantities. The sightline and band columns were
evaluated over all 23\,445\,504 rows of \texttt{v1}; the profile columns over all 64 cells of 300 randomly chosen
box views ($7.9\times10^{7}$ cells). All column densities in the dataset correspond to optically thin dust at
long wavelengths ($\tau_{850}\le0.01$, Sect.~\ref{sec:dataset:verification}). Intensities at
$\le5.8\,\mu\mathrm{m}$ are of order $10^{-108}\,\mathrm{MJy\,sr^{-1}}$ and carry no usable signal.

\begin{table}[ht]
  \centering
  \caption{Distribution of selected quantities: minimum, 1st percentile, median, 99th percentile and maximum.}
  \label{tab:ranges}
  \small
  \begin{tabular}{@{}lllllll@{}}
    \toprule
    Column & Unit & min & 1\% & median & 99\% & max \\
    \midrule
    \multicolumn{7}{@{}l}{\emph{Sightline (all rows)}} \\
    \texttt{N\_col}               & cm$^{-2}$     & $2.4\times10^{18}$ & $9.9\times10^{19}$ & $9.2\times10^{20}$ & $7.1\times10^{21}$ & $2.5\times10^{23}$ \\
    \texttt{Sigma}                & g\,cm$^{-2}$  & $9.3\times10^{-6}$ & $3.8\times10^{-4}$ & $3.6\times10^{-3}$ & $2.7\times10^{-2}$ & $0.97$ \\
    \texttt{temp\_polaris\_mean}  & K             & 8.27 & 12.5 & 14.5 & 16.0 & 54.4 \\
    \texttt{temp\_starforge\_mean}& K             & 6.16 & 11.3 & 13.3 & 15.5 & 53.8 \\
    \texttt{B\_tot\_mean}         & $\mu$G        & 0.96 & 9.6 & 25.3 & 58.3 & 1760 \\
    \texttt{B\_coherence}         & --            & $3.6\times10^{-5}$ & 0.029 & 0.47 & 0.96 & 0.999 \\
    \midrule
    \multicolumn{7}{@{}l}{\emph{Bands (all rows)}} \\
    $I$ (53\,$\mu$m)   & MJy\,sr$^{-1}$ & $3.2\times10^{-4}$ & 0.0074 & 0.018 & 0.072 & $2.1\times10^{6}$ \\
    $I$ (214\,$\mu$m)  & MJy\,sr$^{-1}$ & 0.083 & 3.3 & 20.5 & 87.4 & $1.2\times10^{5}$ \\
    $I$ (849\,$\mu$m)  & MJy\,sr$^{-1}$ & 0.0039 & 0.17 & 1.33 & 8.1 & 1380 \\
    $I$ (1380\,$\mu$m) & MJy\,sr$^{-1}$ & $8.0\times10^{-4}$ & 0.034 & 0.28 & 1.8 & 246 \\
    $p$ (53\,$\mu$m)   & --  & $1.4\times10^{-5}$ & 0.0059 & 0.098 & 0.259 & 0.355 \\
    $p$ (214\,$\mu$m)  & --  & $1.0\times10^{-5}$ & 0.0052 & 0.089 & 0.260 & 0.323 \\
    $p$ (849\,$\mu$m)  & --  & $2.2\times10^{-6}$ & 0.0051 & 0.089 & 0.261 & 0.324 \\
    $I$ (3.4\,$\mu$m)  & MJy\,sr$^{-1}$ & $7\times10^{-136}$ & $3\times10^{-112}$ & $8\times10^{-108}$ & $1\times10^{-104}$ & $3\times10^{-12}$ \\
    \midrule
    \multicolumn{7}{@{}l}{\emph{Profiles (300 box views, all cells)}} \\
    \texttt{dens}           & cm$^{-3}$    & 0.010 & 5.8 & 73.8 & 2310 & $9.9\times10^{5}$ \\
    \texttt{dens\_dust}     & g\,cm$^{-3}$ & $3.9\times10^{-28}$ & $2.3\times10^{-25}$ & $2.9\times10^{-24}$ & $8.9\times10^{-23}$ & $3.8\times10^{-20}$ \\
    \texttt{temp\_polaris}  & K            & 8.37 & 12.7 & 15.0 & 16.1 & 80.0 \\
    \texttt{temp\_starforge}& K            & 5.91 & 11.3 & 13.4 & 15.6 & 66.8 \\
    \texttt{B\_tot}         & $\mu$G       & 0.037 & 7.3 & 26.0 & 69.1 & 2380 \\
    \texttt{B\_pos}         & $\mu$G       & 0.0017 & 2.7 & 20.2 & 60.2 & 2370 \\
    \texttt{sin2gamma\_B}   & --           & $2.6\times10^{-8}$ & 0.019 & 0.77 & 1.00 & 1.00 \\
    \texttt{B\_los}         & $\mu$G       & $-1600$ & $-45.8$ & $-6.6$ & 31.8 & 1320 \\
    \bottomrule
  \end{tabular}
\end{table}

The intensity spectra of all sightlines peak between 160 and $250\,\mu\mathrm{m}$ (at $214\,\mu\mathrm{m}$ for
96.8\%), as expected for dust at 12--16\,K. The polarization fraction changes little with wavelength between
$53\,\mu\mathrm{m}$ and $1.38\,\mathrm{mm}$ (median 0.098 at $53\,\mu\mathrm{m}$, 0.089 at 214 and
$849\,\mu\mathrm{m}$; maximum 0.355).

\subsubsection{Reading the data}

Columns can be read selectively, which keeps memory small; for example, with \texttt{pandas} and \texttt{pyarrow}:
\begin{verbatim}
import pandas as pd, pyarrow.dataset as ds
root = "/scratch/sjafaritabar/starforge_polaris_dataset_v1/data"
cols = ["planck_hfi_353ghz_I_log", "planck_hfi_353ghz_p"] + \
       [f"dens_{k:02d}_log" for k in range(64)]
train = pd.read_parquet(f"{root}/train", columns=cols)
full = ds.dataset(root, format="parquet", partitioning=None)   # lazy, all splits
\end{verbatim}
To obtain images, a file is read one row group at a time (\texttt{pyarrow.parquet.ParquetFile.read\_row\_group})
and the 4096 rows are placed at (\texttt{pix\_row}, \texttt{pix\_col}).

% -----------------------------------------------------------------------------------------------------
\subsection{Verification}
\label{sec:dataset:verification}

Each step has unit tests (step 2: 7, step 3: 13, step 5: 4; all pass), and an independent audit
(\texttt{audit\_pipeline.py}) re-derives every stage from the raw files without using the pipeline's code. It
passed all 15 checks for both dataset versions:
\begin{enumerate}
  \item Every boxed core is in its step-1 catalogue with the same peak, peak particle and mass; box times equal
  the snapshot header times; no two boxes of a snapshot overlap; families re-derived from scratch reproduce all
  317 \texttt{family\_id}s.
  \item Exactly one grid per box; for every box the summaries, grids, star lists, command files (photon numbers,
  radius, dust, wavelengths, views) and POLARIS outputs are mutually consistent, the grid density is unchanged by
  POLARIS, and dust temperatures and alignment radii are within range.
  \item The $850\,\mu\mathrm{m}$ intensity in $\mathrm{MJy\,sr^{-1}}$ matches the optically thin model
  $\sum_k B_\nu(T_k)\,\Delta\tau_k$ computed from the grid, which tests the units, pixel solid angle and distance:
  POLARIS/model has a per-box median of $1.002$ (range $1.001$--$1.005$); $\tau_{850}\le0.01$.
  \item No box has an outlying POLARIS-vs-STARFORGE dust-temperature offset ($>5\sigma$).
  \item Every one of the 1443 columns was recomputed from the raw grids and FITS files with explicit index formulas
  (Table~\ref{tab:views}) for 1080 random sightlines in 36 boxes ($1.56\times10^{6}$ values); the worst
  relative (or, for logarithms and angles, absolute) error is $3.8\times10^{-6}$, i.e.\ single-precision rounding.
  \item \texttt{splits.csv} covers every box once with the step-2 family IDs; no family is in more than one split;
  all families with more than 50 boxes are in training.
  \item Polarization fractions at $\ge50\,\mu\mathrm{m}$ are physical (maximum 0.355), and the spectra of all
  sightlines peak at 60--$500\,\mu\mathrm{m}$ (214\,$\mu$m for 96.8\%).
\end{enumerate}
\texttt{validate\_dataset.py} additionally confirms the row counts, the absence of duplicate sightlines and of
boxes or families in more than one split, and scans every column of every row for non-finite values.

Across the whole dataset, column densities range over $2.4\times10^{18}$--$2.5\times10^{23}\,\mathrm{cm^{-2}}$
(median $9.1\times10^{20}$), density-weighted POLARIS dust temperatures over 8.3--54.4\,K (median 14.5\,K),
density-weighted field strengths over $0.96$--$1760\,\mu\mathrm{G}$ (median $25.7\,\mu\mathrm{G}$), and
\texttt{B\_coherence} over $4\times10^{-5}$--$0.999$ (median 0.48).

% -----------------------------------------------------------------------------------------------------
\subsection{Limitations}
\label{sec:dataset:limitations}

\begin{itemize}
  \item \textbf{Radiation field.} Each box is illuminated by an unattenuated Mathis field with $G_0=1$. In
  STARFORGE the field reaching the box surfaces is attenuated by the surrounding cloud (snapshot 440: equivalent
  $G_0$ median 0.20, 5--95\% range 0.11--0.46). POLARIS dust temperatures are therefore higher than the STARFORGE
  ones, by a per-box median of $+1.49\,\mathrm{K}$ (5--95\%: $+0.29$ to $+2.76\,\mathrm{K}$). Both temperatures
  are in the dataset.
  \item \textbf{Isolated boxes.} Only gas inside the $2.2\,\mathrm{pc}$ cube emits and absorbs; foreground and
  background cloud emission and extinction are absent.
  \item \textbf{No stochastic heating}: bands at $\le25\,\mu\mathrm{m}$ are not reliable.
  \item \textbf{Ideal observations.} No noise, no beam, a single distance (140\,pc) and a pixel equal to a grid
  cell. Noise and beam have to be added when training.
  \item \textbf{No gas kinematics.} Velocities are not included, so velocity-based field-strength estimates
  (e.g.\ Davis--Chandrasekhar--Fermi) cannot be reproduced from the dataset alone.
  \item \textbf{One simulation.} All boxes come from one cloud; families make the evaluation sets independent in
  gas, not in environment. Late snapshots, which contain all stars, are $0.41\,\mathrm{Myr}$ apart.
  \item \textbf{Monte Carlo.} Dust temperatures carry $\approx0.008\,\mathrm{K}$ of Monte Carlo noise, and runs
  are not bitwise reproducible.
\end{itemize}

% -----------------------------------------------------------------------------------------------------
\subsection{Data products and code}
\label{sec:dataset:files}

Code: \texttt{STARFORGE/creating\_the\_dataset/} (\texttt{1.finding\_cores}, \texttt{2.selecting\_boxes},
\texttt{3.preparing\_polaris\_grids} with its \texttt{diagnostics}, \texttt{4.running\_polaris},
\texttt{5.building\_dataset}, \texttt{audit\_pipeline.py}); photon studies in
\texttt{STARFORGE/M2e5\_R30\_Res585/POLARIS/isrf\_photon\_noise\_test} and \texttt{STARFORGE/M2e5\_R30\_Res585/POLARIS/star\_photon\_test}.
Data: snapshots in \texttt{/scratch/STARFORGE/M2e5\_R30\_Res585}; core catalogues in
\texttt{1.finding\_cores/dendrogram\_cores\_v5.2}; boxes in \texttt{2.selecting\_boxes/boxes\_2.2pc};
grids in \texttt{3.preparing\_polaris\_grids/grids\_64}; POLARIS output in
\texttt{/scratch/sjafaritabar/polaris\_runs} ($\approx70\,\mathrm{GB}$); datasets in
\texttt{/scratch/sjafaritabar/starforge\_polaris\_dataset\_v1} and
\texttt{/scratch/sjafaritabar/starforge\_polaris\_dataset\_v1\_stars\_eval}; audit reports in
\texttt{/scratch/sjafaritabar/audit}.
. Required packages in the preamble:
%   \usepackage{amsmath}  \usepackage{booktabs}  \usepackage{graphicx}  \usepackage{longtable}
% References are in dataset_construction.bib (\bibliography{dataset_construction}).
% Every number below is taken from the pipeline code, its logs, manifests and the audit report
% (paths in Sect. "Data products and code"); dates are UTC.
% =====================================================================================================

\section{A synthetic dust-emission and polarization dataset of star-forming cores}
\label{sec:dataset}

This section describes how the training data for this thesis were constructed. Starting from a
magnetohydrodynamic simulation of a star-forming molecular cloud, dense cores were identified in every
simulation snapshot, a $2.2\,\mathrm{pc}$ cube around each selected core was converted into a radiative-transfer
model, and the dust temperature, grain alignment and polarized dust emission of every cube were computed with
the Monte Carlo code POLARIS. The resulting images were combined with the three-dimensional state of the gas
into a table with one row per line of sight. Table~\ref{tab:pipeline} summarises the five steps.

\begin{table}[ht]
  \centering
  \caption{The five steps of the pipeline. Each step is a separate script; versions refer to the
  version strings recorded in each step's output.}
  \label{tab:pipeline}
  \small
  \begin{tabular}{@{}llll@{}}
    \toprule
    Step & Task & Product & Code (version) \\
    \midrule
    1 & Dense cores in every snapshot   & 11\,599 cores in 51 snapshots & \texttt{find\_starforge\_cores\_dendrogram\_v5.py} (5.2.0) \\
    2 & Non-overlapping boxes, families & 1908 boxes, 317 families     & \texttt{select\_boxes.py} \\
    3 & POLARIS grids and commands      & 1908 grids of $64^3$ cells    & \texttt{prepare\_grids.py} (3.5.1) \\
    4 & Radiative transfer              & 1908 runs, 217\,512 images    & \texttt{run\_polaris.py}, POLARIS \\
    5 & Sightline table                 & 23\,445\,504 rows, 1443 columns & \texttt{build\_dataset.py} (1.1.1 / 1.2.0) \\
    \bottomrule
  \end{tabular}
\end{table}

% -----------------------------------------------------------------------------------------------------
\subsection{The STARFORGE simulation}
\label{sec:dataset:sim}

The gas data come from the STARFORGE \citep{Grudic2021} simulation \texttt{M2e5\_R30\_Res585}, run with the
meshless magnetohydrodynamics code GIZMO \citep{Hopkins2015} (header \texttt{GIZMO\_version} 2022). The initial
condition is a uniform-density spherical cloud of $2\times10^{5}\,M_\odot$ and radius $30\,\mathrm{pc}$ (measured:
$1.993\times10^{5}\,M_\odot$ inside $R=30\,\mathrm{pc}$, median density $1.19\times10^{-22}\,\mathrm{g\,cm^{-3}}$)
embedded in a periodic box of side $300\,\mathrm{code\ units}=300.42\,\mathrm{pc}$. The initial three-dimensional
velocity dispersion of the cloud is $5.86\,\mathrm{km\,s^{-1}}$, which corresponds to a virial parameter
$\alpha=2$; the initial gas temperature is $\approx12\,\mathrm{K}$; gas cells have a median mass of
$1.001\times10^{-3}\,M_\odot$ (range $1.001$--$1.301\times10^{-3}\,M_\odot$). The interstellar radiation field
strength parameter in the snapshot header is $1$; solar abundances are $Z_\odot=0.0142$, $Y_\odot=0.2703$.

\paragraph{Units.} The code units are $1.0014\,\mathrm{pc}$ (length), $1.0008\,M_\odot$ (mass) and
$1\,\mathrm{m\,s^{-1}}$ (velocity), so one code time unit is $979.16\,\mathrm{Myr}$. The units were verified
against the initial conditions (\texttt{check\_snapshot\_units.py}): with these units the gravitational constant
derived from the header units agrees with the header value to $3.4\times10^{-4}$, the enclosed mass and density
match the nominal cloud, and the velocity dispersion matches $\alpha=2$. The saved \texttt{MagneticField} is in
tesla, \emph{not} in the unit given by the header field \texttt{Internal\_UnitB\_In\_Gauss}
($2.91\times10^{-9}\,\mathrm{G}$): read in tesla, the initial cloud field is $7.19\,\mu\mathrm{G}$, consistent
(within a factor $1.26$, the uncertainty of the critical mass-to-flux normalisation) with the nominal
normalised mass-to-flux ratio $4.2$, whereas the gauss or header-unit readings are off by four or more
orders of magnitude. POLARIS expects the field in tesla, so the field is passed on unconverted.

\paragraph{Snapshots.} Snapshots \texttt{snapshot\_000} to \texttt{snapshot\_500} in steps of 10 were analysed
(51 snapshots, each $\approx 31\,\mathrm{GB}$ with $\approx 2.5\times10^{8}$ gas cells). Dense cores exist from
snapshot 130 onwards. Snapshots 130--420 are $0.0406\,\mathrm{Myr}$ apart (0.528--1.707\,Myr); snapshots
420--500 are $0.406\,\mathrm{Myr}$ apart (1.707--4.957\,Myr). Sink particles (stars) exist only in the late
snapshots.

% -----------------------------------------------------------------------------------------------------
\subsection{Step 1: dense cores}
\label{sec:dataset:cores}

\paragraph{Molecular density.} For every gas cell the molecular hydrogen number density is
\begin{equation}
  n_{\mathrm{H_2}} = \frac{\rho\, X_\mathrm{H}\, f_\mathrm{neutral}\, f_\mathrm{mol}}{2\,m_\mathrm{H}},\qquad
  X_\mathrm{H} = 1 - Z - Y,
  \label{eq:nh2}
\end{equation}
where $\rho$ is the cell density, $Z$ and $Y$ are the first two columns of \texttt{Metallicity},
$f_\mathrm{neutral}$ is \texttt{NeutralHydrogenAbundance}, $f_\mathrm{mol}$ is \texttt{MolecularMassFraction},
and $m_\mathrm{H}=1.6735575\times10^{-24}\,\mathrm{g}$. Invalid or out-of-range inputs abort the run instead of
being clipped.

\paragraph{Dendrogram.} Cells with $n_{\mathrm{H_2}} > 10^{4}\,\mathrm{cm^{-3}}$ were selected. Two selected cells
are neighbours if either lies within the kernel support radius $h$ (\texttt{SmoothingLength}) of the other,
using periodic distances; the resulting graph was symmetrised. A dendrogram \citep{Rosolowsky2008} was computed
with \texttt{astrodendro} 0.3.1 on this particle graph with
$\texttt{min\_value}=10^{4}\,\mathrm{cm^{-3}}$, $\texttt{min\_delta}=10^{4}\,\mathrm{cm^{-3}}$ and
$\texttt{min\_npix}=100$ cells. Every leaf of the dendrogram, i.e.\ a local maximum of $n_{\mathrm{H_2}}$
with at least 100 cells and a contrast of at least $10^{4}\,\mathrm{cm^{-3}}$ above the level at which it
merges with a neighbour, is a \emph{core}. A leaf is a density peak; it is not required to be gravitationally
bound or associated with a star.

\paragraph{Core properties.} For each core the catalogue records, among others, the peak cell (position,
$n_{\mathrm{H_2}}$, particle ID), mass, volume and radii, shape, internal velocity dispersion, mean gas and dust
temperature, mass-weighted field, the gas-only virial ratio $2(K+U)/|W|$ (direct pair summation with
$10^{-3}\,\mathrm{pc}$ Plummer softening, up to $2\times10^{4}$ cells), and the sink particles near the core,
including the number of sinks within a $2.2\,\mathrm{pc}$ cube centred on the peak.
Step 1 found 11\,599 cores in total (none before snapshot 130; 1038 in snapshot 500) with a total mass of
$4535\,M_\odot$.

% -----------------------------------------------------------------------------------------------------
\subsection{Step 2: boxes and families}
\label{sec:dataset:boxes}

\paragraph{Box selection.} A \emph{box} is an axis-aligned cube of side $S=2.2\,\mathrm{pc}$ centred on a core's
density peak. Within each snapshot, boxes were selected by greedy non-maximum suppression: cores were sorted by
mass (ties, which are common because core masses are multiples of the cell mass, broken by higher peak density,
then lower \texttt{core\_id}), and a core was kept unless its cube overlapped an already kept cube. Two cubes
overlap if the periodic Chebyshev distance between their centres is $<S$. This gives 1908 boxes in 38 snapshots
(130--500; 1 to 104 boxes per snapshot, median 55). The central cores have masses of $0.10$--$9.09\,M_\odot$
(median $0.41\,M_\odot$) and peak densities $n_{\mathrm{H_2}}=2.0\times10^{4}$--$7.0\times10^{9}\,\mathrm{cm^{-3}}$
(median $3.9\times10^{4}\,\mathrm{cm^{-3}}$). Of all cores, 3604 ($1249\,M_\odot$) lie in no box: they were
suppressed by a more massive neighbour without lying inside its cube.

\paragraph{Families.} Boxes of consecutive snapshots largely image the same gas. Two boxes in consecutive
snapshots (consecutive among snapshots that contain boxes) were therefore linked if their cubes overlap
(periodic Chebyshev distance $<S$), and the connected components of these links, found with a union--find
structure, were labelled \texttt{family\_id}. The 1908 boxes form 317 families; the median family has 2 boxes,
the 90th percentile 9.4, the 99th percentile 107, and the largest family has 260 boxes. Six families have more
than 50 boxes. Families are the unit of the train/validation/test split (Sect.~\ref{sec:dataset:splits}).

% -----------------------------------------------------------------------------------------------------
\subsection{Step 3: radiative-transfer models}
\label{sec:dataset:grids}

\paragraph{Grid geometry.} Each box was converted into a POLARIS octree grid that is a uniform cube of
$64^{3}$ cells (octree level 6) of size $\Delta=2.2/64=0.034375\,\mathrm{pc}$. The grid is shifted by half a
cell so that the core's peak cell lies at the centre of cell $(32,32,32)$; centring the cube exactly on the peak
would place the peak on a corner shared by eight cells. Each snapshot was read once and the gas of all its boxes
was cut out in that single pass; boxes were then built in parallel.

\paragraph{Deposition.} Every gas cell whose kernel overlaps the cube (centre within $S/2+h$ of the box centre
along each axis) was deposited with GIZMO's compact cubic-spline kernel,
\begin{equation}
  W(q) \propto \begin{cases} 1-6q^{2}+6q^{3}, & 0\le q<\tfrac12,\\ 2(1-q)^{3}, & \tfrac12\le q<1,\\ 0, & q\ge 1,\end{cases}
  \qquad q = \frac{|\mathbf{x}_c-\mathbf{x}_i|}{h_i},
\end{equation}
sampled at the centres $\mathbf{x}_c$ of all grid cells within the kernel, \emph{including cells outside the cube}.
The weights $w_{ic}=W(q_{ic})/\sum_{c'}W(q_{ic'})$ are normalised over the full kernel and only the in-cube part
is kept, so each grid holds exactly the mass inside the cube and the faces are not under-dense. If a kernel
contains no cell centre ($h_i$ well below $\Delta$), the cell's mass goes to the grid cell containing its centre
(this applied to 1.07\% of the $2.73\times10^{8}$ deposited gas cells). Cell quantities are
\begin{equation}
  \rho_c = \frac{\sum_i m_i w_{ic}}{\Delta^{3}},\qquad
  \langle X\rangle_c = \frac{\sum_i m_i w_{ic} X_i}{\sum_i m_i w_{ic}},
\end{equation}
with $X$ the dust temperature, gas temperature and the three field components (mass-weighted means). A cell
receiving no mass would get a density floor of $10^{-6}$ times the smallest occupied density and the fields of
the nearest occupied cell; this happened for a single cell in a single box ($3.6\times10^{-14}\,M_\odot$).
The grid mass agrees with the mass of the gas cells centred inside the cube to between $-0.12\%$ and $+0.08\%$.
Boxes contain $3.5\times10^{4}$--$7.9\times10^{5}$ contributing gas cells (median $1.3\times10^{5}$) and
$27.6$--$748.6\,M_\odot$ of gas (median $113\,M_\odot$).

\paragraph{Grid fields.} The grid stores (POLARIS field IDs in brackets) the gas number density
$n=\rho/(\mu m_u)$ in $\mathrm{m^{-3}}$ [0], the STARFORGE dust temperature [2], the gas temperature [3] and
$B_x, B_y, B_z$ in tesla [4--6]. The POLARIS branch used reads the gas density only as a number density and
converts it back with $\rho=n\,\mu\,m_u$, where $m_u=1.66053904\times10^{-27}\,\mathrm{kg}$ (POLARIS's
\texttt{m\_H}); the same $n$ and $\mu$ enter the radiative-torque calculation (gas damping, thermal speed).
$\mu$ is therefore the mean mass per gas particle. We use $\mu=2.33$: the mass-weighted median of the particle
mean molecular weight of box gas, computed from the snapshot chemistry, is 2.32--2.35 for
$n_\mathrm{H}=10^{3}$--$10^{6}\,\mathrm{cm^{-3}}$ (diagnostic \texttt{check\_mu.py}, snapshot 500), and
$\rho$, hence the dust mass, is reproduced exactly for any $\mu$.

\paragraph{Dust.} One dust component was used: astrodust \citep{Hensley2023} without iron inclusions
($f_\mathrm{Fe}=0$), with the catalogue's magnetic susceptibility $\chi=4.47\times10^{-4}$ and material density
$2898\,\mathrm{kg\,m^{-3}}$, a power-law size distribution $\mathrm{d}n/\mathrm{d}a\propto a^{-3.5}$
\citep{Mathis1977} between $a_\mathrm{min}=5\,\mathrm{nm}$ and $a_\mathrm{max}=250\,\mathrm{nm}$, and a
dust-to-gas mass ratio of $0.01$ (POLARIS \texttt{<mass\_fraction>}). The ratio is the one STARFORGE itself
uses: its field \texttt{DustToGasRatio\_Local} stores $Z$ times the local dust-to-metals ratio relative to solar,
with a median of $0.01416=Z_\odot$ over all boxes, which corresponds to an absolute dust-to-gas ratio
$0.01\,Z/Z_\odot = 0.01$. The catalogue covers $1\,\mathrm{nm}$--$1\,\mu\mathrm{m}$ in size and
$0.05\,\mu\mathrm{m}$--$5\,\mathrm{mm}$ in wavelength, i.e.\ POLARIS's internal wavelength grid
(100 points, $0.0912$--$2000\,\mu\mathrm{m}$) and all detector wavelengths. With this susceptibility the
Larmor-precession limit of POLARIS's alignment model is never reached: the smallest Larmor-limited grain size
over all $5.0\times10^{8}$ grid cells is $408\,\mathrm{nm}>a_\mathrm{max}$ (\texttt{check\_larmor\_limit.py}).

\paragraph{Radiation sources.} Every box is illuminated by the interstellar radiation field (ISRF) of
\citet{Mathis1983} with strength $G_0=1$, launched from a sphere of 10 times the circumscribed radius of the grid
($10\times\tfrac{\sqrt3}{2}\times2.2\,\mathrm{pc}=19.1\,\mathrm{pc}$). In addition, every sink particle inside the
cube with positive radius and luminosity was added as a blackbody point source at its position, with the
STARFORGE protostellar radius $R_\ast$ and an effective temperature from
$L_\ast = 4\pi R_\ast^{2}\sigma_\mathrm{SB}T_\mathrm{eff}^{4}$. Fifty-two boxes contain stars, 275 sources in total
(up to 39 in one box): $L_\ast=6.8\times10^{-3}$--$1.0\times10^{4}\,L_\odot$ (median $0.16\,L_\odot$),
$R_\ast=0.31$--$20.8\,R_\odot$, $T_\mathrm{eff}=2990$--$27\,900\,\mathrm{K}$, sink masses
$0.02$--$11.7\,M_\odot$. (Step 1 counted 273 sinks in these boxes; in one box, 1851, the half-cell shift of the
grid moves two further sinks inside the cube.) Boxes without stars are heated by the ISRF only.

\paragraph{Photon numbers and Monte Carlo noise.} With radiative-torque alignment, POLARIS launches the
requested number of photons for each of its 100 wavelengths. Production used $10^{7}$ ISRF photons and $10^{6}$
photons per star (per wavelength). These values were chosen from two convergence studies, each running every
configuration twice with independent random numbers and taking the per-cell noise as
$\mathrm{median}\,|x_1-x_2|/\sqrt2$:
(i) ISRF photons $10^{6}$--$10^{7}$ and launch radii 3--30 on box \texttt{snapshot\_440/core\_000395}; at $10^{7}$
photons and radius 10 the dust-temperature noise is $0.0078\,\mathrm{K}$ in the interior, $0.0087\,\mathrm{K}$ on
faces, $0.0103\,\mathrm{K}$ on edges and $0.0137\,\mathrm{K}$ in corners ($0.0081\,\mathrm{K}$ overall), the
alignment radius varies by $0.63\%$, and at $850\,\mu\mathrm{m}$ the noise is $0.026\%$ in $I$, $0.20\%$ in $p$
and $0.035^\circ$ in the polarization angle. Smaller radii concentrate the noise in the corners
($0.073\,\mathrm{K}$ at radius 3), larger radii raise it in the interior ($0.0105\,\mathrm{K}$ at radius 30).
(ii) Star photons $10^{5}$ vs.\ $10^{6}$ on box \texttt{snapshot\_500/core\_000873}: $10^{6}$ halves the
temperature noise in gas below $25\,\mathrm{K}$ (from 0.077--0.084 to 0.039--0.045\,K) and changes the
$53\,\mu\mathrm{m}$ intensity by a median of $2.0\%$ ($0.22\%$ at $850\,\mu\mathrm{m}$).

\paragraph{POLARIS tasks.} Each box has a command file with two tasks, both with radiative-torque alignment
(\texttt{ALIG\_RAT}) and 8 threads:
\begin{enumerate}
  \item \texttt{CMD\_TEMP\_RAT}: Monte Carlo radiative transfer of the ISRF and stellar photons; computes the
  equilibrium dust temperature of every cell and the radiative-torque alignment radius $a_\mathrm{alig}$ (grains
  larger than $a_\mathrm{alig}$ are aligned with the magnetic field). The STARFORGE dust temperature in the input
  grid is replaced by this temperature. Writes \texttt{temp/grid\_temp\_rat.dat}.
  \item \texttt{CMD\_DUST\_EMISSION} on that grid: ray-traced thermal emission of the aligned grains, giving the
  Stokes $I,Q,U,V$, the optical depth and the column density for 38 wavelengths and 3 viewing directions
  (114 detectors) with $64\times64$ pixels each (\texttt{max\_subpixel\_lvl} 0). The detector covers the grid, so
  one pixel corresponds to one grid column ($0.034375\,\mathrm{pc}$); the source distance is
  $140\,\mathrm{pc}$, so a pixel subtends $50.6''$ and
  $\Omega_\mathrm{pix}=(\Delta/d)^{2}=6.03\times10^{-8}\,\mathrm{sr}$.
\end{enumerate}
The 38 wavelengths (Table~\ref{tab:bands}) are those of common far-infrared to millimetre instruments and are
read from one table (\texttt{wavelength\_telescope\_mapping.csv}) by both the grid preparation and the dataset
builder. The three viewing directions are the POLARIS detector rotations
$(\vartheta_1,\vartheta_2)=(0^\circ,0^\circ)$, $(-90^\circ,0^\circ)$ and $(0^\circ,90^\circ)$, which look along the
grid $z$, $y$ and $x$ axes (Sect.~\ref{sec:dataset:geometry}). POLARIS was the \texttt{experimental} branch,
commit \texttt{6c5e1a89c6e14ce806c47bbb87947b98e3f761b0} \citep{Reissl2016}.

% -----------------------------------------------------------------------------------------------------
\subsection{Step 4: radiative-transfer runs}
\label{sec:dataset:runs}

All 1908 boxes were run between 2026-10-31 22:47 and 2026-11-09 12:41 on one machine (AMD EPYC 9554P, 64 cores
/ 128 hardware threads, $377\,\mathrm{GB}$ memory), each box in its own directory, at the lowest scheduling
priority (\texttt{nice 19}), with up to 16 boxes at a time (8 threads each). The runner measured the CPU use of
other users every $15\,\mathrm{s}$ and kept 8 hardware threads free above their peak use over the previous
10 minutes; when the machine was needed, the most recently started boxes were suspended (\texttt{SIGSTOP}) and
resumed 10 minutes after the demand had gone, so no run was ever killed or restarted. Boxes containing stars
were run first. A box took $0.92$--$5.98\,\mathrm{h}$ of wall-clock time (median $1.46\,\mathrm{h}$ at up to 14
concurrent boxes), $2810$ box-hours in total, of which boxes spent $5.3\,\mathrm{h}$ suspended.

Every run was validated before being accepted: the temperature grid exists and all dust temperatures are finite
and within $2$--$3000\,\mathrm{K}$; all 114 detector files exist; $I$, $Q$ and $U$ are finite everywhere; and at
wavelengths $\geq50\,\mu\mathrm{m}$, $I>0$ in every pixel and $Q,U$ are not identically zero. All 1908 runs
finished with exit code 0 and passed; none failed. POLARIS's Monte Carlo is seeded by thread, and the assignment
of photons to threads varies between runs, so runs are statistically but not bitwise reproducible.

% -----------------------------------------------------------------------------------------------------
\subsection{Step 5: the sightline dataset}
\label{sec:dataset:table}

\subsubsection{Geometry and orientation}
\label{sec:dataset:geometry}

Grid arrays are indexed $[i_x,i_y,i_z]$ and images $[\mathrm{row},\mathrm{col}]$, both from 0 to 63. The mapping
between them was determined empirically (\texttt{check\_view\_geometry.py}): for each view, the POLARIS
column-density image was compared with $\sum n\,\Delta$ along each grid axis in all 48 combinations of
line-of-sight axis, transposition and flips; the best orientation matches to a relative difference of
$\approx10^{-13}$. Which end of the line of sight faces the observer cannot be seen in optically thin emission; it
was determined where the dust is optically thick ($21\,\mu\mathrm{m}$, pixels with $\tau>1$), where
$I=\sum_k B_\nu(T_k)\,\Delta\tau_k\,e^{-\tau_k}$ depends on the side from which the absorption is counted
(\texttt{check\_observer\_side.py}). The model built from POLARIS's own temperatures reproduces the image with a
near-constant ratio only if the high-index end of the line-of-sight axis faces the observer, which holds in every
view. Table~\ref{tab:views} gives the result.

\begin{table}[ht]
  \centering
  \caption{Mapping between a dataset row (view $v$, pixel row $r$, pixel column $c$, profile index $k$, with $k=0$
  the cell nearest the observer) and the grid, and the field components in the image frame.}
  \label{tab:views}
  \small
  \begin{tabular}{@{}cccccc@{}}
    \toprule
    View & POLARIS rotation & Line of sight & Image $+$col, $+$row & Grid cell $(i_x,i_y,i_z)$ & $(B_\mathrm{col},B_\mathrm{row},B_\mathrm{los})$ \\
    \midrule
    0 & $(0^\circ,0^\circ)$   & $z$ & $+x$, $+y$ & $(c,\ r,\ 63-k)$    & $(B_x,\ B_y,\ -B_z)$ \\
    1 & $(-90^\circ,0^\circ)$ & $y$ & $+x$, $-z$ & $(c,\ 63-k,\ 63-r)$ & $(B_x,\ -B_z,\ -B_y)$ \\
    2 & $(0^\circ,90^\circ)$  & $x$ & $-z$, $+y$ & $(63-k,\ r,\ 63-c)$ & $(-B_z,\ B_y,\ -B_x)$ \\
    \bottomrule
  \end{tabular}
\end{table}

Angles in the image plane are measured from the image $+$col axis towards $+$row. With Stokes $U$ as written by
POLARIS, the plane-of-sky field direction is the polarization angle plus $90^\circ$: comparing
$\tfrac12\arctan(U,Q)+90^\circ$ with a density-weighted Stokes model of the grid field in this frame (in the more
polarized half of the pixels) gives median differences of $0.2^\circ$--$0.8^\circ$, against $22^\circ$--$54^\circ$ if the
sign of $U$ is reversed. $B_\mathrm{los}>0$ points away from the
observer. Every box and view was checked during the build: $\sum_k n_k\Delta$ along the rotated profiles
reproduces POLARIS's column-density image to a maximum relative difference of $5.4\times10^{-13}$
(the build aborts above $10^{-5}$).

\subsubsection{Contents}

One row is one line of sight: one pixel of one view of one box. The dataset has
$1908\times3\times64\times64=23\,445\,504$ rows and 1443 columns (Table~\ref{tab:columns}).

\begin{table}[ht]
  \centering
  \caption{Columns of the dataset. $k=00\ldots63$ runs along the line of sight from the observer; columns
  marked (+log) also exist as \texttt{<name>\_log} $=\log_{10}$ (NaN where the value is $\le0$). Values are
  \texttt{float32} except $I,Q,U$ (\texttt{float64}) and the integer metadata.}
  \label{tab:columns}
  \small
  \begin{tabular}{@{}p{0.17\textwidth}p{0.47\textwidth}p{0.12\textwidth}r@{}}
    \toprule
    Group & Columns & Unit & Number \\
    \midrule
    metadata & \texttt{snapshot, time\_Myr, core\_id, box\_id, family\_id, n\_stars, split, view, pix\_row, pix\_col} & & 10 \\
    sightline & \texttt{N\_col, Sigma, Sigma\_dust} (+log) & $\mathrm{cm^{-2}}$, $\mathrm{g\,cm^{-2}}$ & 6 \\
              & \texttt{temp\_polaris\_mean, temp\_starforge\_mean} (+log) & K & 4 \\
              & \texttt{B\_tot\_mean} (+log) & $\mu$G & 2 \\
              & \texttt{cos2psi\_B\_mean, sin2psi\_B\_mean, B\_coherence} & -- & 3 \\
    bands ($\times38$) & \texttt{<band>\_I, \_Q, \_U} & MJy\,sr$^{-1}$ & 114 \\
              & \texttt{<band>\_p, \_polangle, \_I\_log, \_p\_log} & --, deg & 152 \\
    profiles ($\times64$) & \texttt{dens\_k, dens\_dust\_k} (+log) & $\mathrm{cm^{-3}}$, $\mathrm{g\,cm^{-3}}$ & 256 \\
              & \texttt{temp\_polaris\_k, temp\_starforge\_k} (+log) & K & 256 \\
              & \texttt{B\_tot\_k, B\_pos\_k} (+log) & $\mu$G & 256 \\
              & \texttt{cos2psi\_B\_k, sin2psi\_B\_k, sin2gamma\_B\_k} & -- & 192 \\
              & \texttt{B\_col\_k, B\_row\_k, B\_los\_k} & $\mu$G & 192 \\
    \midrule
    total & & & 1443 \\
    \bottomrule
  \end{tabular}
\end{table}

\paragraph{Profiles.} For each cell $k$ along the line of sight:
$\texttt{dens}=n$ in $\mathrm{cm^{-3}}$ (gas particles, $\mu=2.33$);
$\texttt{dens\_dust}=0.01\,n\,\mu\,m_u$ in $\mathrm{g\,cm^{-3}}$;
$\texttt{temp\_polaris}$, the dust temperature computed by POLARIS; $\texttt{temp\_starforge}$, the deposited
STARFORGE dust temperature; the field in $\mu$G ($1\,\mathrm{T}=10^{10}\,\mu\mathrm{G}$) in the image frame
(Table~\ref{tab:views}) and the derived quantities
\begin{gather}
  B_\mathrm{pos}^{2}=B_\mathrm{col}^{2}+B_\mathrm{row}^{2},\qquad B_\mathrm{tot}^{2}=B_\mathrm{pos}^{2}+B_\mathrm{los}^{2},\\
  \cos2\psi_B=\frac{B_\mathrm{col}^{2}-B_\mathrm{row}^{2}}{B_\mathrm{pos}^{2}},\qquad
  \sin2\psi_B=\frac{2B_\mathrm{col}B_\mathrm{row}}{B_\mathrm{pos}^{2}},\qquad
  \sin^{2}\gamma_B=\frac{B_\mathrm{pos}^{2}}{B_\mathrm{tot}^{2}},
\end{gather}
where $\psi_B$ is the plane-of-sky field angle and $\gamma_B$ the angle between the field and the line of sight.
The double-angle form removes the $180^\circ$ ambiguity of the field direction, which dust polarization cannot
measure.

\paragraph{Sightline quantities.} $N_\mathrm{col}$ is POLARIS's column-density image converted from
$\mathrm{m^{-2}}$ to $\mathrm{cm^{-2}}$; $\Sigma=N_\mathrm{col}\,\mu m_u$ and $\Sigma_\mathrm{dust}=0.01\,\Sigma$.
Means are density-weighted, $\langle X\rangle=\sum_k n_kX_k/\sum_k n_k$, for both dust temperatures and
$B_\mathrm{tot}$. The field orientation that a polarization map measures is described by
\begin{equation}
  z = \sum_k n_k\,\sin^{2}\gamma_{B,k}\;e^{2i\psi_{B,k}},\qquad
  (\texttt{cos2psi\_B\_mean},\texttt{sin2psi\_B\_mean})=\frac{(\operatorname{Re}z,\operatorname{Im}z)}{|z|},\qquad
  \texttt{B\_coherence}=\frac{|z|}{\sum_k n_k}\in[0,1],
  \label{eq:coherence}
\end{equation}
so \texttt{B\_coherence} combines field tangling along the line of sight and inclination; it correlates with the
$850\,\mu\mathrm{m}$ polarization fraction with $r=0.87$--$0.94$ (Pearson, per file, in the eight files checked
by \texttt{validate\_dataset.py}).

\paragraph{Bands.} POLARIS returns $I,Q,U$ in Jy per pixel; they are converted to $\mathrm{MJy\,sr^{-1}}$ by
dividing by $10^{6}\,\Omega_\mathrm{pix}$. The polarization fraction is $p=\sqrt{Q^{2}+U^{2}}/I$ (NaN where
$I\le0$) and \texttt{polangle} $=\tfrac12\arctan(U,Q)$ in degrees, mapped to $(-90^\circ,90^\circ]$; it is the
E-vector angle, and the inferred field angle is $\texttt{polangle}+90^\circ$. $I,Q,U$ are stored in double
precision because at $3.4$--$5.8\,\mu\mathrm{m}$ the intensities reach $10^{-108}\,\mathrm{MJy\,sr^{-1}}$, below
the single-precision range. POLARIS was run without stochastic heating of small grains, so the 14 bands at
$\le25\,\mu\mathrm{m}$ are flagged \texttt{reliable: false} in the manifest; the 24 bands at
$53\,\mu\mathrm{m}$--$1.38\,\mathrm{mm}$ are reliable. The images contain no observational noise and no beam.

\begin{table}[ht]
  \centering
  \caption{The 38 detector wavelengths and band names (column prefixes). Bands at $\le25\,\mu\mathrm{m}$ are
  flagged unreliable.}
  \label{tab:bands}
  \footnotesize
  \begin{tabular}{@{}rl rl@{}}
    \toprule
    $\lambda$ [$\mu$m] & Band & $\lambda$ [$\mu$m] & Band \\
    \midrule
    3.4   & \texttt{wise\_1}            & 89    & \texttt{sofia\_hawc\_plus\_c} \\
    3.6   & \texttt{spitzer\_irac\_1}   & 90    & \texttt{akari\_fis\_90} \\
    4.5   & \texttt{spitzer\_irac\_2}   & 100   & \texttt{herschel\_pacs\_2\_or\_iras\_100} \\
    4.6   & \texttt{wise\_2}            & 140   & \texttt{akari\_fis\_140} \\
    5.8   & \texttt{spitzer\_irac\_3}   & 154   & \texttt{sofia\_hawc\_plus\_d} \\
    8.0   & \texttt{spitzer\_irac\_4}   & 160   & \texttt{herschel\_pacs\_3} \\
    8.28  & \texttt{msx\_1}             & 214   & \texttt{sofia\_hawc\_plus\_e} \\
    12.0  & \texttt{wise\_3\_or\_iras\_12} & 250 & \texttt{herschel\_spire\_1} \\
    12.13 & \texttt{msx\_2}             & 350   & \texttt{herschel\_spire\_2\_or\_alma\_band\_10\_c43\_8\_or\_planck\_hfi\_857ghz} \\
    14.65 & \texttt{msx\_3}             & 460   & \texttt{alma\_band\_9\_c43\_8} \\
    21.34 & \texttt{msx\_4}             & 500   & \texttt{herschel\_spire\_3} \\
    22.0  & \texttt{wise\_4}            & 550   & \texttt{planck\_hfi\_545ghz} \\
    24.0  & \texttt{spitzer\_mips\_24}  & 690   & \texttt{alma\_band\_8\_c43\_8} \\
    25.0  & \texttt{iras\_25}           & 849   & \texttt{planck\_hfi\_353ghz} \\
    53    & \texttt{sofia\_hawc\_plus\_a} & 870 & \texttt{apex\_laboca} \\
    60    & \texttt{iras\_60}           & 945   & \texttt{alma\_band\_7\_c43\_8} \\
    63    & \texttt{sofia\_hawc\_plus\_b} & 1100 & \texttt{cso\_bolocam} \\
    65    & \texttt{akari\_fis\_65}     & 1300  & \texttt{alma\_band\_6\_c43\_8} \\
    70    & \texttt{herschel\_pacs\_1}  & 1380  & \texttt{planck\_hfi\_217ghz} \\
    \bottomrule
  \end{tabular}
\end{table}

\subsubsection{Train, validation and test splits}
\label{sec:dataset:splits}

Boxes in the same family image largely the same gas a short time apart, so splitting by box or by pixel would put
near-identical data into training and test sets. Whole families were therefore assigned to one split, and all
three views of a box share its split. Families were shuffled with a fixed seed (NumPy \texttt{default\_rng},
seed 20261111) and assigned in turn to the split furthest below its target share (80/10/10\% of the boxes) among
the splits the family fits into; families with more than 50 boxes always go to the training set, so that no single
region dominates a validation or test set. Two versions of the dataset were built from identical data
(Table~\ref{tab:splits}):

\begin{itemize}
  \item \texttt{starforge\_polaris\_dataset\_v1} (builder 1.1.1): the split described above. All 52 boxes with
  stars belong to families assigned to the training set.
  \item \texttt{starforge\_polaris\_dataset\_v1\_stars\_eval} (builder 1.2.0, option \texttt{--assign}): as v1, but
  the four small star-forming families 41 and 226 are moved to validation and 133 and 139 to test; every other box
  keeps its split. Star-forming families 140 and 280 stay in training because they lie only $3.3$ and
  $2.8\,\mathrm{pc}$ from the large training family 2. The non-star validation and test families are identical in
  both versions, so models trained on either can be compared on that common subset.
\end{itemize}

\begin{table}[ht]
  \centering
  \caption{Splits of the two dataset versions (boxes; each box gives $3\times4096$ rows).}
  \label{tab:splits}
  \small
  \begin{tabular}{@{}lrrrrrr@{}}
    \toprule
     & \multicolumn{3}{c}{\texttt{v1}} & \multicolumn{3}{c}{\texttt{v1\_stars\_eval}} \\
    \cmidrule(lr){2-4}\cmidrule(l){5-7}
    Split & boxes & families & boxes with stars & boxes & families & boxes with stars \\
    \midrule
    train      & 1538 (80.6\%) & 207 & 52 & 1458 (76.4\%) & 203 & 36 \\
    validation &  187 (9.8\%)  &  58 &  0 &  216 (11.3\%) &  60 &  8 \\
    test       &  183 (9.6\%)  &  52 &  0 &  234 (12.3\%) &  54 &  8 \\
    \bottomrule
  \end{tabular}
\end{table}

In v1 the largest family makes up 17\% of the training set, 16\% of the validation set and 21\% of the test
set; all three splits span 0.5--5.0\,Myr. All star-forming families lie within the same central region of the
cloud and appear at late times (snapshots 430--500), with their nearest boxes 2.6--4.8\,pc apart; scores on the
star-forming test families therefore measure performance on new cores of that region rather than of a different
cloud.

% -----------------------------------------------------------------------------------------------------
\subsection{The dataset}
\label{sec:dataset:contents}

\subsubsection{What one row contains}

Each row describes one line of sight through one box, seen in one of the three viewing directions. It holds
\begin{enumerate}
  \item \emph{identifiers} that locate the sightline in the simulation (snapshot, time, core, box, family), in the
  split, and in its image (view, pixel row and column);
  \item the \emph{observables} a telescope would measure along this sightline: Stokes $I$, $Q$, $U$, the
  polarization fraction and the polarization angle in 38 bands;
  \item \emph{sightline targets}: quantities integrated or averaged along the line of sight (column and surface
  density, mean dust temperatures, mean field strength, field orientation and coherence);
  \item \emph{profile targets}: the three-dimensional state of the 64 grid cells along the line of sight, ordered
  from the observer: gas and dust density, the two dust temperatures and the magnetic field (strength, orientation
  and components).
\end{enumerate}
Observables and targets of the same row refer to exactly the same column of grid cells, and the profiles are
already in the observer's frame (Table~\ref{tab:views}), so no geometric transformation is needed when using the
data. Because each row group holds one box and view, the rows of a row group reassemble into a $64\times64$ image
with \texttt{pix\_row} and \texttt{pix\_col}, and the profile columns of that image into a $64\times64\times64$
cube; the dataset therefore serves sightline-based (per-pixel) as well as image-based models.

\subsubsection{Files}

The dataset is stored in Apache Parquet format (written by \texttt{pyarrow}, zstd compression); Table~\ref{tab:files}
lists the files. Each split directory holds one file per snapshot that has boxes in that split, and each file holds
one row group of 4096 rows per box and view, so a single box and view can be read without reading the rest of the
file. The data occupy $178\,\mathrm{GB}$ ($166\,\mathrm{GiB}$) per dataset version. Building them from the
POLARIS output takes 27 minutes with 16 processes.

\begin{table}[ht]
  \centering
  \caption{Files of a dataset version (sizes and counts for \texttt{v1}).}
  \label{tab:files}
  \small
  \begin{tabular}{@{}p{0.27\textwidth}p{0.52\textwidth}r@{}}
    \toprule
    Path & Content & Size \\
    \midrule
    \texttt{data/train/<snapshot>.parquet} & 37 files; 1538 boxes, 4614 row groups, 18\,898\,944 rows & 143.5\,GB \\
    \texttt{data/val/<snapshot>.parquet}   & 38 files; 187 boxes, 561 row groups, 2\,297\,856 rows & 17.4\,GB \\
    \texttt{data/test/<snapshot>.parquet}  & 33 files; 183 boxes, 549 row groups, 2\,248\,704 rows & 17.1\,GB \\
    \texttt{splits.csv} & one line per box: \texttt{box} (\texttt{snapshot\_NNN/core\_NNNNNN}), \texttt{family\_id}, \texttt{split} & \\
    \texttt{manifest.json} & builder version and path, creation time, input paths, numbers of boxes, rows and columns,
      geometry conventions, units, the band table with wavelength and \texttt{reliable} flag, split seed, target fractions,
      sizes (and, in \texttt{v1\_stars\_eval}, the family overrides), the largest column-density mismatch, and a QA record per box
      (column-density check per view, non-positive intensities per band) & 612\,kB \\
    \bottomrule
  \end{tabular}
\end{table}

\subsubsection{Columns}

Table~\ref{tab:dictionary} defines every column. In the column names, \texttt{<b>} stands for one of the 38 band
names of Table~\ref{tab:bands} and \texttt{<k>} for the two-digit cell index \texttt{00}\ldots\texttt{63} along
the line of sight (\texttt{00} nearest the observer). A column with suffix \texttt{\_log} contains $\log_{10}$ of the
column without it. Of the 1443 columns, 1320 are single precision (\texttt{float32}), 114 double precision
(\texttt{float64}: all $I$, $Q$, $U$), 8 integers and one (\texttt{split}) a dictionary-encoded string. The
dataset contains no missing or non-finite values: \texttt{validate\_dataset.py} scanned every value of every
column and found none, so the NaN conventions of the builder (logarithm of a non-positive value, $p$ where
$I\le0$) never apply.

\begingroup
\small
\setlength{\LTcapwidth}{\textwidth}
\begin{longtable}{@{}p{0.27\textwidth}p{0.08\textwidth}p{0.11\textwidth}p{0.46\textwidth}@{}}
  \caption{Data dictionary. Type: f32/f64 = single/double-precision float, i8/i16/i32 = integers. Counts in
  brackets give the number of columns of a pattern.}
  \label{tab:dictionary}\\
  \toprule
  Column & Type & Unit & Definition \\
  \midrule
  \endfirsthead
  \multicolumn{4}{@{}l}{\small\emph{Table~\ref{tab:dictionary} continued}}\\
  \toprule
  Column & Type & Unit & Definition \\
  \midrule
  \endhead
  \midrule
  \multicolumn{4}{r@{}}{\small\emph{continued on next page}}\\
  \endfoot
  \bottomrule
  \endlastfoot
  \multicolumn{4}{@{}l}{\emph{Identifiers} (10 columns)} \\
  \texttt{snapshot} & i16 & -- & STARFORGE snapshot number, 130--500 \\
  \texttt{time\_Myr} & f32 & Myr & simulation time of the snapshot, 0.528--4.957 \\
  \texttt{core\_id} & i32 & -- & ID of the box's central core in the step-1 catalogue of that snapshot \\
  \texttt{box\_id} & i32 & -- & unique box ID, 0--1907 (step 2) \\
  \texttt{family\_id} & i32 & -- & cross-snapshot family, 0--316 (Sect.~\ref{sec:dataset:boxes}) \\
  \texttt{n\_stars} & i16 & -- & number of stellar sources inside the box, 0--39 \\
  \texttt{split} & string & -- & \texttt{train}, \texttt{val} or \texttt{test} \\
  \texttt{view} & i8 & -- & viewing direction 0, 1, 2 (Table~\ref{tab:views}) \\
  \texttt{pix\_row}, \texttt{pix\_col} & i8 & -- & image pixel, 0--63 \\
  \midrule
  \multicolumn{4}{@{}l}{\emph{Sightline quantities} (15 columns)} \\
  \texttt{N\_col} (+\texttt{\_log}) & f32 & cm$^{-2}$ & column density of gas particles from POLARIS's column-density image \\
  \texttt{Sigma} (+\texttt{\_log}) & f32 & g\,cm$^{-2}$ & gas surface density, $N_\mathrm{col}\,\mu m_u$ \\
  \texttt{Sigma\_dust} (+\texttt{\_log}) & f32 & g\,cm$^{-2}$ & dust surface density, $0.01\,\Sigma$ \\
  \texttt{temp\_polaris\_mean} (+\texttt{\_log}) & f32 & K & density-weighted mean POLARIS dust temperature \\
  \texttt{temp\_starforge\_mean} (+\texttt{\_log}) & f32 & K & density-weighted mean STARFORGE dust temperature \\
  \texttt{B\_tot\_mean} (+\texttt{\_log}) & f32 & $\mu$G & density-weighted mean field strength $|\mathbf{B}|$ \\
  \texttt{cos2psi\_B\_mean}, \texttt{sin2psi\_B\_mean} & f32 & -- & direction of $z$ (Eq.~\ref{eq:coherence}): plane-of-sky field orientation seen in polarization \\
  \texttt{B\_coherence} & f32 & -- & $|z|/\sum_k n_k$, 0--1 (Eq.~\ref{eq:coherence}) \\
  \midrule
  \multicolumn{4}{@{}l}{\emph{Bands} (7 columns $\times$ 38 bands = 266)} \\
  \texttt{<b>\_I} & f64 & MJy\,sr$^{-1}$ & Stokes $I$ \\
  \texttt{<b>\_Q}, \texttt{<b>\_U} & f64 & MJy\,sr$^{-1}$ & Stokes $Q$, $U$ in the image frame (angles from $+$col towards $+$row) \\
  \texttt{<b>\_p} & f32 & -- & polarization fraction $\sqrt{Q^2+U^2}/I$ \\
  \texttt{<b>\_polangle} & f32 & deg & polarization (E-vector) angle $\tfrac12\arctan(U,Q)$ in $(-90,90]$; field angle $=$ \texttt{polangle} $+90^\circ$ \\
  \texttt{<b>\_I\_log}, \texttt{<b>\_p\_log} & f32 & -- & $\log_{10}$ of $I$ [MJy\,sr$^{-1}$] and of $p$ \\
  \midrule
  \multicolumn{4}{@{}l}{\emph{Profiles} (64 cells each, $k=00$ nearest the observer; 1152 columns)} \\
  \texttt{dens\_<k>} (+\texttt{\_log}) & f32 & cm$^{-3}$ & gas particle density $n=\rho/(\mu m_u)$, $\mu=2.33$ \\
  \texttt{dens\_dust\_<k>} (+\texttt{\_log}) & f32 & g\,cm$^{-3}$ & dust mass density, $0.01\,\rho$ \\
  \texttt{temp\_polaris\_<k>} (+\texttt{\_log}) & f32 & K & dust temperature computed by POLARIS (used for the images) \\
  \texttt{temp\_starforge\_<k>} (+\texttt{\_log}) & f32 & K & STARFORGE dust temperature (mass-weighted deposit) \\
  \texttt{B\_tot\_<k>} (+\texttt{\_log}) & f32 & $\mu$G & field strength $|\mathbf{B}|$ \\
  \texttt{B\_pos\_<k>} (+\texttt{\_log}) & f32 & $\mu$G & plane-of-sky field strength $\sqrt{B_\mathrm{col}^2+B_\mathrm{row}^2}$ \\
  \texttt{cos2psi\_B\_<k>}, \texttt{sin2psi\_B\_<k>} & f32 & -- & plane-of-sky field orientation $\cos2\psi_B$, $\sin2\psi_B$ \\
  \texttt{sin2gamma\_B\_<k>} & f32 & -- & $\sin^2\gamma_B=B_\mathrm{pos}^2/B_\mathrm{tot}^2$ (1: field in the sky plane; 0: along the line of sight) \\
  \texttt{B\_col\_<k>}, \texttt{B\_row\_<k>} & f32 & $\mu$G & field components along image $+$col and $+$row \\
  \texttt{B\_los\_<k>} & f32 & $\mu$G & line-of-sight component, positive away from the observer \\
\end{longtable}
\endgroup

\subsubsection{Value ranges}

Table~\ref{tab:ranges} gives the distribution of the main quantities. The sightline and band columns were
evaluated over all 23\,445\,504 rows of \texttt{v1}; the profile columns over all 64 cells of 300 randomly chosen
box views ($7.9\times10^{7}$ cells). All column densities in the dataset correspond to optically thin dust at
long wavelengths ($\tau_{850}\le0.01$, Sect.~\ref{sec:dataset:verification}). Intensities at
$\le5.8\,\mu\mathrm{m}$ are of order $10^{-108}\,\mathrm{MJy\,sr^{-1}}$ and carry no usable signal.

\begin{table}[ht]
  \centering
  \caption{Distribution of selected quantities: minimum, 1st percentile, median, 99th percentile and maximum.}
  \label{tab:ranges}
  \small
  \begin{tabular}{@{}lllllll@{}}
    \toprule
    Column & Unit & min & 1\% & median & 99\% & max \\
    \midrule
    \multicolumn{7}{@{}l}{\emph{Sightline (all rows)}} \\
    \texttt{N\_col}               & cm$^{-2}$     & $2.4\times10^{18}$ & $9.9\times10^{19}$ & $9.2\times10^{20}$ & $7.1\times10^{21}$ & $2.5\times10^{23}$ \\
    \texttt{Sigma}                & g\,cm$^{-2}$  & $9.3\times10^{-6}$ & $3.8\times10^{-4}$ & $3.6\times10^{-3}$ & $2.7\times10^{-2}$ & $0.97$ \\
    \texttt{temp\_polaris\_mean}  & K             & 8.27 & 12.5 & 14.5 & 16.0 & 54.4 \\
    \texttt{temp\_starforge\_mean}& K             & 6.16 & 11.3 & 13.3 & 15.5 & 53.8 \\
    \texttt{B\_tot\_mean}         & $\mu$G        & 0.96 & 9.6 & 25.3 & 58.3 & 1760 \\
    \texttt{B\_coherence}         & --            & $3.6\times10^{-5}$ & 0.029 & 0.47 & 0.96 & 0.999 \\
    \midrule
    \multicolumn{7}{@{}l}{\emph{Bands (all rows)}} \\
    $I$ (53\,$\mu$m)   & MJy\,sr$^{-1}$ & $3.2\times10^{-4}$ & 0.0074 & 0.018 & 0.072 & $2.1\times10^{6}$ \\
    $I$ (214\,$\mu$m)  & MJy\,sr$^{-1}$ & 0.083 & 3.3 & 20.5 & 87.4 & $1.2\times10^{5}$ \\
    $I$ (849\,$\mu$m)  & MJy\,sr$^{-1}$ & 0.0039 & 0.17 & 1.33 & 8.1 & 1380 \\
    $I$ (1380\,$\mu$m) & MJy\,sr$^{-1}$ & $8.0\times10^{-4}$ & 0.034 & 0.28 & 1.8 & 246 \\
    $p$ (53\,$\mu$m)   & --  & $1.4\times10^{-5}$ & 0.0059 & 0.098 & 0.259 & 0.355 \\
    $p$ (214\,$\mu$m)  & --  & $1.0\times10^{-5}$ & 0.0052 & 0.089 & 0.260 & 0.323 \\
    $p$ (849\,$\mu$m)  & --  & $2.2\times10^{-6}$ & 0.0051 & 0.089 & 0.261 & 0.324 \\
    $I$ (3.4\,$\mu$m)  & MJy\,sr$^{-1}$ & $7\times10^{-136}$ & $3\times10^{-112}$ & $8\times10^{-108}$ & $1\times10^{-104}$ & $3\times10^{-12}$ \\
    \midrule
    \multicolumn{7}{@{}l}{\emph{Profiles (300 box views, all cells)}} \\
    \texttt{dens}           & cm$^{-3}$    & 0.010 & 5.8 & 73.8 & 2310 & $9.9\times10^{5}$ \\
    \texttt{dens\_dust}     & g\,cm$^{-3}$ & $3.9\times10^{-28}$ & $2.3\times10^{-25}$ & $2.9\times10^{-24}$ & $8.9\times10^{-23}$ & $3.8\times10^{-20}$ \\
    \texttt{temp\_polaris}  & K            & 8.37 & 12.7 & 15.0 & 16.1 & 80.0 \\
    \texttt{temp\_starforge}& K            & 5.91 & 11.3 & 13.4 & 15.6 & 66.8 \\
    \texttt{B\_tot}         & $\mu$G       & 0.037 & 7.3 & 26.0 & 69.1 & 2380 \\
    \texttt{B\_pos}         & $\mu$G       & 0.0017 & 2.7 & 20.2 & 60.2 & 2370 \\
    \texttt{sin2gamma\_B}   & --           & $2.6\times10^{-8}$ & 0.019 & 0.77 & 1.00 & 1.00 \\
    \texttt{B\_los}         & $\mu$G       & $-1600$ & $-45.8$ & $-6.6$ & 31.8 & 1320 \\
    \bottomrule
  \end{tabular}
\end{table}

The intensity spectra of all sightlines peak between 160 and $250\,\mu\mathrm{m}$ (at $214\,\mu\mathrm{m}$ for
96.8\%), as expected for dust at 12--16\,K. The polarization fraction changes little with wavelength between
$53\,\mu\mathrm{m}$ and $1.38\,\mathrm{mm}$ (median 0.098 at $53\,\mu\mathrm{m}$, 0.089 at 214 and
$849\,\mu\mathrm{m}$; maximum 0.355).

\subsubsection{Reading the data}

Columns can be read selectively, which keeps memory small; for example, with \texttt{pandas} and \texttt{pyarrow}:
\begin{verbatim}
import pandas as pd, pyarrow.dataset as ds
root = "/scratch/sjafaritabar/starforge_polaris_dataset_v1/data"
cols = ["planck_hfi_353ghz_I_log", "planck_hfi_353ghz_p"] + \
       [f"dens_{k:02d}_log" for k in range(64)]
train = pd.read_parquet(f"{root}/train", columns=cols)
full = ds.dataset(root, format="parquet", partitioning=None)   # lazy, all splits
\end{verbatim}
To obtain images, a file is read one row group at a time (\texttt{pyarrow.parquet.ParquetFile.read\_row\_group})
and the 4096 rows are placed at (\texttt{pix\_row}, \texttt{pix\_col}).

% -----------------------------------------------------------------------------------------------------
\subsection{Verification}
\label{sec:dataset:verification}

Each step has unit tests (step 2: 7, step 3: 13, step 5: 4; all pass), and an independent audit
(\texttt{audit\_pipeline.py}) re-derives every stage from the raw files without using the pipeline's code. It
passed all 15 checks for both dataset versions:
\begin{enumerate}
  \item Every boxed core is in its step-1 catalogue with the same peak, peak particle and mass; box times equal
  the snapshot header times; no two boxes of a snapshot overlap; families re-derived from scratch reproduce all
  317 \texttt{family\_id}s.
  \item Exactly one grid per box; for every box the summaries, grids, star lists, command files (photon numbers,
  radius, dust, wavelengths, views) and POLARIS outputs are mutually consistent, the grid density is unchanged by
  POLARIS, and dust temperatures and alignment radii are within range.
  \item The $850\,\mu\mathrm{m}$ intensity in $\mathrm{MJy\,sr^{-1}}$ matches the optically thin model
  $\sum_k B_\nu(T_k)\,\Delta\tau_k$ computed from the grid, which tests the units, pixel solid angle and distance:
  POLARIS/model has a per-box median of $1.002$ (range $1.001$--$1.005$); $\tau_{850}\le0.01$.
  \item No box has an outlying POLARIS-vs-STARFORGE dust-temperature offset ($>5\sigma$).
  \item Every one of the 1443 columns was recomputed from the raw grids and FITS files with explicit index formulas
  (Table~\ref{tab:views}) for 1080 random sightlines in 36 boxes ($1.56\times10^{6}$ values); the worst
  relative (or, for logarithms and angles, absolute) error is $3.8\times10^{-6}$, i.e.\ single-precision rounding.
  \item \texttt{splits.csv} covers every box once with the step-2 family IDs; no family is in more than one split;
  all families with more than 50 boxes are in training.
  \item Polarization fractions at $\ge50\,\mu\mathrm{m}$ are physical (maximum 0.355), and the spectra of all
  sightlines peak at 60--$500\,\mu\mathrm{m}$ (214\,$\mu$m for 96.8\%).
\end{enumerate}
\texttt{validate\_dataset.py} additionally confirms the row counts, the absence of duplicate sightlines and of
boxes or families in more than one split, and scans every column of every row for non-finite values.

Across the whole dataset, column densities range over $2.4\times10^{18}$--$2.5\times10^{23}\,\mathrm{cm^{-2}}$
(median $9.1\times10^{20}$), density-weighted POLARIS dust temperatures over 8.3--54.4\,K (median 14.5\,K),
density-weighted field strengths over $0.96$--$1760\,\mu\mathrm{G}$ (median $25.7\,\mu\mathrm{G}$), and
\texttt{B\_coherence} over $4\times10^{-5}$--$0.999$ (median 0.48).

% -----------------------------------------------------------------------------------------------------
\subsection{Limitations}
\label{sec:dataset:limitations}

\begin{itemize}
  \item \textbf{Radiation field.} Each box is illuminated by an unattenuated Mathis field with $G_0=1$. In
  STARFORGE the field reaching the box surfaces is attenuated by the surrounding cloud (snapshot 440: equivalent
  $G_0$ median 0.20, 5--95\% range 0.11--0.46). POLARIS dust temperatures are therefore higher than the STARFORGE
  ones, by a per-box median of $+1.49\,\mathrm{K}$ (5--95\%: $+0.29$ to $+2.76\,\mathrm{K}$). Both temperatures
  are in the dataset.
  \item \textbf{Isolated boxes.} Only gas inside the $2.2\,\mathrm{pc}$ cube emits and absorbs; foreground and
  background cloud emission and extinction are absent.
  \item \textbf{No stochastic heating}: bands at $\le25\,\mu\mathrm{m}$ are not reliable.
  \item \textbf{Ideal observations.} No noise, no beam, a single distance (140\,pc) and a pixel equal to a grid
  cell. Noise and beam have to be added when training.
  \item \textbf{No gas kinematics.} Velocities are not included, so velocity-based field-strength estimates
  (e.g.\ Davis--Chandrasekhar--Fermi) cannot be reproduced from the dataset alone.
  \item \textbf{One simulation.} All boxes come from one cloud; families make the evaluation sets independent in
  gas, not in environment. Late snapshots, which contain all stars, are $0.41\,\mathrm{Myr}$ apart.
  \item \textbf{Monte Carlo.} Dust temperatures carry $\approx0.008\,\mathrm{K}$ of Monte Carlo noise, and runs
  are not bitwise reproducible.
\end{itemize}

% -----------------------------------------------------------------------------------------------------
\subsection{Data products and code}
\label{sec:dataset:files}

Code: \texttt{STARFORGE/creating\_the\_dataset/} (\texttt{1.finding\_cores}, \texttt{2.selecting\_boxes},
\texttt{3.preparing\_polaris\_grids} with its \texttt{diagnostics}, \texttt{4.running\_polaris},
\texttt{5.building\_dataset}, \texttt{audit\_pipeline.py}); photon studies in
\texttt{STARFORGE/M2e5\_R30\_Res585/POLARIS/isrf\_photon\_noise\_test} and \texttt{STARFORGE/M2e5\_R30\_Res585/POLARIS/star\_photon\_test}.
Data: snapshots in \texttt{/scratch/STARFORGE/M2e5\_R30\_Res585}; core catalogues in
\texttt{1.finding\_cores/dendrogram\_cores\_v5.2}; boxes in \texttt{2.selecting\_boxes/boxes\_2.2pc};
grids in \texttt{3.preparing\_polaris\_grids/grids\_64}; POLARIS output in
\texttt{/scratch/sjafaritabar/polaris\_runs} ($\approx70\,\mathrm{GB}$); datasets in
\texttt{/scratch/sjafaritabar/starforge\_polaris\_dataset\_v1} and
\texttt{/scratch/sjafaritabar/starforge\_polaris\_dataset\_v1\_stars\_eval}; audit reports in
\texttt{/scratch/sjafaritabar/audit}.
. Required packages in the preamble:
%   \usepackage{amsmath}  \usepackage{booktabs}  \usepackage{graphicx}  \usepackage{longtable}
% References are in dataset_construction.bib (\bibliography{dataset_construction}).
% Every number below is taken from the pipeline code, its logs, manifests and the audit report
% (paths in Sect. "Data products and code"); dates are UTC.
% =====================================================================================================

\section{A synthetic dust-emission and polarization dataset of star-forming cores}
\label{sec:dataset}

This section describes how the training data for this thesis were constructed. Starting from a
magnetohydrodynamic simulation of a star-forming molecular cloud, dense cores were identified in every
simulation snapshot, a $2.2\,\mathrm{pc}$ cube around each selected core was converted into a radiative-transfer
model, and the dust temperature, grain alignment and polarized dust emission of every cube were computed with
the Monte Carlo code POLARIS. The resulting images were combined with the three-dimensional state of the gas
into a table with one row per line of sight. Table~\ref{tab:pipeline} summarises the five steps.

\begin{table}[ht]
  \centering
  \caption{The five steps of the pipeline. Each step is a separate script; versions refer to the
  version strings recorded in each step's output.}
  \label{tab:pipeline}
  \small
  \begin{tabular}{@{}llll@{}}
    \toprule
    Step & Task & Product & Code (version) \\
    \midrule
    1 & Dense cores in every snapshot   & 11\,599 cores in 51 snapshots & \texttt{find\_starforge\_cores\_dendrogram\_v5.py} (5.2.0) \\
    2 & Non-overlapping boxes, families & 1908 boxes, 317 families     & \texttt{select\_boxes.py} \\
    3 & POLARIS grids and commands      & 1908 grids of $64^3$ cells    & \texttt{prepare\_grids.py} (3.5.1) \\
    4 & Radiative transfer              & 1908 runs, 217\,512 images    & \texttt{run\_polaris.py}, POLARIS \\
    5 & Sightline table                 & 23\,445\,504 rows, 1443 columns & \texttt{build\_dataset.py} (1.1.1 / 1.2.0) \\
    \bottomrule
  \end{tabular}
\end{table}

% -----------------------------------------------------------------------------------------------------
\subsection{The STARFORGE simulation}
\label{sec:dataset:sim}

The gas data come from the STARFORGE \citep{Grudic2021} simulation \texttt{M2e5\_R30\_Res585}, run with the
meshless magnetohydrodynamics code GIZMO \citep{Hopkins2015} (header \texttt{GIZMO\_version} 2022). The initial
condition is a uniform-density spherical cloud of $2\times10^{5}\,M_\odot$ and radius $30\,\mathrm{pc}$ (measured:
$1.993\times10^{5}\,M_\odot$ inside $R=30\,\mathrm{pc}$, median density $1.19\times10^{-22}\,\mathrm{g\,cm^{-3}}$)
embedded in a periodic box of side $300\,\mathrm{code\ units}=300.42\,\mathrm{pc}$. The initial three-dimensional
velocity dispersion of the cloud is $5.86\,\mathrm{km\,s^{-1}}$, which corresponds to a virial parameter
$\alpha=2$; the initial gas temperature is $\approx12\,\mathrm{K}$; gas cells have a median mass of
$1.001\times10^{-3}\,M_\odot$ (range $1.001$--$1.301\times10^{-3}\,M_\odot$). The interstellar radiation field
strength parameter in the snapshot header is $1$; solar abundances are $Z_\odot=0.0142$, $Y_\odot=0.2703$.

\paragraph{Units.} The code units are $1.0014\,\mathrm{pc}$ (length), $1.0008\,M_\odot$ (mass) and
$1\,\mathrm{m\,s^{-1}}$ (velocity), so one code time unit is $979.16\,\mathrm{Myr}$. The units were verified
against the initial conditions (\texttt{check\_snapshot\_units.py}): with these units the gravitational constant
derived from the header units agrees with the header value to $3.4\times10^{-4}$, the enclosed mass and density
match the nominal cloud, and the velocity dispersion matches $\alpha=2$. The saved \texttt{MagneticField} is in
tesla, \emph{not} in the unit given by the header field \texttt{Internal\_UnitB\_In\_Gauss}
($2.91\times10^{-9}\,\mathrm{G}$): read in tesla, the initial cloud field is $7.19\,\mu\mathrm{G}$, consistent
(within a factor $1.26$, the uncertainty of the critical mass-to-flux normalisation) with the nominal
normalised mass-to-flux ratio $4.2$, whereas the gauss or header-unit readings are off by four or more
orders of magnitude. POLARIS expects the field in tesla, so the field is passed on unconverted.

\paragraph{Snapshots.} Snapshots \texttt{snapshot\_000} to \texttt{snapshot\_500} in steps of 10 were analysed
(51 snapshots, each $\approx 31\,\mathrm{GB}$ with $\approx 2.5\times10^{8}$ gas cells). Dense cores exist from
snapshot 130 onwards. Snapshots 130--420 are $0.0406\,\mathrm{Myr}$ apart (0.528--1.707\,Myr); snapshots
420--500 are $0.406\,\mathrm{Myr}$ apart (1.707--4.957\,Myr). Sink particles (stars) exist only in the late
snapshots.

% -----------------------------------------------------------------------------------------------------
\subsection{Step 1: dense cores}
\label{sec:dataset:cores}

\paragraph{Molecular density.} For every gas cell the molecular hydrogen number density is
\begin{equation}
  n_{\mathrm{H_2}} = \frac{\rho\, X_\mathrm{H}\, f_\mathrm{neutral}\, f_\mathrm{mol}}{2\,m_\mathrm{H}},\qquad
  X_\mathrm{H} = 1 - Z - Y,
  \label{eq:nh2}
\end{equation}
where $\rho$ is the cell density, $Z$ and $Y$ are the first two columns of \texttt{Metallicity},
$f_\mathrm{neutral}$ is \texttt{NeutralHydrogenAbundance}, $f_\mathrm{mol}$ is \texttt{MolecularMassFraction},
and $m_\mathrm{H}=1.6735575\times10^{-24}\,\mathrm{g}$. Invalid or out-of-range inputs abort the run instead of
being clipped.

\paragraph{Dendrogram.} Cells with $n_{\mathrm{H_2}} > 10^{4}\,\mathrm{cm^{-3}}$ were selected. Two selected cells
are neighbours if either lies within the kernel support radius $h$ (\texttt{SmoothingLength}) of the other,
using periodic distances; the resulting graph was symmetrised. A dendrogram \citep{Rosolowsky2008} was computed
with \texttt{astrodendro} 0.3.1 on this particle graph with
$\texttt{min\_value}=10^{4}\,\mathrm{cm^{-3}}$, $\texttt{min\_delta}=10^{4}\,\mathrm{cm^{-3}}$ and
$\texttt{min\_npix}=100$ cells. Every leaf of the dendrogram, i.e.\ a local maximum of $n_{\mathrm{H_2}}$
with at least 100 cells and a contrast of at least $10^{4}\,\mathrm{cm^{-3}}$ above the level at which it
merges with a neighbour, is a \emph{core}. A leaf is a density peak; it is not required to be gravitationally
bound or associated with a star.

\paragraph{Core properties.} For each core the catalogue records, among others, the peak cell (position,
$n_{\mathrm{H_2}}$, particle ID), mass, volume and radii, shape, internal velocity dispersion, mean gas and dust
temperature, mass-weighted field, the gas-only virial ratio $2(K+U)/|W|$ (direct pair summation with
$10^{-3}\,\mathrm{pc}$ Plummer softening, up to $2\times10^{4}$ cells), and the sink particles near the core,
including the number of sinks within a $2.2\,\mathrm{pc}$ cube centred on the peak.
Step 1 found 11\,599 cores in total (none before snapshot 130; 1038 in snapshot 500) with a total mass of
$4535\,M_\odot$.

% -----------------------------------------------------------------------------------------------------
\subsection{Step 2: boxes and families}
\label{sec:dataset:boxes}

\paragraph{Box selection.} A \emph{box} is an axis-aligned cube of side $S=2.2\,\mathrm{pc}$ centred on a core's
density peak. Within each snapshot, boxes were selected by greedy non-maximum suppression: cores were sorted by
mass (ties, which are common because core masses are multiples of the cell mass, broken by higher peak density,
then lower \texttt{core\_id}), and a core was kept unless its cube overlapped an already kept cube. Two cubes
overlap if the periodic Chebyshev distance between their centres is $<S$. This gives 1908 boxes in 38 snapshots
(130--500; 1 to 104 boxes per snapshot, median 55). The central cores have masses of $0.10$--$9.09\,M_\odot$
(median $0.41\,M_\odot$) and peak densities $n_{\mathrm{H_2}}=2.0\times10^{4}$--$7.0\times10^{9}\,\mathrm{cm^{-3}}$
(median $3.9\times10^{4}\,\mathrm{cm^{-3}}$). Of all cores, 3604 ($1249\,M_\odot$) lie in no box: they were
suppressed by a more massive neighbour without lying inside its cube.

\paragraph{Families.} Boxes of consecutive snapshots largely image the same gas. Two boxes in consecutive
snapshots (consecutive among snapshots that contain boxes) were therefore linked if their cubes overlap
(periodic Chebyshev distance $<S$), and the connected components of these links, found with a union--find
structure, were labelled \texttt{family\_id}. The 1908 boxes form 317 families; the median family has 2 boxes,
the 90th percentile 9.4, the 99th percentile 107, and the largest family has 260 boxes. Six families have more
than 50 boxes. Families are the unit of the train/validation/test split (Sect.~\ref{sec:dataset:splits}).

% -----------------------------------------------------------------------------------------------------
\subsection{Step 3: radiative-transfer models}
\label{sec:dataset:grids}

\paragraph{Grid geometry.} Each box was converted into a POLARIS octree grid that is a uniform cube of
$64^{3}$ cells (octree level 6) of size $\Delta=2.2/64=0.034375\,\mathrm{pc}$. The grid is shifted by half a
cell so that the core's peak cell lies at the centre of cell $(32,32,32)$; centring the cube exactly on the peak
would place the peak on a corner shared by eight cells. Each snapshot was read once and the gas of all its boxes
was cut out in that single pass; boxes were then built in parallel.

\paragraph{Deposition.} Every gas cell whose kernel overlaps the cube (centre within $S/2+h$ of the box centre
along each axis) was deposited with GIZMO's compact cubic-spline kernel,
\begin{equation}
  W(q) \propto \begin{cases} 1-6q^{2}+6q^{3}, & 0\le q<\tfrac12,\\ 2(1-q)^{3}, & \tfrac12\le q<1,\\ 0, & q\ge 1,\end{cases}
  \qquad q = \frac{|\mathbf{x}_c-\mathbf{x}_i|}{h_i},
\end{equation}
sampled at the centres $\mathbf{x}_c$ of all grid cells within the kernel, \emph{including cells outside the cube}.
The weights $w_{ic}=W(q_{ic})/\sum_{c'}W(q_{ic'})$ are normalised over the full kernel and only the in-cube part
is kept, so each grid holds exactly the mass inside the cube and the faces are not under-dense. If a kernel
contains no cell centre ($h_i$ well below $\Delta$), the cell's mass goes to the grid cell containing its centre
(this applied to 1.07\% of the $2.73\times10^{8}$ deposited gas cells). Cell quantities are
\begin{equation}
  \rho_c = \frac{\sum_i m_i w_{ic}}{\Delta^{3}},\qquad
  \langle X\rangle_c = \frac{\sum_i m_i w_{ic} X_i}{\sum_i m_i w_{ic}},
\end{equation}
with $X$ the dust temperature, gas temperature and the three field components (mass-weighted means). A cell
receiving no mass would get a density floor of $10^{-6}$ times the smallest occupied density and the fields of
the nearest occupied cell; this happened for a single cell in a single box ($3.6\times10^{-14}\,M_\odot$).
The grid mass agrees with the mass of the gas cells centred inside the cube to between $-0.12\%$ and $+0.08\%$.
Boxes contain $3.5\times10^{4}$--$7.9\times10^{5}$ contributing gas cells (median $1.3\times10^{5}$) and
$27.6$--$748.6\,M_\odot$ of gas (median $113\,M_\odot$).

\paragraph{Grid fields.} The grid stores (POLARIS field IDs in brackets) the gas number density
$n=\rho/(\mu m_u)$ in $\mathrm{m^{-3}}$ [0], the STARFORGE dust temperature [2], the gas temperature [3] and
$B_x, B_y, B_z$ in tesla [4--6]. The POLARIS branch used reads the gas density only as a number density and
converts it back with $\rho=n\,\mu\,m_u$, where $m_u=1.66053904\times10^{-27}\,\mathrm{kg}$ (POLARIS's
\texttt{m\_H}); the same $n$ and $\mu$ enter the radiative-torque calculation (gas damping, thermal speed).
$\mu$ is therefore the mean mass per gas particle. We use $\mu=2.33$: the mass-weighted median of the particle
mean molecular weight of box gas, computed from the snapshot chemistry, is 2.32--2.35 for
$n_\mathrm{H}=10^{3}$--$10^{6}\,\mathrm{cm^{-3}}$ (diagnostic \texttt{check\_mu.py}, snapshot 500), and
$\rho$, hence the dust mass, is reproduced exactly for any $\mu$.

\paragraph{Dust.} One dust component was used: astrodust \citep{Hensley2023} without iron inclusions
($f_\mathrm{Fe}=0$), with the catalogue's magnetic susceptibility $\chi=4.47\times10^{-4}$ and material density
$2898\,\mathrm{kg\,m^{-3}}$, a power-law size distribution $\mathrm{d}n/\mathrm{d}a\propto a^{-3.5}$
\citep{Mathis1977} between $a_\mathrm{min}=5\,\mathrm{nm}$ and $a_\mathrm{max}=250\,\mathrm{nm}$, and a
dust-to-gas mass ratio of $0.01$ (POLARIS \texttt{<mass\_fraction>}). The ratio is the one STARFORGE itself
uses: its field \texttt{DustToGasRatio\_Local} stores $Z$ times the local dust-to-metals ratio relative to solar,
with a median of $0.01416=Z_\odot$ over all boxes, which corresponds to an absolute dust-to-gas ratio
$0.01\,Z/Z_\odot = 0.01$. The catalogue covers $1\,\mathrm{nm}$--$1\,\mu\mathrm{m}$ in size and
$0.05\,\mu\mathrm{m}$--$5\,\mathrm{mm}$ in wavelength, i.e.\ POLARIS's internal wavelength grid
(100 points, $0.0912$--$2000\,\mu\mathrm{m}$) and all detector wavelengths. With this susceptibility the
Larmor-precession limit of POLARIS's alignment model is never reached: the smallest Larmor-limited grain size
over all $5.0\times10^{8}$ grid cells is $408\,\mathrm{nm}>a_\mathrm{max}$ (\texttt{check\_larmor\_limit.py}).

\paragraph{Radiation sources.} Every box is illuminated by the interstellar radiation field (ISRF) of
\citet{Mathis1983} with strength $G_0=1$, launched from a sphere of 10 times the circumscribed radius of the grid
($10\times\tfrac{\sqrt3}{2}\times2.2\,\mathrm{pc}=19.1\,\mathrm{pc}$). In addition, every sink particle inside the
cube with positive radius and luminosity was added as a blackbody point source at its position, with the
STARFORGE protostellar radius $R_\ast$ and an effective temperature from
$L_\ast = 4\pi R_\ast^{2}\sigma_\mathrm{SB}T_\mathrm{eff}^{4}$. Fifty-two boxes contain stars, 275 sources in total
(up to 39 in one box): $L_\ast=6.8\times10^{-3}$--$1.0\times10^{4}\,L_\odot$ (median $0.16\,L_\odot$),
$R_\ast=0.31$--$20.8\,R_\odot$, $T_\mathrm{eff}=2990$--$27\,900\,\mathrm{K}$, sink masses
$0.02$--$11.7\,M_\odot$. (Step 1 counted 273 sinks in these boxes; in one box, 1851, the half-cell shift of the
grid moves two further sinks inside the cube.) Boxes without stars are heated by the ISRF only.

\paragraph{Photon numbers and Monte Carlo noise.} With radiative-torque alignment, POLARIS launches the
requested number of photons for each of its 100 wavelengths. Production used $10^{7}$ ISRF photons and $10^{6}$
photons per star (per wavelength). These values were chosen from two convergence studies, each running every
configuration twice with independent random numbers and taking the per-cell noise as
$\mathrm{median}\,|x_1-x_2|/\sqrt2$:
(i) ISRF photons $10^{6}$--$10^{7}$ and launch radii 3--30 on box \texttt{snapshot\_440/core\_000395}; at $10^{7}$
photons and radius 10 the dust-temperature noise is $0.0078\,\mathrm{K}$ in the interior, $0.0087\,\mathrm{K}$ on
faces, $0.0103\,\mathrm{K}$ on edges and $0.0137\,\mathrm{K}$ in corners ($0.0081\,\mathrm{K}$ overall), the
alignment radius varies by $0.63\%$, and at $850\,\mu\mathrm{m}$ the noise is $0.026\%$ in $I$, $0.20\%$ in $p$
and $0.035^\circ$ in the polarization angle. Smaller radii concentrate the noise in the corners
($0.073\,\mathrm{K}$ at radius 3), larger radii raise it in the interior ($0.0105\,\mathrm{K}$ at radius 30).
(ii) Star photons $10^{5}$ vs.\ $10^{6}$ on box \texttt{snapshot\_500/core\_000873}: $10^{6}$ halves the
temperature noise in gas below $25\,\mathrm{K}$ (from 0.077--0.084 to 0.039--0.045\,K) and changes the
$53\,\mu\mathrm{m}$ intensity by a median of $2.0\%$ ($0.22\%$ at $850\,\mu\mathrm{m}$).

\paragraph{POLARIS tasks.} Each box has a command file with two tasks, both with radiative-torque alignment
(\texttt{ALIG\_RAT}) and 8 threads:
\begin{enumerate}
  \item \texttt{CMD\_TEMP\_RAT}: Monte Carlo radiative transfer of the ISRF and stellar photons; computes the
  equilibrium dust temperature of every cell and the radiative-torque alignment radius $a_\mathrm{alig}$ (grains
  larger than $a_\mathrm{alig}$ are aligned with the magnetic field). The STARFORGE dust temperature in the input
  grid is replaced by this temperature. Writes \texttt{temp/grid\_temp\_rat.dat}.
  \item \texttt{CMD\_DUST\_EMISSION} on that grid: ray-traced thermal emission of the aligned grains, giving the
  Stokes $I,Q,U,V$, the optical depth and the column density for 38 wavelengths and 3 viewing directions
  (114 detectors) with $64\times64$ pixels each (\texttt{max\_subpixel\_lvl} 0). The detector covers the grid, so
  one pixel corresponds to one grid column ($0.034375\,\mathrm{pc}$); the source distance is
  $140\,\mathrm{pc}$, so a pixel subtends $50.6''$ and
  $\Omega_\mathrm{pix}=(\Delta/d)^{2}=6.03\times10^{-8}\,\mathrm{sr}$.
\end{enumerate}
The 38 wavelengths (Table~\ref{tab:bands}) are those of common far-infrared to millimetre instruments and are
read from one table (\texttt{wavelength\_telescope\_mapping.csv}) by both the grid preparation and the dataset
builder. The three viewing directions are the POLARIS detector rotations
$(\vartheta_1,\vartheta_2)=(0^\circ,0^\circ)$, $(-90^\circ,0^\circ)$ and $(0^\circ,90^\circ)$, which look along the
grid $z$, $y$ and $x$ axes (Sect.~\ref{sec:dataset:geometry}). POLARIS was the \texttt{experimental} branch,
commit \texttt{6c5e1a89c6e14ce806c47bbb87947b98e3f761b0} \citep{Reissl2016}.

% -----------------------------------------------------------------------------------------------------
\subsection{Step 4: radiative-transfer runs}
\label{sec:dataset:runs}

All 1908 boxes were run between 2026-10-31 22:47 and 2026-11-09 12:41 on one machine (AMD EPYC 9554P, 64 cores
/ 128 hardware threads, $377\,\mathrm{GB}$ memory), each box in its own directory, at the lowest scheduling
priority (\texttt{nice 19}), with up to 16 boxes at a time (8 threads each). The runner measured the CPU use of
other users every $15\,\mathrm{s}$ and kept 8 hardware threads free above their peak use over the previous
10 minutes; when the machine was needed, the most recently started boxes were suspended (\texttt{SIGSTOP}) and
resumed 10 minutes after the demand had gone, so no run was ever killed or restarted. Boxes containing stars
were run first. A box took $0.92$--$5.98\,\mathrm{h}$ of wall-clock time (median $1.46\,\mathrm{h}$ at up to 14
concurrent boxes), $2810$ box-hours in total, of which boxes spent $5.3\,\mathrm{h}$ suspended.

Every run was validated before being accepted: the temperature grid exists and all dust temperatures are finite
and within $2$--$3000\,\mathrm{K}$; all 114 detector files exist; $I$, $Q$ and $U$ are finite everywhere; and at
wavelengths $\geq50\,\mu\mathrm{m}$, $I>0$ in every pixel and $Q,U$ are not identically zero. All 1908 runs
finished with exit code 0 and passed; none failed. POLARIS's Monte Carlo is seeded by thread, and the assignment
of photons to threads varies between runs, so runs are statistically but not bitwise reproducible.

% -----------------------------------------------------------------------------------------------------
\subsection{Step 5: the sightline dataset}
\label{sec:dataset:table}

\subsubsection{Geometry and orientation}
\label{sec:dataset:geometry}

Grid arrays are indexed $[i_x,i_y,i_z]$ and images $[\mathrm{row},\mathrm{col}]$, both from 0 to 63. The mapping
between them was determined empirically (\texttt{check\_view\_geometry.py}): for each view, the POLARIS
column-density image was compared with $\sum n\,\Delta$ along each grid axis in all 48 combinations of
line-of-sight axis, transposition and flips; the best orientation matches to a relative difference of
$\approx10^{-13}$. Which end of the line of sight faces the observer cannot be seen in optically thin emission; it
was determined where the dust is optically thick ($21\,\mu\mathrm{m}$, pixels with $\tau>1$), where
$I=\sum_k B_\nu(T_k)\,\Delta\tau_k\,e^{-\tau_k}$ depends on the side from which the absorption is counted
(\texttt{check\_observer\_side.py}). The model built from POLARIS's own temperatures reproduces the image with a
near-constant ratio only if the high-index end of the line-of-sight axis faces the observer, which holds in every
view. Table~\ref{tab:views} gives the result.

\begin{table}[ht]
  \centering
  \caption{Mapping between a dataset row (view $v$, pixel row $r$, pixel column $c$, profile index $k$, with $k=0$
  the cell nearest the observer) and the grid, and the field components in the image frame.}
  \label{tab:views}
  \small
  \begin{tabular}{@{}cccccc@{}}
    \toprule
    View & POLARIS rotation & Line of sight & Image $+$col, $+$row & Grid cell $(i_x,i_y,i_z)$ & $(B_\mathrm{col},B_\mathrm{row},B_\mathrm{los})$ \\
    \midrule
    0 & $(0^\circ,0^\circ)$   & $z$ & $+x$, $+y$ & $(c,\ r,\ 63-k)$    & $(B_x,\ B_y,\ -B_z)$ \\
    1 & $(-90^\circ,0^\circ)$ & $y$ & $+x$, $-z$ & $(c,\ 63-k,\ 63-r)$ & $(B_x,\ -B_z,\ -B_y)$ \\
    2 & $(0^\circ,90^\circ)$  & $x$ & $-z$, $+y$ & $(63-k,\ r,\ 63-c)$ & $(-B_z,\ B_y,\ -B_x)$ \\
    \bottomrule
  \end{tabular}
\end{table}

Angles in the image plane are measured from the image $+$col axis towards $+$row. With Stokes $U$ as written by
POLARIS, the plane-of-sky field direction is the polarization angle plus $90^\circ$: comparing
$\tfrac12\arctan(U,Q)+90^\circ$ with a density-weighted Stokes model of the grid field in this frame (in the more
polarized half of the pixels) gives median differences of $0.2^\circ$--$0.8^\circ$, against $22^\circ$--$54^\circ$ if the
sign of $U$ is reversed. $B_\mathrm{los}>0$ points away from the
observer. Every box and view was checked during the build: $\sum_k n_k\Delta$ along the rotated profiles
reproduces POLARIS's column-density image to a maximum relative difference of $5.4\times10^{-13}$
(the build aborts above $10^{-5}$).

\subsubsection{Contents}

One row is one line of sight: one pixel of one view of one box. The dataset has
$1908\times3\times64\times64=23\,445\,504$ rows and 1443 columns (Table~\ref{tab:columns}).

\begin{table}[ht]
  \centering
  \caption{Columns of the dataset. $k=00\ldots63$ runs along the line of sight from the observer; columns
  marked (+log) also exist as \texttt{<name>\_log} $=\log_{10}$ (NaN where the value is $\le0$). Values are
  \texttt{float32} except $I,Q,U$ (\texttt{float64}) and the integer metadata.}
  \label{tab:columns}
  \small
  \begin{tabular}{@{}p{0.17\textwidth}p{0.47\textwidth}p{0.12\textwidth}r@{}}
    \toprule
    Group & Columns & Unit & Number \\
    \midrule
    metadata & \texttt{snapshot, time\_Myr, core\_id, box\_id, family\_id, n\_stars, split, view, pix\_row, pix\_col} & & 10 \\
    sightline & \texttt{N\_col, Sigma, Sigma\_dust} (+log) & $\mathrm{cm^{-2}}$, $\mathrm{g\,cm^{-2}}$ & 6 \\
              & \texttt{temp\_polaris\_mean, temp\_starforge\_mean} (+log) & K & 4 \\
              & \texttt{B\_tot\_mean} (+log) & $\mu$G & 2 \\
              & \texttt{cos2psi\_B\_mean, sin2psi\_B\_mean, B\_coherence} & -- & 3 \\
    bands ($\times38$) & \texttt{<band>\_I, \_Q, \_U} & MJy\,sr$^{-1}$ & 114 \\
              & \texttt{<band>\_p, \_polangle, \_I\_log, \_p\_log} & --, deg & 152 \\
    profiles ($\times64$) & \texttt{dens\_k, dens\_dust\_k} (+log) & $\mathrm{cm^{-3}}$, $\mathrm{g\,cm^{-3}}$ & 256 \\
              & \texttt{temp\_polaris\_k, temp\_starforge\_k} (+log) & K & 256 \\
              & \texttt{B\_tot\_k, B\_pos\_k} (+log) & $\mu$G & 256 \\
              & \texttt{cos2psi\_B\_k, sin2psi\_B\_k, sin2gamma\_B\_k} & -- & 192 \\
              & \texttt{B\_col\_k, B\_row\_k, B\_los\_k} & $\mu$G & 192 \\
    \midrule
    total & & & 1443 \\
    \bottomrule
  \end{tabular}
\end{table}

\paragraph{Profiles.} For each cell $k$ along the line of sight:
$\texttt{dens}=n$ in $\mathrm{cm^{-3}}$ (gas particles, $\mu=2.33$);
$\texttt{dens\_dust}=0.01\,n\,\mu\,m_u$ in $\mathrm{g\,cm^{-3}}$;
$\texttt{temp\_polaris}$, the dust temperature computed by POLARIS; $\texttt{temp\_starforge}$, the deposited
STARFORGE dust temperature; the field in $\mu$G ($1\,\mathrm{T}=10^{10}\,\mu\mathrm{G}$) in the image frame
(Table~\ref{tab:views}) and the derived quantities
\begin{gather}
  B_\mathrm{pos}^{2}=B_\mathrm{col}^{2}+B_\mathrm{row}^{2},\qquad B_\mathrm{tot}^{2}=B_\mathrm{pos}^{2}+B_\mathrm{los}^{2},\\
  \cos2\psi_B=\frac{B_\mathrm{col}^{2}-B_\mathrm{row}^{2}}{B_\mathrm{pos}^{2}},\qquad
  \sin2\psi_B=\frac{2B_\mathrm{col}B_\mathrm{row}}{B_\mathrm{pos}^{2}},\qquad
  \sin^{2}\gamma_B=\frac{B_\mathrm{pos}^{2}}{B_\mathrm{tot}^{2}},
\end{gather}
where $\psi_B$ is the plane-of-sky field angle and $\gamma_B$ the angle between the field and the line of sight.
The double-angle form removes the $180^\circ$ ambiguity of the field direction, which dust polarization cannot
measure.

\paragraph{Sightline quantities.} $N_\mathrm{col}$ is POLARIS's column-density image converted from
$\mathrm{m^{-2}}$ to $\mathrm{cm^{-2}}$; $\Sigma=N_\mathrm{col}\,\mu m_u$ and $\Sigma_\mathrm{dust}=0.01\,\Sigma$.
Means are density-weighted, $\langle X\rangle=\sum_k n_kX_k/\sum_k n_k$, for both dust temperatures and
$B_\mathrm{tot}$. The field orientation that a polarization map measures is described by
\begin{equation}
  z = \sum_k n_k\,\sin^{2}\gamma_{B,k}\;e^{2i\psi_{B,k}},\qquad
  (\texttt{cos2psi\_B\_mean},\texttt{sin2psi\_B\_mean})=\frac{(\operatorname{Re}z,\operatorname{Im}z)}{|z|},\qquad
  \texttt{B\_coherence}=\frac{|z|}{\sum_k n_k}\in[0,1],
  \label{eq:coherence}
\end{equation}
so \texttt{B\_coherence} combines field tangling along the line of sight and inclination; it correlates with the
$850\,\mu\mathrm{m}$ polarization fraction with $r=0.87$--$0.94$ (Pearson, per file, in the eight files checked
by \texttt{validate\_dataset.py}).

\paragraph{Bands.} POLARIS returns $I,Q,U$ in Jy per pixel; they are converted to $\mathrm{MJy\,sr^{-1}}$ by
dividing by $10^{6}\,\Omega_\mathrm{pix}$. The polarization fraction is $p=\sqrt{Q^{2}+U^{2}}/I$ (NaN where
$I\le0$) and \texttt{polangle} $=\tfrac12\arctan(U,Q)$ in degrees, mapped to $(-90^\circ,90^\circ]$; it is the
E-vector angle, and the inferred field angle is $\texttt{polangle}+90^\circ$. $I,Q,U$ are stored in double
precision because at $3.4$--$5.8\,\mu\mathrm{m}$ the intensities reach $10^{-108}\,\mathrm{MJy\,sr^{-1}}$, below
the single-precision range. POLARIS was run without stochastic heating of small grains, so the 14 bands at
$\le25\,\mu\mathrm{m}$ are flagged \texttt{reliable: false} in the manifest; the 24 bands at
$53\,\mu\mathrm{m}$--$1.38\,\mathrm{mm}$ are reliable. The images contain no observational noise and no beam.

\begin{table}[ht]
  \centering
  \caption{The 38 detector wavelengths and band names (column prefixes). Bands at $\le25\,\mu\mathrm{m}$ are
  flagged unreliable.}
  \label{tab:bands}
  \footnotesize
  \begin{tabular}{@{}rl rl@{}}
    \toprule
    $\lambda$ [$\mu$m] & Band & $\lambda$ [$\mu$m] & Band \\
    \midrule
    3.4   & \texttt{wise\_1}            & 89    & \texttt{sofia\_hawc\_plus\_c} \\
    3.6   & \texttt{spitzer\_irac\_1}   & 90    & \texttt{akari\_fis\_90} \\
    4.5   & \texttt{spitzer\_irac\_2}   & 100   & \texttt{herschel\_pacs\_2\_or\_iras\_100} \\
    4.6   & \texttt{wise\_2}            & 140   & \texttt{akari\_fis\_140} \\
    5.8   & \texttt{spitzer\_irac\_3}   & 154   & \texttt{sofia\_hawc\_plus\_d} \\
    8.0   & \texttt{spitzer\_irac\_4}   & 160   & \texttt{herschel\_pacs\_3} \\
    8.28  & \texttt{msx\_1}             & 214   & \texttt{sofia\_hawc\_plus\_e} \\
    12.0  & \texttt{wise\_3\_or\_iras\_12} & 250 & \texttt{herschel\_spire\_1} \\
    12.13 & \texttt{msx\_2}             & 350   & \texttt{herschel\_spire\_2\_or\_alma\_band\_10\_c43\_8\_or\_planck\_hfi\_857ghz} \\
    14.65 & \texttt{msx\_3}             & 460   & \texttt{alma\_band\_9\_c43\_8} \\
    21.34 & \texttt{msx\_4}             & 500   & \texttt{herschel\_spire\_3} \\
    22.0  & \texttt{wise\_4}            & 550   & \texttt{planck\_hfi\_545ghz} \\
    24.0  & \texttt{spitzer\_mips\_24}  & 690   & \texttt{alma\_band\_8\_c43\_8} \\
    25.0  & \texttt{iras\_25}           & 849   & \texttt{planck\_hfi\_353ghz} \\
    53    & \texttt{sofia\_hawc\_plus\_a} & 870 & \texttt{apex\_laboca} \\
    60    & \texttt{iras\_60}           & 945   & \texttt{alma\_band\_7\_c43\_8} \\
    63    & \texttt{sofia\_hawc\_plus\_b} & 1100 & \texttt{cso\_bolocam} \\
    65    & \texttt{akari\_fis\_65}     & 1300  & \texttt{alma\_band\_6\_c43\_8} \\
    70    & \texttt{herschel\_pacs\_1}  & 1380  & \texttt{planck\_hfi\_217ghz} \\
    \bottomrule
  \end{tabular}
\end{table}

\subsubsection{Train, validation and test splits}
\label{sec:dataset:splits}

Boxes in the same family image largely the same gas a short time apart, so splitting by box or by pixel would put
near-identical data into training and test sets. Whole families were therefore assigned to one split, and all
three views of a box share its split. Families were shuffled with a fixed seed (NumPy \texttt{default\_rng},
seed 20261111) and assigned in turn to the split furthest below its target share (80/10/10\% of the boxes) among
the splits the family fits into; families with more than 50 boxes always go to the training set, so that no single
region dominates a validation or test set. Two versions of the dataset were built from identical data
(Table~\ref{tab:splits}):

\begin{itemize}
  \item \texttt{starforge\_polaris\_dataset\_v1} (builder 1.1.1): the split described above. All 52 boxes with
  stars belong to families assigned to the training set.
  \item \texttt{starforge\_polaris\_dataset\_v1\_stars\_eval} (builder 1.2.0, option \texttt{--assign}): as v1, but
  the four small star-forming families 41 and 226 are moved to validation and 133 and 139 to test; every other box
  keeps its split. Star-forming families 140 and 280 stay in training because they lie only $3.3$ and
  $2.8\,\mathrm{pc}$ from the large training family 2. The non-star validation and test families are identical in
  both versions, so models trained on either can be compared on that common subset.
\end{itemize}

\begin{table}[ht]
  \centering
  \caption{Splits of the two dataset versions (boxes; each box gives $3\times4096$ rows).}
  \label{tab:splits}
  \small
  \begin{tabular}{@{}lrrrrrr@{}}
    \toprule
     & \multicolumn{3}{c}{\texttt{v1}} & \multicolumn{3}{c}{\texttt{v1\_stars\_eval}} \\
    \cmidrule(lr){2-4}\cmidrule(l){5-7}
    Split & boxes & families & boxes with stars & boxes & families & boxes with stars \\
    \midrule
    train      & 1538 (80.6\%) & 207 & 52 & 1458 (76.4\%) & 203 & 36 \\
    validation &  187 (9.8\%)  &  58 &  0 &  216 (11.3\%) &  60 &  8 \\
    test       &  183 (9.6\%)  &  52 &  0 &  234 (12.3\%) &  54 &  8 \\
    \bottomrule
  \end{tabular}
\end{table}

In v1 the largest family makes up 17\% of the training set, 16\% of the validation set and 21\% of the test
set; all three splits span 0.5--5.0\,Myr. All star-forming families lie within the same central region of the
cloud and appear at late times (snapshots 430--500), with their nearest boxes 2.6--4.8\,pc apart; scores on the
star-forming test families therefore measure performance on new cores of that region rather than of a different
cloud.

% -----------------------------------------------------------------------------------------------------
\subsection{The dataset}
\label{sec:dataset:contents}

\subsubsection{What one row contains}

Each row describes one line of sight through one box, seen in one of the three viewing directions. It holds
\begin{enumerate}
  \item \emph{identifiers} that locate the sightline in the simulation (snapshot, time, core, box, family), in the
  split, and in its image (view, pixel row and column);
  \item the \emph{observables} a telescope would measure along this sightline: Stokes $I$, $Q$, $U$, the
  polarization fraction and the polarization angle in 38 bands;
  \item \emph{sightline targets}: quantities integrated or averaged along the line of sight (column and surface
  density, mean dust temperatures, mean field strength, field orientation and coherence);
  \item \emph{profile targets}: the three-dimensional state of the 64 grid cells along the line of sight, ordered
  from the observer: gas and dust density, the two dust temperatures and the magnetic field (strength, orientation
  and components).
\end{enumerate}
Observables and targets of the same row refer to exactly the same column of grid cells, and the profiles are
already in the observer's frame (Table~\ref{tab:views}), so no geometric transformation is needed when using the
data. Because each row group holds one box and view, the rows of a row group reassemble into a $64\times64$ image
with \texttt{pix\_row} and \texttt{pix\_col}, and the profile columns of that image into a $64\times64\times64$
cube; the dataset therefore serves sightline-based (per-pixel) as well as image-based models.

\subsubsection{Files}

The dataset is stored in Apache Parquet format (written by \texttt{pyarrow}, zstd compression); Table~\ref{tab:files}
lists the files. Each split directory holds one file per snapshot that has boxes in that split, and each file holds
one row group of 4096 rows per box and view, so a single box and view can be read without reading the rest of the
file. The data occupy $178\,\mathrm{GB}$ ($166\,\mathrm{GiB}$) per dataset version. Building them from the
POLARIS output takes 27 minutes with 16 processes.

\begin{table}[ht]
  \centering
  \caption{Files of a dataset version (sizes and counts for \texttt{v1}).}
  \label{tab:files}
  \small
  \begin{tabular}{@{}p{0.27\textwidth}p{0.52\textwidth}r@{}}
    \toprule
    Path & Content & Size \\
    \midrule
    \texttt{data/train/<snapshot>.parquet} & 37 files; 1538 boxes, 4614 row groups, 18\,898\,944 rows & 143.5\,GB \\
    \texttt{data/val/<snapshot>.parquet}   & 38 files; 187 boxes, 561 row groups, 2\,297\,856 rows & 17.4\,GB \\
    \texttt{data/test/<snapshot>.parquet}  & 33 files; 183 boxes, 549 row groups, 2\,248\,704 rows & 17.1\,GB \\
    \texttt{splits.csv} & one line per box: \texttt{box} (\texttt{snapshot\_NNN/core\_NNNNNN}), \texttt{family\_id}, \texttt{split} & \\
    \texttt{manifest.json} & builder version and path, creation time, input paths, numbers of boxes, rows and columns,
      geometry conventions, units, the band table with wavelength and \texttt{reliable} flag, split seed, target fractions,
      sizes (and, in \texttt{v1\_stars\_eval}, the family overrides), the largest column-density mismatch, and a QA record per box
      (column-density check per view, non-positive intensities per band) & 612\,kB \\
    \bottomrule
  \end{tabular}
\end{table}

\subsubsection{Columns}

Table~\ref{tab:dictionary} defines every column. In the column names, \texttt{<b>} stands for one of the 38 band
names of Table~\ref{tab:bands} and \texttt{<k>} for the two-digit cell index \texttt{00}\ldots\texttt{63} along
the line of sight (\texttt{00} nearest the observer). A column with suffix \texttt{\_log} contains $\log_{10}$ of the
column without it. Of the 1443 columns, 1320 are single precision (\texttt{float32}), 114 double precision
(\texttt{float64}: all $I$, $Q$, $U$), 8 integers and one (\texttt{split}) a dictionary-encoded string. The
dataset contains no missing or non-finite values: \texttt{validate\_dataset.py} scanned every value of every
column and found none, so the NaN conventions of the builder (logarithm of a non-positive value, $p$ where
$I\le0$) never apply.

\begingroup
\small
\setlength{\LTcapwidth}{\textwidth}
\begin{longtable}{@{}p{0.27\textwidth}p{0.08\textwidth}p{0.11\textwidth}p{0.46\textwidth}@{}}
  \caption{Data dictionary. Type: f32/f64 = single/double-precision float, i8/i16/i32 = integers. Counts in
  brackets give the number of columns of a pattern.}
  \label{tab:dictionary}\\
  \toprule
  Column & Type & Unit & Definition \\
  \midrule
  \endfirsthead
  \multicolumn{4}{@{}l}{\small\emph{Table~\ref{tab:dictionary} continued}}\\
  \toprule
  Column & Type & Unit & Definition \\
  \midrule
  \endhead
  \midrule
  \multicolumn{4}{r@{}}{\small\emph{continued on next page}}\\
  \endfoot
  \bottomrule
  \endlastfoot
  \multicolumn{4}{@{}l}{\emph{Identifiers} (10 columns)} \\
  \texttt{snapshot} & i16 & -- & STARFORGE snapshot number, 130--500 \\
  \texttt{time\_Myr} & f32 & Myr & simulation time of the snapshot, 0.528--4.957 \\
  \texttt{core\_id} & i32 & -- & ID of the box's central core in the step-1 catalogue of that snapshot \\
  \texttt{box\_id} & i32 & -- & unique box ID, 0--1907 (step 2) \\
  \texttt{family\_id} & i32 & -- & cross-snapshot family, 0--316 (Sect.~\ref{sec:dataset:boxes}) \\
  \texttt{n\_stars} & i16 & -- & number of stellar sources inside the box, 0--39 \\
  \texttt{split} & string & -- & \texttt{train}, \texttt{val} or \texttt{test} \\
  \texttt{view} & i8 & -- & viewing direction 0, 1, 2 (Table~\ref{tab:views}) \\
  \texttt{pix\_row}, \texttt{pix\_col} & i8 & -- & image pixel, 0--63 \\
  \midrule
  \multicolumn{4}{@{}l}{\emph{Sightline quantities} (15 columns)} \\
  \texttt{N\_col} (+\texttt{\_log}) & f32 & cm$^{-2}$ & column density of gas particles from POLARIS's column-density image \\
  \texttt{Sigma} (+\texttt{\_log}) & f32 & g\,cm$^{-2}$ & gas surface density, $N_\mathrm{col}\,\mu m_u$ \\
  \texttt{Sigma\_dust} (+\texttt{\_log}) & f32 & g\,cm$^{-2}$ & dust surface density, $0.01\,\Sigma$ \\
  \texttt{temp\_polaris\_mean} (+\texttt{\_log}) & f32 & K & density-weighted mean POLARIS dust temperature \\
  \texttt{temp\_starforge\_mean} (+\texttt{\_log}) & f32 & K & density-weighted mean STARFORGE dust temperature \\
  \texttt{B\_tot\_mean} (+\texttt{\_log}) & f32 & $\mu$G & density-weighted mean field strength $|\mathbf{B}|$ \\
  \texttt{cos2psi\_B\_mean}, \texttt{sin2psi\_B\_mean} & f32 & -- & direction of $z$ (Eq.~\ref{eq:coherence}): plane-of-sky field orientation seen in polarization \\
  \texttt{B\_coherence} & f32 & -- & $|z|/\sum_k n_k$, 0--1 (Eq.~\ref{eq:coherence}) \\
  \midrule
  \multicolumn{4}{@{}l}{\emph{Bands} (7 columns $\times$ 38 bands = 266)} \\
  \texttt{<b>\_I} & f64 & MJy\,sr$^{-1}$ & Stokes $I$ \\
  \texttt{<b>\_Q}, \texttt{<b>\_U} & f64 & MJy\,sr$^{-1}$ & Stokes $Q$, $U$ in the image frame (angles from $+$col towards $+$row) \\
  \texttt{<b>\_p} & f32 & -- & polarization fraction $\sqrt{Q^2+U^2}/I$ \\
  \texttt{<b>\_polangle} & f32 & deg & polarization (E-vector) angle $\tfrac12\arctan(U,Q)$ in $(-90,90]$; field angle $=$ \texttt{polangle} $+90^\circ$ \\
  \texttt{<b>\_I\_log}, \texttt{<b>\_p\_log} & f32 & -- & $\log_{10}$ of $I$ [MJy\,sr$^{-1}$] and of $p$ \\
  \midrule
  \multicolumn{4}{@{}l}{\emph{Profiles} (64 cells each, $k=00$ nearest the observer; 1152 columns)} \\
  \texttt{dens\_<k>} (+\texttt{\_log}) & f32 & cm$^{-3}$ & gas particle density $n=\rho/(\mu m_u)$, $\mu=2.33$ \\
  \texttt{dens\_dust\_<k>} (+\texttt{\_log}) & f32 & g\,cm$^{-3}$ & dust mass density, $0.01\,\rho$ \\
  \texttt{temp\_polaris\_<k>} (+\texttt{\_log}) & f32 & K & dust temperature computed by POLARIS (used for the images) \\
  \texttt{temp\_starforge\_<k>} (+\texttt{\_log}) & f32 & K & STARFORGE dust temperature (mass-weighted deposit) \\
  \texttt{B\_tot\_<k>} (+\texttt{\_log}) & f32 & $\mu$G & field strength $|\mathbf{B}|$ \\
  \texttt{B\_pos\_<k>} (+\texttt{\_log}) & f32 & $\mu$G & plane-of-sky field strength $\sqrt{B_\mathrm{col}^2+B_\mathrm{row}^2}$ \\
  \texttt{cos2psi\_B\_<k>}, \texttt{sin2psi\_B\_<k>} & f32 & -- & plane-of-sky field orientation $\cos2\psi_B$, $\sin2\psi_B$ \\
  \texttt{sin2gamma\_B\_<k>} & f32 & -- & $\sin^2\gamma_B=B_\mathrm{pos}^2/B_\mathrm{tot}^2$ (1: field in the sky plane; 0: along the line of sight) \\
  \texttt{B\_col\_<k>}, \texttt{B\_row\_<k>} & f32 & $\mu$G & field components along image $+$col and $+$row \\
  \texttt{B\_los\_<k>} & f32 & $\mu$G & line-of-sight component, positive away from the observer \\
\end{longtable}
\endgroup

\subsubsection{Value ranges}

Table~\ref{tab:ranges} gives the distribution of the main quantities. The sightline and band columns were
evaluated over all 23\,445\,504 rows of \texttt{v1}; the profile columns over all 64 cells of 300 randomly chosen
box views ($7.9\times10^{7}$ cells). All column densities in the dataset correspond to optically thin dust at
long wavelengths ($\tau_{850}\le0.01$, Sect.~\ref{sec:dataset:verification}). Intensities at
$\le5.8\,\mu\mathrm{m}$ are of order $10^{-108}\,\mathrm{MJy\,sr^{-1}}$ and carry no usable signal.

\begin{table}[ht]
  \centering
  \caption{Distribution of selected quantities: minimum, 1st percentile, median, 99th percentile and maximum.}
  \label{tab:ranges}
  \small
  \begin{tabular}{@{}lllllll@{}}
    \toprule
    Column & Unit & min & 1\% & median & 99\% & max \\
    \midrule
    \multicolumn{7}{@{}l}{\emph{Sightline (all rows)}} \\
    \texttt{N\_col}               & cm$^{-2}$     & $2.4\times10^{18}$ & $9.9\times10^{19}$ & $9.2\times10^{20}$ & $7.1\times10^{21}$ & $2.5\times10^{23}$ \\
    \texttt{Sigma}                & g\,cm$^{-2}$  & $9.3\times10^{-6}$ & $3.8\times10^{-4}$ & $3.6\times10^{-3}$ & $2.7\times10^{-2}$ & $0.97$ \\
    \texttt{temp\_polaris\_mean}  & K             & 8.27 & 12.5 & 14.5 & 16.0 & 54.4 \\
    \texttt{temp\_starforge\_mean}& K             & 6.16 & 11.3 & 13.3 & 15.5 & 53.8 \\
    \texttt{B\_tot\_mean}         & $\mu$G        & 0.96 & 9.6 & 25.3 & 58.3 & 1760 \\
    \texttt{B\_coherence}         & --            & $3.6\times10^{-5}$ & 0.029 & 0.47 & 0.96 & 0.999 \\
    \midrule
    \multicolumn{7}{@{}l}{\emph{Bands (all rows)}} \\
    $I$ (53\,$\mu$m)   & MJy\,sr$^{-1}$ & $3.2\times10^{-4}$ & 0.0074 & 0.018 & 0.072 & $2.1\times10^{6}$ \\
    $I$ (214\,$\mu$m)  & MJy\,sr$^{-1}$ & 0.083 & 3.3 & 20.5 & 87.4 & $1.2\times10^{5}$ \\
    $I$ (849\,$\mu$m)  & MJy\,sr$^{-1}$ & 0.0039 & 0.17 & 1.33 & 8.1 & 1380 \\
    $I$ (1380\,$\mu$m) & MJy\,sr$^{-1}$ & $8.0\times10^{-4}$ & 0.034 & 0.28 & 1.8 & 246 \\
    $p$ (53\,$\mu$m)   & --  & $1.4\times10^{-5}$ & 0.0059 & 0.098 & 0.259 & 0.355 \\
    $p$ (214\,$\mu$m)  & --  & $1.0\times10^{-5}$ & 0.0052 & 0.089 & 0.260 & 0.323 \\
    $p$ (849\,$\mu$m)  & --  & $2.2\times10^{-6}$ & 0.0051 & 0.089 & 0.261 & 0.324 \\
    $I$ (3.4\,$\mu$m)  & MJy\,sr$^{-1}$ & $7\times10^{-136}$ & $3\times10^{-112}$ & $8\times10^{-108}$ & $1\times10^{-104}$ & $3\times10^{-12}$ \\
    \midrule
    \multicolumn{7}{@{}l}{\emph{Profiles (300 box views, all cells)}} \\
    \texttt{dens}           & cm$^{-3}$    & 0.010 & 5.8 & 73.8 & 2310 & $9.9\times10^{5}$ \\
    \texttt{dens\_dust}     & g\,cm$^{-3}$ & $3.9\times10^{-28}$ & $2.3\times10^{-25}$ & $2.9\times10^{-24}$ & $8.9\times10^{-23}$ & $3.8\times10^{-20}$ \\
    \texttt{temp\_polaris}  & K            & 8.37 & 12.7 & 15.0 & 16.1 & 80.0 \\
    \texttt{temp\_starforge}& K            & 5.91 & 11.3 & 13.4 & 15.6 & 66.8 \\
    \texttt{B\_tot}         & $\mu$G       & 0.037 & 7.3 & 26.0 & 69.1 & 2380 \\
    \texttt{B\_pos}         & $\mu$G       & 0.0017 & 2.7 & 20.2 & 60.2 & 2370 \\
    \texttt{sin2gamma\_B}   & --           & $2.6\times10^{-8}$ & 0.019 & 0.77 & 1.00 & 1.00 \\
    \texttt{B\_los}         & $\mu$G       & $-1600$ & $-45.8$ & $-6.6$ & 31.8 & 1320 \\
    \bottomrule
  \end{tabular}
\end{table}

The intensity spectra of all sightlines peak between 160 and $250\,\mu\mathrm{m}$ (at $214\,\mu\mathrm{m}$ for
96.8\%), as expected for dust at 12--16\,K. The polarization fraction changes little with wavelength between
$53\,\mu\mathrm{m}$ and $1.38\,\mathrm{mm}$ (median 0.098 at $53\,\mu\mathrm{m}$, 0.089 at 214 and
$849\,\mu\mathrm{m}$; maximum 0.355).

\subsubsection{Reading the data}

Columns can be read selectively, which keeps memory small; for example, with \texttt{pandas} and \texttt{pyarrow}:
\begin{verbatim}
import pandas as pd, pyarrow.dataset as ds
root = "/scratch/sjafaritabar/starforge_polaris_dataset_v1/data"
cols = ["planck_hfi_353ghz_I_log", "planck_hfi_353ghz_p"] + \
       [f"dens_{k:02d}_log" for k in range(64)]
train = pd.read_parquet(f"{root}/train", columns=cols)
full = ds.dataset(root, format="parquet", partitioning=None)   # lazy, all splits
\end{verbatim}
To obtain images, a file is read one row group at a time (\texttt{pyarrow.parquet.ParquetFile.read\_row\_group})
and the 4096 rows are placed at (\texttt{pix\_row}, \texttt{pix\_col}).

% -----------------------------------------------------------------------------------------------------
\subsection{Verification}
\label{sec:dataset:verification}

Each step has unit tests (step 2: 7, step 3: 13, step 5: 4; all pass), and an independent audit
(\texttt{audit\_pipeline.py}) re-derives every stage from the raw files without using the pipeline's code. It
passed all 15 checks for both dataset versions:
\begin{enumerate}
  \item Every boxed core is in its step-1 catalogue with the same peak, peak particle and mass; box times equal
  the snapshot header times; no two boxes of a snapshot overlap; families re-derived from scratch reproduce all
  317 \texttt{family\_id}s.
  \item Exactly one grid per box; for every box the summaries, grids, star lists, command files (photon numbers,
  radius, dust, wavelengths, views) and POLARIS outputs are mutually consistent, the grid density is unchanged by
  POLARIS, and dust temperatures and alignment radii are within range.
  \item The $850\,\mu\mathrm{m}$ intensity in $\mathrm{MJy\,sr^{-1}}$ matches the optically thin model
  $\sum_k B_\nu(T_k)\,\Delta\tau_k$ computed from the grid, which tests the units, pixel solid angle and distance:
  POLARIS/model has a per-box median of $1.002$ (range $1.001$--$1.005$); $\tau_{850}\le0.01$.
  \item No box has an outlying POLARIS-vs-STARFORGE dust-temperature offset ($>5\sigma$).
  \item Every one of the 1443 columns was recomputed from the raw grids and FITS files with explicit index formulas
  (Table~\ref{tab:views}) for 1080 random sightlines in 36 boxes ($1.56\times10^{6}$ values); the worst
  relative (or, for logarithms and angles, absolute) error is $3.8\times10^{-6}$, i.e.\ single-precision rounding.
  \item \texttt{splits.csv} covers every box once with the step-2 family IDs; no family is in more than one split;
  all families with more than 50 boxes are in training.
  \item Polarization fractions at $\ge50\,\mu\mathrm{m}$ are physical (maximum 0.355), and the spectra of all
  sightlines peak at 60--$500\,\mu\mathrm{m}$ (214\,$\mu$m for 96.8\%).
\end{enumerate}
\texttt{validate\_dataset.py} additionally confirms the row counts, the absence of duplicate sightlines and of
boxes or families in more than one split, and scans every column of every row for non-finite values.

Across the whole dataset, column densities range over $2.4\times10^{18}$--$2.5\times10^{23}\,\mathrm{cm^{-2}}$
(median $9.1\times10^{20}$), density-weighted POLARIS dust temperatures over 8.3--54.4\,K (median 14.5\,K),
density-weighted field strengths over $0.96$--$1760\,\mu\mathrm{G}$ (median $25.7\,\mu\mathrm{G}$), and
\texttt{B\_coherence} over $4\times10^{-5}$--$0.999$ (median 0.48).

% -----------------------------------------------------------------------------------------------------
\subsection{Limitations}
\label{sec:dataset:limitations}

\begin{itemize}
  \item \textbf{Radiation field.} Each box is illuminated by an unattenuated Mathis field with $G_0=1$. In
  STARFORGE the field reaching the box surfaces is attenuated by the surrounding cloud (snapshot 440: equivalent
  $G_0$ median 0.20, 5--95\% range 0.11--0.46). POLARIS dust temperatures are therefore higher than the STARFORGE
  ones, by a per-box median of $+1.49\,\mathrm{K}$ (5--95\%: $+0.29$ to $+2.76\,\mathrm{K}$). Both temperatures
  are in the dataset.
  \item \textbf{Isolated boxes.} Only gas inside the $2.2\,\mathrm{pc}$ cube emits and absorbs; foreground and
  background cloud emission and extinction are absent.
  \item \textbf{No stochastic heating}: bands at $\le25\,\mu\mathrm{m}$ are not reliable.
  \item \textbf{Ideal observations.} No noise, no beam, a single distance (140\,pc) and a pixel equal to a grid
  cell. Noise and beam have to be added when training.
  \item \textbf{No gas kinematics.} Velocities are not included, so velocity-based field-strength estimates
  (e.g.\ Davis--Chandrasekhar--Fermi) cannot be reproduced from the dataset alone.
  \item \textbf{One simulation.} All boxes come from one cloud; families make the evaluation sets independent in
  gas, not in environment. Late snapshots, which contain all stars, are $0.41\,\mathrm{Myr}$ apart.
  \item \textbf{Monte Carlo.} Dust temperatures carry $\approx0.008\,\mathrm{K}$ of Monte Carlo noise, and runs
  are not bitwise reproducible.
\end{itemize}

% -----------------------------------------------------------------------------------------------------
\subsection{Data products and code}
\label{sec:dataset:files}

Code: \texttt{STARFORGE/creating\_the\_dataset/} (\texttt{1.finding\_cores}, \texttt{2.selecting\_boxes},
\texttt{3.preparing\_polaris\_grids} with its \texttt{diagnostics}, \texttt{4.running\_polaris},
\texttt{5.building\_dataset}, \texttt{audit\_pipeline.py}); photon studies in
\texttt{STARFORGE/M2e5\_R30\_Res585/POLARIS/isrf\_photon\_noise\_test} and \texttt{STARFORGE/M2e5\_R30\_Res585/POLARIS/star\_photon\_test}.
Data: snapshots in \texttt{/scratch/STARFORGE/M2e5\_R30\_Res585}; core catalogues in
\texttt{1.finding\_cores/dendrogram\_cores\_v5.2}; boxes in \texttt{2.selecting\_boxes/boxes\_2.2pc};
grids in \texttt{3.preparing\_polaris\_grids/grids\_64}; POLARIS output in
\texttt{/scratch/sjafaritabar/polaris\_runs} ($\approx70\,\mathrm{GB}$); datasets in
\texttt{/scratch/sjafaritabar/starforge\_polaris\_dataset\_v1} and
\texttt{/scratch/sjafaritabar/starforge\_polaris\_dataset\_v1\_stars\_eval}; audit reports in
\texttt{/scratch/sjafaritabar/audit}.
